\documentclass[12pt]{article}

\usepackage{fontspec}
\setmainfont{Latin Modern Roman}

\usepackage[french]{babel}

\usepackage[
  letterpaper,
  left=3cm,
  right=5.5cm,
  top=2.5cm,
  bottom=2.5cm,
  marginparwidth=3.2cm,
  marginparsep=6mm
]{geometry}

\usepackage{setspace}
\onehalfspacing

\setlength{\parindent}{0pt}
\setlength{\parskip}{0.5\baselineskip}

\title{Journal\\[0.5em]\large Journal de Bernhard Albrecht Stettler. Volume 1}
\author{}
\date{De l'an 1774 jusqu'\`a 1786}

\begin{document}

\maketitle

Le but qu’on se propose généralement en écrivant le journal de sa vie, est de parvenir plus
facilement à une parfaite connaissance de soi même. Le fidèle tableau de nos actions, de nos
jugements, et de nos pensées, nous fait connaître les mouvements secrets de notre âme, ils se
développent avec franchise et se présentent à notre imagination sous leur vrai point de vue.
L’agréable souvenir que réveille en nous le récit des événements passés de notre vie, le plaisir
qu’éprouveront nos descendants à la lecture des faits historiques, les leçons qu’ils puiseront dans
nos malheurs, sont les fruits de ce long travail... Aiguillonné par ces motifs je n’ai point redouté
les obstacles attachés à un ouvrage aussi étendu ; la haine de l’oisiveté et l’amour des occupations
utiles me les feront aisément vaincre. J’écris dans la langue française, par ce qu’elle m’est plus
étrangére que la maternelle, cet exercice me perfectionnera et donnera plus de facilité à mon style
pour l’avenir Je commence mon journal seulement à ma 25ème année, c’est-à-dire le 1er novembre 1799
Cormondrèche ; une conversation que j’ai eue ces jours passés avec M. le Baron de Chambrier,me fait
prendre cette résolution ; j’en dirai d’avantage quand je serai arrivé à cette époque. Je veux
maintenant reculer à ma première jeunesse et retracer tous les événements majeurs de ces temps, dont
la plupart ont beaucoup influé sur mes destinées présentes. Je suis né à Frienisberg le 6 mars 1774.
Trois ans de ma vie se sont écoulés dans ce séjour, je puis dire dans le bonheur, car je vivais
inconnu à moi-même. Mon père dont la préfecture de ce baillage allait à sa fin, se rendit à Berne
avec toute sa famille, je passais toute mon enfance alternativement dans cette ville ou à Kirchberg,
une campagne de mes parents. En hiver je fréquentais le collège, en été des précepteurs étaient
chargés de mon éducation. Les dix premieères années mes occupations se bornèrent à apprendre les
principes de notre religion, ceux de la langua latine, et du grec, je renonçais à cette dernière
langue aussitôt qu’il fut décidé que je ne me vouerais pas au Saint ministère. Le souvenir le plus
désagréable qui me reste de ce temps de mon adolescence, est la désunion constante entre mon frère
cadet et moi. Que de sujets d’amertume et de douleurs pour nos parents, qui ne désiraient que le
bonheur de leurs enfants ; nos torts étaient sans doute réciproques, les miens étaient plus graves
comme aîné, et je me les reprocherai toujours. La suite nous a prouvé, et j’aurai occasion d’en
parler, que ce peu d’accord entre nous

\marginpar{\footnotesize\raggedright 1787. 1788. 1789.} provenait d’une incompatibilité d’humeur et
de caractére. Il est un grand nombre de familles où les discordes entre frères et soeurs sont très
ordinaires, on se querelle pour s’aimer davantage après ; mais la nôtre était malheureusement d’une
toute autre espèce, et son influence a fixé ma carriére future. En 1787 je quittai le collège, la
derniere année m’a été préjudiciable ; un régent vieux et infirme ne sut profiter des bonnes
dispositions de ses écoliers, et en sortant de sa classe j’étais beaucoup plus ignorant que lorsque
j’y entrai. C’est dans cette année que l’institut politique prit son origine ; les personnes les
plus éclairées furent choisies pour mettre la dernière main à l’instruction de la jeunesse
patricienne de Berne. Mon père conçut les meilleures espérances d’un établissement à la fondation
duquel il avait même contribué, j’y fis mon entrée au mois de novembre et je l’ai fréquenté avec
assez de succès, pendant trois hivers consécutifs. Je commençais en même temps mes leçons dans la
langue française, des progrès rapides assurèrent à mon maître l’amitié et la confiance de mes
parents, mais ils eurent à s’en repentir par la suite. Tandis que tous leurs efforts tendaient à me
frayer le chemin du bonheur, par des conseils sages, un choix de livres instructifs et religieux, et
surtout par un exemple sans tache ; cet homme, vil corrupteur de l’innocente jeunesse sut entraver
leurs bienveillants desseins. à peine étais-je parvenu à cet âge où les passions humaines se
développent, la conduite de mon maître changea ; il me dit un soir : maintenant mon cher ami, que
vous êtes plus avancé dans l’étude de la langue française, je ne veux plus vous occuper de thèmes et
de grammaire, notre leçon doit devenir une conversation suivie, nous repasserons ainsi d’une manière
agréable, l’histoire, la géographie vous vous accoutumerez à parler par ce moyen sur tous les objets
qui se présentent, et le but de vos parents sera accompli. Les premiers jours se passèrent en effet
comme il me l’avait dit ; peu à peu nos entretiens ne furent plus ceux de l’instruction et à la fin
ils devinrent ceux de la débauche et du vice. Je me garderai bien de souiller ce papier avec le
détail des honteux moyens que cet homme dépravé mit en jeu pour conduire mon jeune coeur à
l’égarement ; je pense avec effroi à la joie impudente qu’il manifestait en s’apercevant de
l’irritation que causaient sur mes nerfs ses infâmes propos, et avec quelle satisfaction il
m’enseignait l’art du soulagement funeste qu’on se donne à soi-même. Mais ce cruel séducteur ne
borna par là son crime. Il me fit connaître dans sa maison une servante, et me donna tous les
ren–seignements

\marginpar{\footnotesize\raggedright 1789. 1790.} pour m’en procurer la jouissance ; la réussite ne
fut pas difficile, et je fis sous ses auspices le premier pas au libertinage. C’était à la fin de
1789 je me faisais instruire dans le même temps par mon beau-frére L. À 11. heures du matin mon âme
s’abandonnait avec plaisir aux douceurs d’une occupation aussi salutaire, et le soir à 5 heures le
vice reprenait son empire. Cependant les cris de ma conscience se faisaient vivement entendre, mes
premières erreurs furent suivies du plus sincère repentir ; le soir avant de me coucher je priais le
ciel de me préserver des tentations ; je me regardais au miroir, m’imaginant que la débauche aurait
déjà ravagé mes traits, tant j’étais ignorant et simple. Tout ce que j’ai pu gagner sur moi, était
de dire à mes parents que je croyais a présent mes leçons dans la langue française superflues ; je
sus en même temps empêcher qu’on mît mon frére cadet à la même école, sans découvrir mes raisons.
Mais il s’agissait encore de me corriger du vice ; j’avais sacrifié les femmes à l’autre misérable
jouissance qui avilit l’homme et le met au rang des brutes. Pendant plusieurs mois toutes mes peines
n’eurent point de succès ; je ne me laissai Dieu la raison pas décourager, la raison vint à mon
secours et peu à peu elle triompha au point que je n’eus plus à redouter de fâcheuses suites pour
l’avenir. Je crois que c’est à peu près dans la même époque que m’arriva un accident qui m’était des
plus pénibles. Nous faisions tous les soirs entre amis une partie au reversi, tantôt chez l’un
tantôt chez l’autre ; je m’apercevais que plusieurs de mes camarades ne jouaient pas avec la loyauté
requise. Envisageant la chose plutôt sous un aspect indifférent que sérieux, j’imitais les mêmes
ruses, et je ne sais si c’est par plus de témérité ou moins d’adresse que mes camarades ont surpris
mes coups ; ces Messieurs ne firent semblant de rien, et le jour après, au lieu de les voir tous
chez moi, je trouve un billet signé de tous, où ils se déclarent mes plus implacables ennemis. Ce
coup inattendu m’affligea in– finiment, j’aimais beaucoup mes camarades, et je ne pouvais m’imaginer
comme ils pouvaient rompre une liaison, pour des fautes où ils étaient mes séducteurs et mes
complices ; J’en fis le rapport à mon père, qui eut la bonté de me dicter une lettre de
raccommodement, elle n’a pas eu d’effet, par ce qu’un de ces amis se donna toutes les peines
d’engager les autres, à ne pas se réconcilier, quoiqu’il voyait qu’ils y étaient trés disposés.
Telle était ma situation, au commen–cement de l’an 1790. Abandonné de mes amis, libertin
inconsidéré, je paraissais vouloir tromper les espérances des meilleurs parents, qui ignoraient ce
défaut

\marginpar{\footnotesize\raggedright 1790.} mais la providence a eu soin de mes jours. Heureusement
que mes études me laissaient fort peu de tems à la dissipation ; je passai les matins à l’institut
et les après-midi chez mon beau-frére D. Il n’y avait que les soirées qui furent quelquefois assez
mal employées. Je fus arraché à ce genre de vie, par la résolution que prit mon père de me placer au
pays de Vaud. Je partis le 29 avril 1790 pour me rendre à Morges, chez M. le secrétaire baillival
Pache ; j’eus le bonheur de m’y faire des amis, j’en avais un grand besoin, la privation ... des
premiers me rendit plus prudent envers les autres. Je passais six mois fort agréablement ; des
courses à Lausanne, à Aubonne et Rolle, un voyage autour du lac, des bals et des fêtes et la société
des femmes, varièrent les plaisirs de ce temps-là. Cependant l’amour vint troubler le repos de mon
coeur ; ce n’était point pour la première fois que je sentais sa flèche ; j’aimais à Berne, si on
peut aimer avant de connaître l’amour ; mais le nouveau sentiment prit la place de l’autre, et je me
livrais entièrement à ce nouvel objet ; l’aimer et le lui dire fut l’affaire du moment ; je lui
parlais en même temps de mariage, mais l’envie m’en passa un jour, que je vis qu’elle avait ses
petites mains pleines de verrues. Âge heureux ! où les impressions que nous recevons de ce charmant
sexe, disparaissent avec les circonstances qui les font naître. (Cette dernière phrase est sans
logique.) Pendant mon séjour à Morges, un de mes amis me dit une grossiereté ; j’en fis des plaintes
à mes autres camarades, qui me donnèrent l’idée d’arranger un duel aux pistolets, chargés simplement
à poudre ; ils se proposèrent d’être nos secondants et de faire la chose eux-mêmes. Ils cherchèrent
en conséquence des pistolets et fabriquaient des cartouches à balles, se proposant de la mordre en
chargeant ; j’écrivis un billet d’invitation à mon adversaire, comme il était destiné au service
étranger, je savais d’avance qu’il ne re– fuserait pas ; il me répondit qu’il serait prêt à l’heure
indiquée, et comme il demeurait avec un de ceux qui préparaient la farce, il le pria d’être son
second, l’autre y consentit après quelques difficultés, qui servirent à mieux cacher l’entreprise.
Nous nous acheminâmes vers le bois, la tranquillité et le sang froid de l’adversaire qui fumait sa
pipe comme si de rien n’était, nous déconcerta un peu ; on tira les numéros, il eut l’un, et se mit
en posture à une distance

\marginpar{\footnotesize\raggedright 1790. 1791.} de 15 pas ; j’avoue que malgré que je savais qu’il
ne m’en arriverait rien, j’avais une peur du diable ; le coup partit avec un fort petit bruit, il y
avait trop peu de poudre, je fis le généreux en tirant en l’air ; nous nous embrassâmes et il fut
convenu que les ennemis réconciliés payeraient un goûter à Messieurs les secondants. J’ignore encore
dans ce moment si mon adversaire a su de quoi il s’agissait. Le mois de septembre fut le dernier que
je passais au pays de Vaud ; la maladie de M. Pache causait mon prompt retour ; j’ai logé pendant 15
jours au château, chez Mr le Baillif de Ryhiner, qui m’a comblé d’honnêtetés. La satisfaction de
revoir mes parents qui se trouvaient à Kirchberg était extrême, nous rentrâmes en ville en novembre.
En décembre et une partie de janvier 1791 se passa à Ursenbach ; mon père venait d’hériter le
ministre Aechler un cousin fort éloigné. En revenant à Berne, je continuais à fréquenter l’institut.
Le 6 de mars en commençant ma 16ème année, on me plaça dans la Chancellerie, plutôt pour me faire
travailler, que de suivre cette carrière. Il faut que mes parents se soient aperçus que la
séparation entre mon cadet et moi produisait un bon effet ; il est certain que pendant ma 1ère
absence nous nous écrivions des lettres dont le contenu était des épanchements réciproques d’une
tendre amitié, et des promesses solennelles de ne plus nous rendre nos jours amers par la discorde
et la désunion ; mais nous ne tînmes point encore parole. Mon père résolut de mettre mon frère au
pays de Vaud ; nous partîmes le 1er de mai, mon père et ses deux fils, pour conduire le cadet chez
M. le Doyen Bourgeois à Saint Cierge ; nous passâmes par Payerne et Yverdon, et nous avions le plus
grand plaisir de voyager avec un père aussi chéri. On nous reçut à Saint Cierge le mieux du monde,
nous y laissâmes mon frère, et prîmes la route de Moudon, pour nous en retourner à Berne. Peu de
temps après, je crois au mois de juin, on alla à Kirchberg. Dans les belles soirées de l’été, à la
promenade, mes parents s’occupaient du sort futur de leurs enfants. Mes deux soeurs étaient établies
; mon frère aîné, dont l’amour pour l’étude devait faire un homme de lettres, venait de quitter
l’université de Göttingen, il était attendu à la maison aussitôt que ses voyages à Berlin à Vienne
et à Francfort, où il voulait assister au couronnement de l’empereur Léopold, seraient à leur fin.
Mes les deux fils cadets demandaient encore les sollicitudes

\marginpar{\footnotesize\raggedright 1791.} des meilleurs parents. Une lettre de mon père reçut dans
ce temps de M. Stettler, Lieutenant Colonel au régiment de Rochmondet, au service du roi de
Sardaigne où il lui offre une place de sous-lieutenant pour un de ses fils, le mit dans une espèce
de perplexitude. Ennemi de l’état militaire, par des exemples mal-heureux qui le touchaient de près
; le brouillard qui paraissait déjà obscurcir alors l’horizon politique, faisaient redouter à mes
parents une carrière qui sous une trompeuse apparence cache les désagréments qui en sont
insépa–rables. Ils crurent cependant devoir m’informer de l’espèce d’établissement qui venait de
s’offrir, ma mère me parla de la lettre en question, et sans y mettre la moindre importance elle me
dit : je pense que l’idée de servir ne te vient pas ? Ma résolution contraire mit le comble à son
étonnement, elle me conseilla de réfléchir mûrement avant de me décider ; mon père me fit voir les
dangers attachés à cette vocation sans s’opposer à mes voeux de la suivre. Rendu à moi-même je
sentis toute l’importance du parti que je me proposais à prendre ; la brouillerie avec mes amis, la
désagréable impression qui en était réstée au public, le peu d’harmonie dans laquelle je vivais avec
mon frère, la déclaration de mon père de ne point me placer dans la chancellerie où je me plaisais,
la joie de revoir et de vivre avec R. Stettler de Bipp mon ancien ami, et enfin la vanité de porter
l’uniforme, furent les motifs qui me déterminèrent à choisir le métier des armes ; j’en fis la
déclaration à mes parents, ils ne contrarièrent point ma volonté, mon père écrivit au Lieutenant
Colonel Stettler, qu’un de ses fils se décidait à accepter la place qu’il avait eu la complaisance
de lui offrir ; Il n’y en avait pas de vacante ce moment là, mais on me promit la 1ère qui
vaquerait. Les affaires d’état appelèrent mon père en ville au mois d’août, nous allâmes nous y
établir. J’arborais avec empressement la cocarde bleue sur un chapeau militaire ; un camp de
plaisance était alors tout près de Berne, j’y allais tous les jours et je pris ainsi insensiblement
le goût de mon état. Ce qui ne contribua pas peu à me le rendre agréable, fut l’arrivée de mon ami
Stettler, qui reçut un congé de quelques semaines. Dans le temps que son bataillon était en garnison
à Carouge près de

\marginpar{\footnotesize\raggedright 1791.} Genève. Tout ce qu’on m’avait dit de pénible du métier
fut dissipé en un moment ; Stettler partit pour rejoindre son corps, je brûlais dès lors d’envie de
le suivre, il me promit de m’instruire aussitôt qu’un officier viendrait à quitter d’une manière ou
d’une autre. Au milieu de septembre ma famille retourna à Kirchberg pour y passer l’automne, je la
suivis ; pendant ce temps je m’occupais à me familiariser avec le pays que je devais habiter sous
peu ; Busching en fait une description si avantageuse qu’en la lisant, je fus au comble de la joie ;
mon imagination était si frappée de ma carrière future, que je me rappelle d’avoir dit un soir à la
cure de K[irchberg] en grande compagnie à une partie de jeu : depuis que je suis en Piémont je n’ai
jamais eu d’aussi mauvaises cartes ; Tout le monde se moqua de moi, en me questionant beaucoup sur
ce joli pays. Cependant l’époque où je devais réaliser mes voeux s’approchait à grands pas ; nous
rentrâmes en ville au commencement de novembre je crus devoir faire encore un cours de mathématiques
avant mon départ, M. Dorner fut appelé, je pris aussi des leçons dans la langue italienne du vieux
Cérichelli, cette étude fit mes délices. Il ne s’agissait plus que de savoir quand il y aurait une
vacance au régiment, on m’avait offert d’y entrer comme cadet mais cela ne m’agréa point. Dans cette
continuelle attente je trouve un soir en arrivant à la maison une lettre sur ma table, je reconnus
l’écriture de mon ami Stettler, des espérances naquirent dans mon coeur, et j’y vois en effet que la
mort en duel du lieutenant Porta faisait vaquer une place. L’idée d’entrer au corps sous d’aussi
funestes auspices, troubla pour un moment mon plaisir extrême ; mes parents profitèrent d’une aussi
malheureuse circonstance pour me faire les plus paternelles exhortations, elles sont toujours
restées

\marginpar{\footnotesize\raggedright 1791 et 1792.} - profondément gravées dans mon coeur. Enfin cet
événement décida de mon sort ; mon père écrivit au brigadier Rochmondet pour lui demander la
sous-lieutenance dans toutes les formes, sa réponse arriva fort vite, le brevet fut demandé à la
cour, et le colonel du régiment me donna l’agré–ment de rester encor quelque temps chez moi, pour
prendre tous les arrangements possibles pour mon départ. Il y avait alors plusieurs officiers
bernois en semestre, de ce nombre étaient le lieutenant-colonel Stettler, les capitaines Ernst,
Weber, et Bucher, et le lieutenant d’Herbort. Mon cousin Bucher prit particulièrement le plus grand
intérêt à moi, il me conduisit chez tous ces officiers, il me donna tous les renseignements pour me
mettre en uniforme, il arrangea son retour au corps de manière à me prendre avec lui. Ainsi se finit
l’année 1791. Le commencement du mois de janvier 1792 se passa en préparatifs de départ ; tous les
métiers étaient en mouvement pour achever mon équipement, on me fit même allez chez le peintre
Maurer, qui traça mes traits avec des craies en couleurs, le portrait fut très ressemblant. C’était
toujours une fête pour moi de mettre l’uniforme, je ne m’imaginais point alors, que quelques années
après je le poserai avec le même plaisir que je l’ai pris, comme on le verra par la suite. Mon
départ fut fixé au 25 courant, le soir du 24 je pris congé de ma famille en versant des larmes. Nous
nous mîmes en voyage le lendemain à 5 heures du matin, Messieurs Bucher, Herbort et mon frère aîné
qui m’accompagna jusqu’a Lausanne. Nous comptions aller jusqu’à Moudon le 1er jour, mais la Broye
avait si grossie par la fonte des neiges, qu’elle s’est formée en torrent et nous a fermé le passage
à Marnans ; la poste en voulant passer ce jour-là se renversa, le carrosse était encore au milieu de
l’eau à notre arrivée. L’inondation nous arrêta jusqu’au 27, nous allâmes diner à Lausanne ; le
cousin Bucher invita le capitaine du régiment Mazars l’aîné, cet officier me donna une singulière
idée du

\marginpar{\footnotesize\raggedright 1792.} corps, il ne m’adressa pas la parole. Après diner il
fallut me séparer de mon frère, je le chargeais encore de mille choses pour les personnes chéries
que j’avais quittées et mes larmes coulèrent de nouveau. Nous arrivâmes le 28 au soir à Genève, nous
y passâmes le 29. Je fis ce jour plusieurs petites emplettes. M. le capitaine Jayet qui commandait
alors le détachement de Chêne, dîna avec nous à la Couronne, il démentit l’opinion que M. Mazars
m’avait fait concevoir, par des manières amicales et honnêtes Le 30 janvier est le jour remarquable
de mon entrée au régiment. Nous sommmes arrivés à la Bonneville, où était la garnison du 3ème
battaillon, un moment avant midi ; en descendant à l’auberge des Trois Maures nous y trouvâmes
plusieurs officiers qui voulaient se mettre à table ; le cousin Bucher me présenta et je fus très
bien reçu, on nous invita à dîner. Je vis arriver à ma grande joie mon cher ami Stettler ; le ton
amical et gai qui régnait pendant le repas parmi les officiers me plut infiniment. Je vis partir
avec le plus grand regret Messieurs Bucher et Herbort qui durent aller joindre leurs battaillons
respectifs à Annecy. Le premier m’avait servi de père pendant tout le voyage, il me donna les
meilleurs conseils pour ma conduite future, il me fit le choix des officiers dont je devais le plus
rechercher l’amitié, et finit par m’inviter à venir passer quelques jours près de lui, pendant le
carnaval prochain, pour faire connaissance avec les autres officiers du régiment. Herbort, qui
m’était un peu parent du côté de ma mère m’a aussi donné pendant notre course maintes preuves
d’amitié, il voyait avec plaisir mon goût pour l’occupation et me dit très sérieusement, que j’étais
un jeune homme dont il espérait faire quelque chose, il me promit de me montrer plusieurs pièces
curieuses quand je viendrai le voir à Annecy. Ces Messieurs partis, l’ami Stettler me conduisit chez
lui, il m’assigna une chambre dans son logement, il me proposa d’avoir son domestique, le fifre
Borgna, en commun, et me témoigna tout ce qu’on peut attendre d’un bon, vrai et sincère ami ; c’est
de lui que je reçus les premières directions pour me mettre au fait du service militaire, il
m’indiqua le Sieur sergent Roselli pour instructeur d’exercice, il m’ accompagna au commencement à
l’appel de ma compagnie, à la visite d’hôpital,

\marginpar{\footnotesize\raggedright 1792.} c’est du meilleur de mon coeur que je lui rends ici,
l’hommage de ma reconnaissance et de ma gratitude. Son amitié infatigable ne se borna point à ce que
j’ai cité plus haut, il voulut encore me procurer ses amis, il me présenta à tous, leur parla en ma
faveur les priant de partager avec moi l’amitié qu’ils lui avaient vouée jusqu’ici. Je ne veux dans
ce journal caractériser les différentes personnes avec lesquelles le hasard m’a souvent contraint à
vivre ; je me contenterais de parler de celles qui m’ont intéressé le plus. M. le
capitaine-lieutenant Muret fut de ce nombre, son esprit, sa bonne humeur et sa politesse doivent lui
assurer le suffrage de tous ceux qui le connaissent, c’est avec lui que j’ai d’abord désiré de
former une liaison moins générale que le sont celles des régiments. Je ne fus pas content de mon
capitaine M. Bégoz, homme d’un certain âge qui me paraissait froid, sombre, mélancolique, je lui
préférais Jonquiére le cadet capitaine-lieutenant aux grenadiers ; Stettler, Muret et lui étaient
les officiers que je fréquentais le plus à Bonneville ; j’y fis aussi la connaissance du Père
Ramini, un Sarde, et aumonier du régiment, et celle de Caniggia chirurgien du corps. Le 2 février
j’écrivis pour la 1ère fois à mes parents, je leur fis des détails sur mon nouveau genre de vie, ils
apprirent avec plaisir que j’étais satisfait de mon sort. Il est sûr que je me trouvais fort heureux
; entièrement maître de moi-même, entouré d’amis qui m’étaient chers, dans une agréable petite
ville, il me restait peu à désirer. Quoique nous étions alors en hiver le temps était fort doux, il
n’y eut pas de jours que nous n’allions promener. Le bruit de mon arrivée parvint à la Roche, autre
petite ville du Faucigny, où une centurie du régiment était en garnison sous les ordres du major
Zehender. Mon ami de Watteville d’Oberhoffen vint me voir, et m’invita d’aller à la Roche, d’où nous
n’étions éloignés qu’à deux lieues ; Stettler et moi louâmes des chevaux, et dans une heure nous
fûmes à notre destination. Les officiers de ces deux compagnies me reçurent très amicalement,
surtout les quatre frères Zehender qui étaient logés ensemble, ils nous logèrent, et nous
conduisirent dans les maisons de Livet et Détoire. Le ton sans gêne et vraiment amical sur lequel

\marginpar{\footnotesize\raggedright 1792.} on se trouvait dans ces deux familles nobles, me donna
une bonne idée de la societé en Savoie ; je félicitais mes camarades d’avoir autant de ressources
dans ce vilain trou de la Roche. Nous n’avions à Bonneville que les maisons du comte d’Écluse où on
était sur un pied gênant, et celle de M. Dumont où il y avait une petite vieille femme qui avait de
l’esprit, mais à l’âge de 70 ans ; elle raffolait d’un officier du régiment qui se trouvait alors
détaché. Être bernois et sociable étaient de suffisants titres à l’amitié de Messieurs Zehender, je
connaissais le cadet depuis fort longtems, et la connaissance avec les aînés fut bientôt faite ; ils
m’obligèrent de leur promettre de venir les voir souvent ; leur premier accueil m’avait été trop
agréable pour ne pas désirer à y retourner. Pendant notre visite, Watteville était froid et réservé,
et je lui en aurais fait des reproches, si on ne m’avait dit à temps que l’amour en était la cause.
Nous nous séparâmes de nos amis, pour re–joindre nos pénates, Stettler me dépeignit en chemin le
différent caractère des officiers que nous avions vus, et se loua beaucoup de l’hospitalité de
Messieurs Zehender. De retour à Bonneville, je trouve une lettre du cher cousin Bucher qui me dit
que M. le brigadier Rochmondet permît que je vins à Annecy, et qu’on se réjouissait de m’y voir ;
peu de tems après je partis pour me rendre à cette amicale invitation ; c’était au milieu du
carnaval ; je fis le voyage avec Dietzi et Manton, un officier du régiment de Genevois, que j’avais
trouvé à la Roche ; nous nous arrêtâmes à Annecy-le-Vieux, où les parents de Manton occupaient un
beau château ; on nous y servit quelques restaurations, et nous continuâmes notre route. Arrivés à
Annecy sur les 5 heures du soir, Dietzi me conduisit dans son logement pour y faire un peu de
toilette, afin d’oser paraître à un bal qui devait avoir lieu le même soir ; je fus très content de
mon perruquier, qui me prodiga le titre de chevalier, je pensais en effet que les chevaliers de ce
pays étaient donc à peu près faits comme moi. Lorsque nous fûmes prêts, je priai Dietzi de
m’accompagner chez mon cousin Bucher, je ne le trouvai pas à sa maison, nous le rencontrâmes en rue
; du moment qu’il me vit il me combla de nouveau des mêmes bienfaits dont j’avais

\marginpar{\footnotesize\raggedright 1792.} dejà eu à me louer. Il s’empara de moi pour me présenter
à tous les officiers du régiment ; nous entrâmes au café, où j’eus bien des révérences à faire ; sur
le tard on s’en fut au bal, où la même cérémonie de saluer ceux des officiers que je n’avais par
encoe vus, recommença. Je me réjouissais beaucoup de voir Jonathan de Graffenried, que j’avais connu
à Berne, je l’abordais d’un air de connaissance, qui me fut rendu froidement et avec protection ;
j’en ai été fort étonné. La manière dont je fus traité par tous les officiers du corps, me plut
beaucoup, je me trouvais soulagé de cette première apparition ; en me pressa de danser, mais je n’en
avais nullement envie, je me retirai dans la chambre à côté, près du feu. Content d’être un moment
seul, je me livrais à mille réflexions sur ma situation présente, j’étais heureux, tout riait à mon
imagination, et ce bonheur n’a pas discontinué pendant mon séjour à Annecy. Après le bal, le cher
cousin vint me prendre, pour nous en aller chez lui, où un petit joli souper nous attendait ; il me
questi-onna beaucoup sur tout ce que je venais de voir, et il jouissait de ma nouveauté. Le
lendemain il me mena au quartier, il s’y prépara une cérémonie qui fit une douloureuse impression
sur mon coeur ; huit soldats passè́rent par les verges, la journée était froide, et le sang qui
couvrait leur dos déchirés se glaçait à mesure qu’il sortait du corps. Les officiers qui me virent
sensible me consolaient en me disant que je verrais bien autre. Après cette cruelle scène je fis
avec mon cousin la tournée des chambrées de sa compagnie qui était la plus belle du régiment, il
recommanda la propreté à ses gens ; nous mangeâmes de leur soupe que j’ai trouvé bonne. Le reste de
la matinée se passa en visites chez tous les officiers ; je commençais par le chef qui me reçut très
bien et m’invita à dîner, ce que firent tour à tour la plupart des officiers ; ils me rendirent tous
ma visite. Je suis resté à Annecy environ 15 jours, je n’y vis point Herbort, il était en
détachement, cet espace s’écoula avec une rapidité étonnante, la variation des jours était celle des
plaisirs. J’ai fait dans ce temps la connaissance de la maison de Seyssel où mon coeur aurait couru
de grands risques

\marginpar{\footnotesize\raggedright 1792.} si mon départ ne les avait prévenu. La Bonneville me
parut une solitude après avoir quitté Annecy ; je donnais à Stettler des nouvelles de ses camarades,
j’avais même fait connaissance avec plusieurs officiers, qui lui étaient encore inconnus, quoiqu’il
y avait déjà près de trois ans qu’il était au corps. Les 20 premiers jours de mars se passèrent fort
tranquillement ; la mort de l’Empereur Léopold nous affecta tous. Le 21 du mois je fus commandé
d’aller en détachement à Avusy petit village à 3 lieues de Genève, et à 8 de la Bonneville ; je
m’acheminais de bon matin, à cheval, pour aller dîner à Carouge, où je vis pour la 1ère fois des
troupes piémontaises ; je trouvais à l’auberge des Trois Rois huit officiers de la légion légère ;
ils me permirent de me mettre à leur table pour dîner ; pendant tout le repas je n’entendis que leur
jargon qui m’était absolument inconnu, peu avant de se lever de table, un de ces Messieurs me
demande si je savais le piémontais ? je lui répondis que non, ce qui le surprit beaucoup. Après midi
j’allai prendre les ordres du commandant de Carouge le marquis de Varax de qui dépendait le
détachement d’Avusy. Ce Monsieur là me reçut fort mal ; il trouvait mauvais de ce qu’on envoyait un
jeune homme sans expérience, commander un poste qu’il jugeait de la plus haute importance, et poussa
l’honnêteté au point de me témoigner de l’envie de me renvoyer au corps, je lui dis que je ne
désirais pas mieux, que j’avais été très fâché de quitter la garnison ; il se décida cependant à me
laisser partir et me remit une lettre pour l’officier que je devais relever. En arrivant à Avusy, M.
Bourgeois me reçut fort bien, il me donna les plus amples instructions pour me tirer avec honneur
d’une commission que M. de Varax m’avait fait paraître fort périlleuse. M. de Launay gentilhomme
savoyard, établi à Avusy consentit que j’occupâsse le même appartement de mon prédecesseur et me
prit en pension. Sur le soir, Bourgeois me quitta, bien content de rejoindre la Bonneville, où la
Dumont l’attendait avec une vive impatience ; et me voici installé commandant d’un détachement d’une
vingtaine d’hommes ; ce titre me plut beaucoup, j’en soutenais le lustre avec grand soin, je ne
sortais jamais sans avoir mon ordonnance avec moi, et souvent en allant à la visite de ma troupe je
faisais prendre

\marginpar{\footnotesize\raggedright 1792.} les armes, j’observais une discipline sévère, qui ne
plut pas à mes vieux invalides. Les dimanches j’en menais une partie à Chancy, village genevois pour
assister au culte, et j’avais assez de confiance de ne pas craindre qu’ils me désertassent,
quoiqu’ils étaient sur terre franche. Je me trouvais assez bien dans la maison de Launay, mon
appartement était commode, les personnes de la famille me témoignaient de l’amitié, Mme surtout
m’accablait de carresses, femme de près de 40 ans avec un époux de 60 elle voulait tirer parti d’un
jeune homme de 18 ans, mais cela ne prit pas, tant par sa laideur que par ma timidité. J’aurais été
content de la pension, si je n’avais été obligé de manger continuellement maigre à cause du carême,
et ce qui me surprenait, était de voir assaisonner la salade avec du sucre en place de sel ; malgré
la scrupuleuse attention de M. de Launay de ne manger que du maigre chez lui, nous allions
quelquefois nous pro–mener à Passari une belle et bonne auberge sur Genève, où il se chargeait la
conscience d’un bon rôti de veau, me recommandant très fort de n’en rien dire à sa femme ; je crois
qu’il craignait qu’un autre fois elle voudrait aussi venir avec nous. Le premier mois se passait
agréablement à Avusy, je donnais des leçons d’écriture aux enfants de M. de Launay, je lisais des
livres d’un cabinet de lecture de Genève où je m’étais abonné, j’allais voir de temps en temps ces
dames de la Grave et le général de ce nom, qui demeuraient dans notre voisinage ; un mauvais cheval
de la maison était à ma disposition pour aller faire des promenades à Carouge et Genève, finalement
M. de Launay m’apprit à jouer au berlan, dans le dessein de me vider le gousset, je pris la passion
de ce jeu ; après dîner nous nous mettions quelquefois à cheval sur un tronc d’arbre qui était sous
un maronnier devant la maison en guise de banc, sans nous séparer jusqu’à nuit tombante ; d’autres
fois le soir venaient plusieurs riches paysans genevois pour faire la partie, alors on jouait gros
jeu, mais fort honnêtement, j’ai toujours eu le bonheur à m’en tirer avantageusement. La privation
de mes amis, le souvenir des fêtes et des plaisirs d’Annecy commencèrent à se faire sentir
fortement, ma solitude surtout me rappelait

\marginpar{\footnotesize\raggedright 1792.} le bonheur dont je jouissais au sein de ma patrie et de
ma famille ; et ces différents sentiments me causèrent l’étrange maladie du pays. Je résolus tout à
coup le 23 avril d’aller à Carouge demander à M. de Varax la permission d’aller faire une course
jusqu’à Morges ; je partis d’abord après dîner ; un officier de la légion légère grand ami de Mme de
Launay fut mis dans le secret sur son conseil, il devait me remplacer pendant mon absence. Arrivé à
Carouge je vis M. Hardisson, (nom de cet officier) qui se prêta de bonne grâce à ma proposition ;
mais il fallait avant tout l’agrément du commandant. M. de Varax me refusa tout net, il me dit : Si
vous voulez vous faire remplacer par un officier de votre corps, j’y consens. Comme je sentais fort
bien l’extravagance de mon projet, quoique j’avais dit à M. de Varax que des affaires pressantes me
demandaient à Morges ; je vis qu’il y fallait renoncer ; je sortis trè́s mécontent de M. le marquis
et m’ache–minais sur Genève ; il me vint une idée hardie de m’en aller sans permission, j’y trouvais
fort peu d’obstacles et je m’en fus louer un cheval pour partir le lendemain. Avant de quitter Avusy
j’avais fort recommandé au sergent du détachement (Dummermuth) de veiller sur mes gens, je leur dis
moi-même : mes amis tenez vous tranquilles jusqu’à mon retour il y aura pour boire, je pensais quand
même M. Hardisson viendrait à ma place, il s’occuperait plus de Mme de Launay que de mes soldats.
Tranquillement logé au Petit Maure, j’attendais avec impatience le matin pour me mettre en voyage ;
le temps me favorisa c’était le 24 courant le cheval était excellent à 2 heures après midi je fus à
Morges sans débrider ; à mon passage à Versoix on me fit un peu peur en me chantant le ça ira, et
les aristocrates à la lanterne, je piquais des deux pour l’échapper. Tout content d’être dans ma
patrie je me fis servir un bon dîner à la Couronne, après quoi je me rendis chez mon cousin
Wyttenbach receveur général des péages, où on me reçut à bras ouverts ; la manière dont j’étais venu
les étonna un peu, mais je leur donnais des paroles de consolation.

\marginpar{\footnotesize\raggedright 1792.} Ma réception au château chez M. le Bailli Ryhiner fut
moins satisfaisante ; ancien capitaine au régiment où je servais, M. de Ryhiner me fit plusieurs
questions, il s’apper–çut à mes réponses que jétais absent d’un détachement sans permission, il
n’eut rien de plus empressé que de me signifier de repartir sur-le-champ, et de ne plus faire de
pareils coups qui pourraient facilement m’attirer des arrêts de 6 mois dans une citadelle ; j’en
conçus beaucoup d’inquiétude, je passais tout le 25 à Morges, et le 26 je repris le chemin de mon
poste, j’y arrivais à 6 heures du soir – apprenant à ma grande satisfaction qu’il ne s’y était rien
passé de nouveau. M. de Varax m’avait fort recommandé de lui écrire souvent, afin qu’il sût si je me
trouvais à mon poste ; je me mis à mon bureau sur-le-champ pour lui mander que tout se passait ici
dans le plus parfait ordre. Est-ce un effet de sa bonté ? ou ne l’a-t-il pas su ? je l’ignore, mais
je n’eus jamais aucun désagrément de cette imprudence ; elle me guérit par contre du mal du pays.
Les derniers jours d’avril et une grande partie du mois de mai se passèrent à Avusy de la manière
dont j’ai parlé plus haut ; le maigre disparut et la table fut meilleure ; je payais pour les quatre
repas et le logement 2 louis d’or par mois. Le 18 de mai je suis allé voir les deux battaillons du
régiment qui avaient quitté Annecy pour se rendre à Saint-Julien près de Genève, je fis connaissance
avec des officiers que je n’avais pas encore vus, et je renouvelais celle des autres ; ils me
comblè́rent de nouveau d’honnêtetés ; le cousin Bucher me parut alors plus réservé, il avait une
jeune fille avec lui, et il hésita de m’inviter à souper redoutant l’influence des attraits de la
jolie Caterine. Le 20 je reçus ordre de partir avec mon détachement pour me rendre à Douvaine et de
là à Thonon, où se trouvait le 1er battaillon, depuis son départ de la Roche et de la Bonneville. Je
quittai Avusy sans regrets, le voyage se fit très heureusement, mon ordonnance et mon domestique en
même temps m’avait grossièrement désobéi ; pour le punir je le fis conduire par un caporal et quatre
hommes comme prisonnier jusqu’à la garnison ; l’officier qui m’avait relevé, et dont j’ai oublié le
nom, me donna ce conseil. Cette mauvaise conduite entraîna la disgrâce de ce domestique, Stettler le
remplaça par Witz qui nous servit assez longtemps et assez mal. J’arrivai à Thonon le 22 mai,
content de me retrouver avec mes amis et camarades ; Stettler me dit qu’il avait arrêté un logement
pour nous deux au couvent de la Visitation dans un plain-pied, j’y fis porter mes hardes.

\marginpar{\footnotesize\raggedright 1792.} Je trouvai à Thonon le jeune de Luternau, que j’avais
connu à Berne, il se conduisit à mon égard comme Graffenried, mais peu à peu il fut charmé de
s’humaniser davantage. Il y avait un battaillon des troupes légères de garnison avec nous, un
certain Chevalier Serra Lieutenant de ce corps m’a bien scandalisé par ses moeurs ; Reiberti était
par contre un brave garçon ; Gonzan et Cossato de bons enfants. Nous passions notre temps très
agréablement, on nous avait fait faire connaissance avec la maison Vignier Mangé, Stettler était
bien amoureux de Mlle Jenni ; nous recevions aussi des politesses de M. le baron des Étoles, et de
la maison du Boule ; M. le capitaine Jayet avait alors son épouse avec lui, nous y allions passer
souvent une partie de la journée. J’étais alors le plus lié avec Stettler, Watteville et Zehender,
Luternau partit pour Saint-Julien avec Alric l’aîné ; une de nos plus agréables jouissances de ce
temps était de boire de la bière ou du vin d’Arbois, en revenant de l’exercice, au café de la petite
Lyonnaise ; je m’amusais aussi à jouer au billard, M. le capitaine Lieutenant Nicole me gagna
quelques louis, cette perte me fit cesser le jeu pour quelque temps. Je crois que c’était le 26
juin, jour anniversaire du roi de Sardaigne, que nous fûmes invités par nos camarades de
Saint-Julien d’aller à un bal qu’ils se proposaient de donner ; Stettler et moi partîmes seuls de
Thonon, et nous n’eumes point de regrets, la fête était superbe ; on commença à danser à 9 heures du
soir, sous une magnifique tente, ornée avec beaucoup de goût ; des dames françaises, des Suissesses,
des Genevoises et ce qu’il y avait de mieux en Savoie, embellit ce bal, il était parfaitement bien
servi, et nos officiers tous mis avec la dernière élégance y brillèrent, on dansa jusqu’au jour.
Nous nous en retournâmes à Thonon, Stettler et moi, trè́s satisfaits de notre partie de plaisirs.
C’est dans le courant de juillet si je ne me trompe, que le Régiment bernois de Watteville sortit de
France ; nous apprîmes à Thonon le jour qu’il arriverait sur le territoire suisse, et nous allâmes
plusieurs attendre son entrée dans les États de la République de Berne, à Nyon. M. le trésorier de
Gingins

\marginpar{\footnotesize\raggedright 1792.} fut envoyé par Leurs Excellences pour recevoir ce
régiment. Le Souverain ordonna de traiter les officiers de ce corps, M. de Gingins nous invita à ce
repas, il était bon, mais il ne répondait par à ce que l’on attendait ; pendant le dîner je faisais
tous mes efforts pour ne pas dormir, ayant passé une nuit blanche sur le lac j’avais de la peine à
me défendre du sommeil, au dessert il me passa, et nous nous réveillâmes tous pour faire du
désordre, on jeta assiettes, verres, bouteilles, couteaux et tout ce qui se trouva sur la table, par
la fenêtre ; j’avoue à ma honte que j’ai excellé dans cette manoeuvre. M. Stettler capitaine dans
Watteville me fit d’amicaux reproches de ce que j’avais augmenté les frais du souverain – je lui
répondis : Plus le compte sera gros, plus Leurs Excellences croiront que nous avons été bien
traités, et plus Elles auront de plaisir. On parla même à Berne du tapage qui s’était fait, et on
eut la bonté de me citer comme un des premiers auteurs. Nous nous rembarquâmes sur le soir et nous
sommes arrivés de nuit a Thonon. C’est à peu près dans ce temps que je fus commandé en détachement,
pour allez relever Watteville qui était à Nernier, petit village au bord du lac de Genève et
vis-à-vis de Nyon. L’ami Stettler prit un cheval, et m’accompagna à mon poste ; avant de partir, M.
le lieutenant-colonel Stettler, qui commandait le battaillon, me défendit sous peine de trois mois
d’arrêts de ne jamais aller à Nyon. Aussitôt arrivés à Nernier, Watteville me remit les consignes,
nous dînâmes et après il fut résolu que nous prendrions un bateau pour aller faire une tournée en
Suisse. Pour éviter le châtiment qui m’était réservé en cas que la chose parvint aux oreilles de M.
Stettler, je mis l’habit de chasse de Watteville et je me fis passer pour le domestique de ces
Messieurs. J’ai si bien pris la tournure de mon rôle, qu’en commandant le goûter à l’auberge de la
Couronne à Nyon, un jeune valet m’accosta et me prit pour un de ses camarades il a même fallu toute
l’autorité de mes maîtres supposés, pour éviter que la dispute ne devînt sérieuse, ce jeune homme ne
voulait pas se laisser

\marginpar{\footnotesize\raggedright 1792.} détromper, le tour nous fit beaucoup rire. Mes amis s’en
retournèrent à Thonon et moi à Nernier ; M. le colonel n’a rien su, ou n’a rien voulu savoir de mon
imprudence. Arrivé à mon poste je me mis au lit, je ne pus fermer l’oeil, le matelat était d’une
dureté horrible ; le matin le sergent vint au rapport pour prendre mes ordres, je lui dis de
commander huit hom̃mes pour venir défaire mon matelat, ils vinrent, arrangèrent le crin, le
recousirent, et je dormis à merveille cette nuit. Quelques jours aprè́s je me mis sur eau avec mon
sergent, deux caporaux et trois soldats, le vent vint si fort qu’il nous poussa jusqu’à Rolle, je
fus dîner avec mon sergent et je donnais aux autres de l’argent pour boire. Sur le tard je voulais
repartir, mais le lac se trouva innavigable ; je courus dans les auberges pour ramasser mes gens,
qui, se trouvant sur terre libre, refusent de m’obéir ; le sergent se donna pendant longtemps
d’inutiles peines, il y réussit à la fin et nous prîmes la route de Nyon espérant d’y trouver un
bateau pour nous conduire à Nernier, nous fûmes obligés de laisser le nôtre à Rolle. Sur le chemin,
le caporal Desgouttes, qui était ivre, faisait l’insolent, je le menaçais de lui passer l’épée à
travers le corps, il devint furieux et voulut m’attaquer, mais les autres me défendirent sabre à la
main. Arrivés à Nyon de nuit, j’eus toutes les peines de trouver un batelier qui voulut me ramener à
Nernier, à la fin il s’en présente un ; en entrant dans le bateau je tombe, et je me crotte de tout
mon long, le vent était toujours fort ; nous débarquâmes après les 10 heures du soir, le vieux
soldat qui me servait m’annonça à ma grande satisfaction que pendant son commandement il ne s’était
rien passé de nouveau ; je fis d’abord mettre mon caporal au fer ; le lendemain j’envoyais prendre
le bateau que j’avais laissé à Rolle, il arriva le soir fort heureusement et cette course
inconsidérée n’eut aucune fâcheuse suite. Après huit jours de séjour dans ce malheureux village, où
l’ennui m’aurait fait commettre d’autres extravagances, je reçus l’ordre de me rendre à Douvaine, à
la fin de juillet, sous les ordres du capitaine Bégoz qui y commandait. J’étais fort content du
change ; Douvaine est un grand village sur la route de Thonon à Genève ; je fus logé à l’auberge où
j’étais assez bien.

\marginpar{\footnotesize\raggedright 1792.} M. Bégoz m’envoya à Thonon prendre le prêt, je pris un
cheval, à mon retour j’attachais derrière la selle l’argent qui était tout en petite monnaie, le sac
s’ouvrit et je semais ainsi une vingtaine de livres tournois. Un jour mon capitaine me dit qu’il
avait envie d’aller en ville ; je voulus profiter de son absence pour voir Mlle Marion sa petite
maîtresse ; mes vues étaient innocentes dès le commencement, peu à peu une chose donna l’autre,
enfin l’humanité était sur le point de vaincre la raison, on se défendit tant soit peu, mais au même
moment entra le domestique de la maison qui nous surveillait, et le beau projet fut en poudre ; je
ne fis plus de tentative à l’avenir. M. Guiot Syndic du village chez lequel je dînais de temps en
temps, me procura la connaissance de M. Quisard de Massongy ; par le secours de quelques bonnes
bouteilles de vin, ce brave Savoyard devint mon ami, il m’invita à venir le voir à sa campagne qui
n’était qu’à une demi-heure de Douvaine ; je m’y rendis et n’eus pas lieu de m’en repentir ; hélas !
c’est à Massongy que mon coeur conçut les premières impressions du plus vrai amour, la charmante
Julie fille cadette de M. Quisard m’enchaîna dès l’instant que j’eus le bonheur de la voir, si
j’avais dû peindre alors l’idole de mon coeur, chaque trait aurait trahi l’impétuosité de mes feux ;
aujourd’hui plus calme et moins épris je pourrais être vrai. Julie avait 17 ans et toute la
fraîcheur de cet âge, animée par de grands yeux bleus qui embellissaient une figure charmante, sa
taille était svelte, sa belle chevelure d’un blond cendré, sa voix, ses traits formaient le plus
aimable tout. Le marquis de Seyssel, jeune étourdi, officier aux dragons de Chablais, était peu
avant moi le volage adorateur de tant de charmes ; j’en parlais à Julie, elle m’avoua qu’en effet le
marquis avait fréquenté la maison, mais qu’on ne pensait guère plus à ceux qui nous oublient. Je fis
tous mes efforts pour gagner le coeur de mon aimable amie, je courti–sais le père et la mère
(seconde femme de M. Quisard et belle mère de Julie), je réussis à merveille, on me pria de
m’envisager comme le fils de la maison ; et il y avait toujours un couvert pour moi à leur table.
Tous les soirs je m’en retournais à mon poste, enivré d’amour ;

\marginpar{\footnotesize\raggedright 1792.} mon vieux soldat s’apperçut facilement qu’il s’était
opéré une révolution en moi depuis mon séjour à Nernier ; il me dit un jour très naïvement :
Monsieur vous êtes sûrement amoureux, sans cela vous ne seriez pas sans cesse dans cet autre
village, je lui en fis l’aveu, et depuis ce temps il s’appliqua infinement mieux à mes affaires, il
nettoyait mieux les habits, il cirait les bottes avec plus de soin, et quand je ne me mettais pas à
sa fantaisie, il me disait : M. ceci ne plaira pas à votre belle. Dans cette époque c’est-à-dire au
mois d’août 92 on parlait déjà de la guerre de la France avec le roi de Sardaigne, les officiers
furent prévenus de prendre les mesures nécessaires pour entrer en campagne en cas de besoin. La
première chose fut d’acheter des chevaux ; je ne voulais par me presser d’abord, mais l’affaire du
10 Août qui a épouvanté l’Europe et mis les Suisses au désespoir, me fit envisager la guerre comme
inévitable, et je crus ne plus pouvoir me dispenser de l’achat d’un cheval ; j’en fis l’acquisition
pour 12 louis. Fort peu en peine des événements politiques, je m’occupais plus de ma chère Julie que
du sort des États ; mon joli petit cheval que je menais bien, ayant passé six mois au manège,
m’emporta tous les après-midi chez M. Quisard ; j’étais alors l’homme du monde le plus heureux, je
jouissais du rare bonheur d’aimer et d’être aimé ; je pressais souvent ma Julie contre mon sein et
lui donnais des baisers ardents sur la bouche, quelque–fois même je portais ma main hardie sur sa
gorge charmante, mais le reproche amical : laisse-la, elle se flétrit, commandait le ménagement.
Quoique initié dans le secret des jouissances défendues, l’idée ne m’en vint jamais, les faveurs que
je goûtais me paraissaient délicieuses, je ne croyais point à la possibilité d’en découvrir de plus
douces, et un sentiment indéfinissable m’inspirait du respect pour un fruit que d’autres auraient
cueilli avec transport. Un jour que des badinages innocents nous procuraient toujours de nouveaux
charmes, Julie se mit à courir, je la suivis sans pouvoir l’attraper, elle se sauva enfin dans un
petit cabinet, et se jeta tout en riant sur le lit, je la joignis et fermais la porte. Dans ce
moment elle était absolument en mon pouvoir ; je me glorifie encore aujourd’hui du mouvement
vertueux qui me saisit alors, Julie sentait fort bien le péril auquel elle s’était innocemment
exposée, je vis son trouble, et lui dis : comment ! ma chère amie pouvez-vous vous méfier de moi,
croyez-vous que je veuille ainsi abuser de tout ce que j’ai de plus cher au monde ? je la relevais
du lit, et nos larmes coulèrent. J’avoue

\marginpar{\footnotesize\raggedright 1792.} J’avoue qu’aujourd’hui en me rappelant ma conduite, elle
me paraît étrange, je ne sais à quel sentiment l’attribuer, était-ce vertu ? était-ce timidité ? je
l’ignore. Il est certain qu’échappés à cette tentation nous nous en aimions davantage, mais on vint
tout à coup troubler ce bel amour. M. Quisard avait pendant longtemps été le spectateur tranquille
de nos beaux feux, il ne paraissait en concevoir aucune inquiétude ; confiant, amical, et prévenant
avec moi, je commen-çais à me douter qu’il s’imaginait peut-être que je pourrais avoir des vûes de
mariage sur sa fille. Un dimanche Julie m’avait donné rendez-vous au jardin, elle me dit le samedi
soir : Venez demain matin de bonne heure ; tout le monde sera à la messe, moi j’en sortirai, et je
viendrai vous joindre au cabinet du jardin, où nous serons seuls. Je me rendis à mon poste à l’heure
indiquée ; j’attendais, j’attendais, mais en vain ; la messe finit, et au lieu de Julie je vois
arriver le père d’un air réfléchi ; je lui dis M. je vous attendais ici, pensant que vous y
viendriez au sortir de l’église ; il me fit sentir qu’il savait fort bien que je n’étais pas là pour
lui, et me tint ce surprenant discours : Monsieur, je vois naître une inclination dans le coeur de
ma fille, mes sollicitudes paternelles veulent savoir quel est votre but en la lui inspirant, je
vous crois trop honnête homme pour en redouter un, qui ne fut noble, cependant les jeunes gens de
votre âge, et de votre état, s’abandonnent quelquefois imprudemment à leurs passions. Pendant le
temps que vous avez fréquenté ma maison j’ai appris à vous connaître favorablement, au point que je
vous offre ma fille en mariage, persuadé que vous la rendriez heureuse ; j’ai 100 mille livres
tournois de Savoie qui appartiennent à mes deux enfants, je vous en donne la moitié -si vous voulez
épouser ma fille ; ce parti ne vous convient-il pas, je suis forcé de vous prier de ne plus mettre
le pied dans ma maison ; pardonnez une franchise que l’amour paternel et mes devoirs de père
m’imposent. M. Quisard s’arrête et attend avec impa–tience ma réponse, et voilà à peu près ce que je
lui dis, peut-être avec plus d’embarras que je le retrace ici : il est vrai Monsieur, que j’aime
Mlle votre fille, et je me trouverais heureux de couronner mon amour par les liens

\marginpar{\footnotesize\raggedright 1792.} du mariage, mais l’injuste rigueur de mon sort y met des
obstacles insurmontables ; je lui dis alors que les lois de mon pays me faisaient perdre mes biens,
ma bourgeoisie, et mon état en épousant une femme catholique ; cela suffit pour le convaincre que
cette union ne pouvait se faire. Là-dessus M. Quisard me pria de me soumettre à ce qu’il m’avait
dit, je suis fâché Monsieur de me séparer de vous, ajouta-t-il, portez vous bien, oubliez Julie, et
pensez quelquefois à un ami qui vous aime et vous estime. Je lui promis d’obéir, mais je voulais
voir mon amante pour la dernière fois, je ne pus l’obtenir, et je partis de Massongy vers les midi
sans avoir vu une âme de la maison, que le maître. Je m’en retournais à Douvaine, stupéfait de tout
ce qui venait de se passer, ma raison approuvait ce digne père, tandis que mon amour le condamnait ;
après quelques heures de réflexions, je résolus de faire mon possible d’oublier Julie. Pusieurs
jours avant ce bizarre entretien j’avais acheté à Genève quelques petits bijoux que je destinais à
ma belle, j’en disposais d’abord autrement et me proposais de les donner à ma soeur D. quand j’irais
en Suisse. Ce qui contribua le plus à me rendre le sacrifice de mon amour moins pénible, fut la fin
de mon détachement ; on vint me relever au commencement de septembre rendu à mes amis, au milieu
d’une très agréable garnison, j’oubliais les moments délicieux que j’avais passés auprès de Julie,
et je ne l’ai plus revue, depuis ce certain samedi soir, mais dans ce moment encore il m’en reste un
doux souvenir, et je me suis arrêté longtemps à ce court espace de ma vie, parce qu’il a été un des
plus heureux de ceux que j’ai parcourus. Peu après mon retour à Thonon je demandais à M. le
lieutenant colonel Stettler la permission de m’absenter pour quelques jours, il me l’accorda, je me
proposais de m’embarquer à Évian, et d’aller faire une tournée à Lausanne chez ma tante Porta, soeur
de mon père ; j’arrivais à Évian sur le soir, et je commandais un bateau, le lac était trop agité
pour le passer et une pluie abondante en faisait passer l’envie ; je m’en fus à la grande auberge où
je demandais à souper ; à table d’hôte je fis connaissance avec un abbé français honnête au-dessus
de toute expression ; il me dit qu’il avait aussi dessein de s’embarquer pour Lausanne, et que nous
ferions le voyage ensemble ; en nous retirant, il ordonna au domestique de nous donner une chambre à
deux

\marginpar{\footnotesize\raggedright 1792.} lits, puisque nous devions faire le même voyage le
lendemain, et me demanda, n’est-ce pas Monsieur ? je répondis : avec plaisir. Arrivé à notre
appartement il ferma la porte à clef, et commença une conversation indifférente, il ouvre un buffet
d’où il sort une tabatière remplie de mor-ceaux carrés de bon chocolat d’Alphonse, il m’en offrit et
j’en pris. Le couvercle de cette boîte se détacha, l’abbé me fit voir une peinture infâme, qui
m’ouvrit les yeux sur l’homme avec lequel je me trouvais ; je lui dis, voilà ce que je n’aurais par
cherché dans la poche d’un ecclésias–tique ; comment ? me répondit-il avec enthousiasme, ne
sommes-nous pas des hommes comme tous les autres ? il voulait me faire des caresses et m’engager de
coucher avec lui ; je me saisis de mon épée le menaçant de la lui passer à travers le corps, s’il ne
me laissait tranquille, je la pris même nue dans mon lit. Toutes ces mesures déconcertèrent les
sales et abominables projets de M. l’Abbé, il me dit : Que vous faites l’enfant Monsieur le
Chevalier on voit bien qu’il y a peu de temps que vous vivez dans le monde ; je ne répondis rien et
me tenais tranquille dans le lit, je dormais fort peu, le matin je me levais de bonne heure, je
sortis de la chambre le comblant d’injures. Le lac était encore impraticable, je me rendis chez M.
Bourgeois qui commandait le détachement d’Évian, je lui fis le récit de mon aventure, il se mit à
rire comme un fou. Allons déjeuner là-dessus, me dit-il, et bien copi–eusement, car nous ne dînerons
qu’après 3 heures, je veux te mener mon cher ami, chez M. Benford, un Anglais qui a un million de
rentes, et chez lequel ton abbé ferait peut-être fortune, je crois qu’ils ont les mêmes goûts. Nous
y fûmes en effet ; avant dîner nous entendîmes une dame de cour d’une princesse russe, qui était
venue de Lausanne ce jour-là, avec le prince Camille et un grand monde, toucher du clavecin et
chanter comme un ange ; elle était belle comme l’amour. On servit le dîner qui fut splendide, après
il y eut concert, sur les cinq heures nous montâmes dans un cabriolet

\marginpar{\footnotesize\raggedright 1792.} que deux chevaux arabes emportèrent comme au vol au bois
d’Évian, où sous une magnifique tente tout ce beau monde se réunit pour boire le thé. Benford ne
paraissait nullement jouir de tant de délices, il en témoignait la plus grande indifférence. Après 6
heures on s’en retourna à Évian, je fis ma révérence à Benford, il dit à Bourgeois, vous savez,
commandant, que votre couvert est toujours mis. Je passai la nuit chez Bourgeois, beaucoup plus
tranquillement que chez l’Abbé, nous parlions pendant plusieurs heures de tout ce que nous avions vu
de beau dans la journée ; le lendemain je repris le chemin de Thonon, ma permission était déjà plus
qu’à moitié écoulée. Je passais les jours de septembre très agréablement ; mon cheval me prouvait
surtout beaucoup de plaisirs ; Stettler en avait aussi reçu un de Suisse, et tous les soirs nous
faisions de charmantes courses. Pendant notre séjour à Thonon, nous sommes allés une fois chez les
Chartreux à Ripailles, où nous avons fait un excellent repas servi en maigre. Le 22 de ce mois nous
apprîmes tout à coup la nouvelle de l’invasion de l’armée française en Savoie ; les troupes du roi
de Sardaigne n’avaient fait aucune résistance, elles se retiraient en grande hâte dans le Piémont,
par le passage presqu’inpraticable des Bauges. Notre battaillon, et les deux compagnies de la légion
légère, qui com–posaient la garnison de Thonon attendaient des ordres avec la plus grande impatience
; il n’en vint point et le danger approcha. Nos commandants résolurent de partir pour Meillerie et
d’y prendre une position ; arrivés à ce village, dépourvu de tout, on envoya à Vevey cher–cher des
vivres. Deux capitaines et deux subalternes furent détachés à Évian pour observer les mouvements qui
se faisaient de ces côtés ; le capitaine Muret arriva dans cette ville, il trouve le sénat
rassem–blé, et les membres décorés de la cocarde tricolore ; armé d’un pistolet il s’approche du
président et lui ordonne de poser ce signe de rébellion, et d’en signifier autant aux assesseurs ;
ils obéirent tous, et Muret les harangua les exhortant à l’attachement et à la soumission de leur
légitime souverain, jusqu’à ce qu’une force impérieuse en disposerait autre–ment.

\marginpar{\footnotesize\raggedright 1792.} Dans le même temps, je fus envoyé avec le chevalier
Cossato,lieutenant de la légion, du côté de la Bonneville, pour aller à la découverte de l’ennemi ;
Cossato quoique mon aîné me pria de commander ; nous parcourûmes les vallées qui séparent le
Chablais du Faucigny. Arrivés de nuit à Bonne un grand village à deux petites lieues de Bonneville,
nous descendîmes chez le curé, qui nous dit qu’il serait dangereux de se porter en avant. J’envoyai
un homme de confiance à Bonneville, il me rapporta que tout s’y préparait à recevoir les Français,
qui étaient attendus le même jour. Le 24 je résolus de partir avec mon compagnon pour instruire les
capitaines qui étaient à Évian de ce que je venais d’apprendre. Un abbé nommé Gerdil nous accompagna
jusqu’à une lieue d’Évian, où il se sépara de nous. J’ai toujours souffert des maux de dents pendant
cette course. En entrant dans cette ville nous trouvâmes tous les bourgeois sous les armes, avec la
cocarde nationale ; la garde arrête mon cheval et me dit, voilà encore un de ces insolents
Piémontais qui se sauve, je leur dis, Messieurs je suis Suisse, par conséquent libre de naissance,
j’ai servi votre roi avec fidélité je le servirais encore et je veux aller joindre mon corps ; il y
a longtemps qu’il a foutu le camp me répondit- on, faites-en autant et partez ; je les priai de
laisser passer de même mon compagnon de voyage, ils y consentirent après lui avoir dit un tas de
sottises. Nous piquâmes des deux, heureux d’avoir échappé, en arrivant à un bois peu distant de
Meillerie, nous rencontrâmes huit hommes munis d’énormes bâtons qui venaient à nous en courant, ils
donnèrent de grands coups sur la croupe de nos chevaux et nous laissèrent courir je crois que
c’étaient des bandits échappés des prisons. Arrivés à Meillerie on nous dit que notre troupe s’était
embarquée le jour au–paravant ; fort indisposés contre nos chefs, qui nous avaient ainsi abandonnés,
nous cherchâmes un bateau pour aller à Vevey où nos corps s’étaient rendus ; après bien des peines
nous pûmes nous embarquer, mais sans nos chevaux ; qu’un homme s’offrit à retirer jusqu’à ce qu’il
se présenterait une occasion pour nous les faire passer en Suisse. Une pluie affreuse nous surprit
sur le lac ; mon pauvre Cassato, qui n’avait jamais été en bateau, était dans des angoisses

\marginpar{\footnotesize\raggedright 1792.} horribles. Nous arrivâmes enfin après les 8 heures du
soir à Vevey, tout trempés ; on nous conduisit au château, chez M. le Bailli de Watteville, où se
trouvaient nos chefs, ils furent très contents de nous revoir sain et sauf on leur avait dit que les
patriotes nous avaient lanterné à notre passage à Évian. Il est vrai que ce n’est pas leur faute si
nous ne l’avons été. On voulut nous faire asseoir au château, mais nous étions dans un état qui ne
nous permit point de salir les beaux meubles du bailli. L’aide-major vint me dire que j’étais invité
à souper chez M. Roux, je m’y rendis, et on me traita parfaitement bien ; je devais aller loger chez
M. Hugonin à la Tour, après souper un domestique m’y conduisit, et dans cette maison je fus aussi
supérieurement traité. Le lendemain 26 septembre toute la troupe partit pour se rendre à Bex ; les
deux colonels d’Antignan et Stettler s’établirent au château d’Aigle, chez M. de Diessbach, où ils
auraient volontiers fini la campagne ; le sort me plaça avec M. le Major Zehender, chez le maître de
poste Riccoux, un vrai galant homme qui nous combla de bienfaits ; j’ai oublié de dire qu’à notre
passage à Aigle la ville traita le corps d’officiers et donna du vin aux soldats. Leurs Excellences
de Berne ordonnèrent au régiment de continuer sa route pour rentrer dans les É́tats du roi de
Sardaigne ; le canton du Valais après bien des difficultés, nous permit de traverser ses États, en
ôtant la platine à nos fusils, ce qui ne fut pas rigoureusement observé. Les recruteurs espagnols
nous débauchèrent une infinité de soldats qui jusqu’à ce moment ont été fidèles à leurs drapeaux. Le
1er octobre nous passâmes le Saint-Bernard. Je ne veux point passer sous silence, un plaisir inouï
que j’ai eu sur la route de Villeneuve à Aigle ; fort inquiet sur le sort de mon cheval, j’ai fait
partir une barque d’abord à mon arrivée en Suisse, pour aller le prendre, ainsi que celui de
Cossato, je m’impatientais infiniment du retour de mon domestique, et au moment ou j’étais très
fatigué de faire la route à pied, on m’amène mon cheval très bien portant. Les pères charitables de
l’hospice du Grand-Saint-Bernard reçurent fort bien les 600 hommes, jamais depuis l’existence de
leur généreux établissement ils n’avaient eu un aussi grand nombre de passagers à la fois, ils nous
traitèrent au mieux ; le

\marginpar{\footnotesize\raggedright 1792.} soldat a fait preuve dans cette circonstance, que le
sentiment d’honneur et de la recon– naissance ont rarement l’accè́s de son coeur ; ils ont volé à
ces généreux prêtres tout ce qu’ils ont pu emporter, en pain, fromage, et autres, les Piémontais
surtout se sont distingués, mais le roi a vengé une aussi vile conduite en donnant à l’hospice un
dédommagement considérable, et bien au-dessus des sacrifices qu’ils avaient faits. Nous arrivâmes le
3 octobre à la cité d’Aoste où se trouvaient tous les fuyards de l’armée piémontaise ; les officiers
avaient perdu leurs équipages, ils étaient bien déguenillés ; nous plaignîmes beaucoup ceux du
régiment et chacun de nous s’empressa de les soulager de son mieux. Après quelque séjour à Aoste on
mit l’armée en cantonnements ; le 3ème battaillon dont je faisais nombre, fut destiné à Courmayeur,
les restes du régiment restè́rent dans des villages voisins plus près d’Aoste. Nous nous trouvions
passablement bien dans notre nouvelle garnison les officiers étaient tous logés dans la grande
maison qui servait en été aux personnes qui venaient boire les eaux de Courmayeur qui sont
sulphureuses et très salubres. Un salon vaste nous réunissait dès le commencement du jour, chacun
faisait à la cheminée son café, on dînait ensemble et le soir plusieurs parties de jeux servirent à
nos amusements. Je fus arraché à ce joli séjour pour aller en détachement à Combales une montagne
qui séparait le Piémont de la Savoie tout près du Mont Blanc, sous les ordres de M. le
Capitaine-lieutenant Bergier. C’était au milieu de novembre le temps était mauvais, les monts
couverts de neige, et le poste ne présentait d’autre refuge qu’une mauvaise baraque de bergers, qui
ne mettait nullement à l’abri des injures de l’air. En entrant dans ce triste lieu, je vis le
Capitaine-lieutenant Muret suspendu sur un matelas, qui était attaché aux poutres du toit par de
longues cordes ; il se fit descendre pour nous recevoir, et nous remit le poste avec beaucoup de
plaisir ; j’y passais huit jours, aussi désagréablement que possible M. Bergier me faisait
patrouiller au milieu de la nuit, par des rocs impraticables, son air malpropre, sale, et puant
n’harmoniait guère avec moi, j’attendais impatiemment l’époque de ma délivrance, elle arriva enfin,
et je me trouvais fort heureux de m’en retourner à Courmayeur – le tour de fatigue me toucha
quelques temps après, mais j’engageais

\marginpar{\footnotesize\raggedright 1792 et 93.} un enseigne du régiment de le faire pour moi, en
lui payant une pièce de 24 livres tournois. C’est à Courmayeur que j’eus pour la première fois un
différend avec mon ami Stettler, il crut comme mon aîné en service, quoique du même grade, avoir le
droit de choisir son appartement sur le nombre, qui était destiné aux subalternes, son oncle qui
commandait alors, décida contre lui, on tira au sort, et nous nous raccommodâmes. En décembre 1792
le régiment envoya le Capitaine Bucher en recrue à Berne ; tant par complaisance pour mes parents,
que par amitié pour moi, il me demanda au Brigadier Rochmondet, pour l’aider au recrutement. J’en
avais témoigné la plus grande envie à M. le Lieutenant-colonel Stettler, qui me badinait beaucoup,
de ce que je voulais aller en Suisse, après un séjour de 11 mois au corps, tandis que d’autres
officiers ne demandaient des semes–tres que le 3ème hiver ; je sentais la justice de ses
raisonnements et je renon–çai à mes projets ; une lettre du chef du corps qui ordonnait à M.
Stettler de me laisser partir, me mit au comble de ma joie. Je quittai Courmayeur le 7 du mois,
Messieurs Bucher et Mazars m’attendaient à Aoste, pour nous mettre en voyage ensemble ; nous
partîmes le 10 j’eus beaucoup à souffrir des sar–casmes du Capitaine Mazars, qui était fort jaloux
qu’un jeune officier bernois s’en allât en semestre, ou en recrue, après n’avoir passé qu’un jour au
régiment comme il disait ; il me tourmentait toujours, mais le plaisir de revoir mes chers parents
effaçait ces petits chagrins ; le cousin Bucher était en échange plus complaisant, plus honnête,
plus amical que jamais. Nous arrivâmes à Berne le 16 décembre Mazars nous quitta à Lausanne à mon
grand plaisir, quoique sur la fin il fut un peu plus doux. On m’accueillit à la maison paternelle
avec bienveillance, et je me croyais l’homme le plus heureux du monde ; je m’aperçus cependant bien
vite que onze mois d’ab– sence n’avaient pas suffi pour détruire l’opinion peu favorable que le
public avait conçu à mon égard, comme je l’ai dit plus haut ; mes anciens amis me traitèrent encore
en froid ; un jour me présentant en uniforme dans une société de dames des plus brillantes de Berne,
on eut la grossièreté de ne me point offrir de carte pour le jeu, aussi n’y remis-je pas le pied. Au
commencement de 1793 mon frère aîné me proposa de m’introduire dans sa société j’y consentis avec
plaisir, et m’y trouvais fort bien ; c’était ma 1ère

\marginpar{\footnotesize\raggedright 1793.} entrée dans le monde, en Suisse. À peu près dans ce
temps arriva mon ami Stettler, qui avait reçu un congé de quelques mois ; j’en eus une vraie
satis–faction ; l’espérance de me reconcilier par sa médiation avec mes anciens amis, commençait à
renaître, mais ses soins n’eurent pas les succès qu’on en attendait. Pour mon malheur je fréquentais
beaucoup mon ami de Luternau qui était à Berne, comme prisonnier de guerre, et c’est lui qui
m’entraîna au libertinage ; presque toutes les soirées se passaient dans les arcades où nous étions
à la poursuite des créatures, que nous faisions aller dans des mauvais lieux pour les voir
librement, quelquefois B. de Graffenriéd et mon ami Stettler, et souvent Sinner, étaient de la
partie, je voyais ce dernier très fréquemment. Je passais mes matinées chez mon cousin Bucher, où
les affaires du recrutement m’occupaient un peu, j’allais l’après-midi chez ma soeur D. faire une
agréable lecture. J’achetai dans ce temps-là un beau cheval normand pour 18 louis d’or mon père me
permit de faire venir sa chaise en ville pour varier mes promenades à cheval ou en cabriolet,
Luternau avait aussi son équipage, ainsi nous étions fort souvent en course. Le petit cheval que
j’avais acheté en Savoie était ruiné, je l’ai confié à mon départ à un officier du régiment, pour le
remettre, pendant mon absence. Un séjour de plusieurs semaines chez mon beau-frère à Schupfen au
commencement de mars, m’a un peu distrait de mes dissipations, je pris tout à coup la résolution de
partir de là, après avoir passé une nuit remplie d’inquiétudes ; des symptômes de vérole les ont
causées. Je quittai Schupfen assez brusquement pour aller consulter un chirurgien à Berne ; mes
hôtes ne pouvaient s’imaginer les raisons d’un départ aussi précipité. Mon mal ne paraissait pas
aussi grave que je le craignais, M. Théolier, chirurgien-major du ci-devant régiment de Watteville
m’assurait que ce n’était qu’un échauffement, et me prescrivit une espèce de baume. Dans ce temps
l’ami Stettler et le jeune Wyttenbach, un nouvel officier, résolurent de joindre le corps ;
Luternau, Sinner, et moi voulûmes les accompagner jusqu’à Fribourg ; on partit à cheval, et on fit
grande ribotte avant de se séparer ; mon mal empira tellement que je sentis des douleurs aigües à
mon arrivée à Berne. Le 6 du mois d’avril un

\marginpar{\footnotesize\raggedright 1793.} incendie éclata dans notre maison, le toit et tout ce
qui se trouvait sur le galetas devint la proie des flammes. Occupant une chambre du 3ème étage, je
crus prudent d’en sortir les effets, je les mis tous dans le lieu des privés où ils furent hors de
danger, on rit beaucoup de cette singulière idée. Les réparations à faire à notre maison
nécessitèrent mes parents d’entrer dans une autre, où on me donna le plain-pied, par ce moyen je me
croyais à même de soigner mon incommodité à l’insu de mes parents, mais leur surveillance trompa mes
espérances ; ma mère voulut absolument connaître le sujet de ma maladie elle en conçut des motifs
bien plus graves que ceux qui existaient en effet, et mon médecin ne put se refuser à lui tout
découvrir. Dès le moment que j’appris que mes parents connaissaient la source de mes maux, j’en
étais vivement affligé, et je souffrais plus moralement que physiquement ; mon père me fit des
reproches paternels, qui firent la plus forte impression sur mon coeur. Au bout de huit jours je fus
radicalement guéri, et je pressai le cousin Bucher de faire en sorte que nous puissions bientôt
partir pour le régiment, il comprit fort bien le sujet de mon empressement et le départ fut fixé au
1er mai. Le chirurgien m’apporta son compte qui pour un traitement de dix jours, se montait à six
louis d’or, j’eus recours à mon père pour satisfaire ce Juif, et je me souviendrais toute ma vie,
qu’en sortant cette somme du bureau, mon père me la jeta sur la table et me dit : „Voilà de l’argent
que je regrette.” Ma plume s’arrête longtepms à la description de ce fâcheux événement de ma vie ;
mon coeur se rappelle vivement des tourments qui en ont été la suite et il s’applaudit de l’heureux
effet qu’il a produit ; c’est à peu près depuis cette époque, où instruit par le malheur, j’ai
renoncé au libertinage. Je partis de Berne avec moins de regrets que les autres fois ; le voyage se
fit agréablement avec le cousin Bucher et un jeune Gruber qui entra au régiment. Après huit jours de
marche nous arrivâmes heureusement à Ivrée où notre corps se trouvait en cantonnement, les trois
voyageurs se logèrent ensemble ; nous mangions alors chez Mme Tignon au Lion d’or, à la meilleure
table que j’ai connue, pour 3 louis d’or dîner et souper. Le temps s’écoulait dans les plaisirs
pendant le séjour ; plusieurs jeunes gens entrèrent alors au régiment ; il était d’usage de traiter
le corps d’officiers et nous eûmes ainsi plusieurs excelletns repas ; l’exercice matin et soir, des
pro–menades à cheval, et des petites courses à Chavran, étaient nos occcupations journalières, et
chacun faisait ses petits préparatifs pour entrer en campagne.

\marginpar{\footnotesize\raggedright 1793.} Le 25 juin, Messieurs de Luternau, le ministre Ris et
moi résolûmes d’aller à Turin pour voir le lendemain la fête de la naissance du roi, qui se
célébrait toutes les années avec beaucoup de pompe. Je fus épris du faste de la cour et de la beauté
rare de cette capitale, nous y passâmes deux jours et nous vîmes tout ce que la ville offre de
remarquable. De retour à Ivrée le régiment reçut l’ordre de partir pour la vallée d’Aoste ; nous
sortîmes de notre garnison avec 2000 hommes. Le régiment était superbe, au grand complet et presque
tous Suisses. On le destina à camper au Freney, un bois situé à la pente des glaciers du Mont Blanc
; M. Rochmondet qui avait le commandement des troupes stationnées de ces côtés, me nomma son aide de
camp. Deux mois se passèrent dans ce camp, sans que rien d’intéressant survînt ; je me trouvais fort
bien à la table du Brigadier et je n’avais d’autres chagrins que ceux que me causèrent la jalousie
de plusieurs de mes camarades qui ne voyaient pas avec plaisir la bienveillance que m’accordait le
chef. Au milieu du mois d’août une forte dysenterie me força de quitter le camp ; je me fis
transporter dans un village voisin de Courmayeur à Antraigues. Au commencement de septembre se fit
l’invasion de nos troupes en Savoie ; je n’eus pas le bonheur de partager la gloire que le régiment
s’acquit alors, je ne pus le rejoindre qu’à la fin de septembre avec les deux colonels Tschiffeli
qui étaient restés malades à Aoste. Arrivé à Moûtiers où se trouvait le quartier général de Son
Altesse Royale le duc de Montferrat, et du Général autrichien d’Argenteau, le Brigadier Rochmondet
qui avait été promu au grade de Major-général me reçut froidement, et me dit que ma longue absence
l’avait obligé de se pourvoir d’un autre aide de camp, je pris mon parti, et je fis dresser ma tente
dans le camp qui était devant la ville. Quelques gardes aux avant-postes entre autres celle avec
Jonquière l’aîné à la Chapelle en avant de Moûtiers, plusieurs courses et de petits détachements
étaient le service que nous faisions dans cette place. J’ignore le quantième du mois que la retraite
fut ordonnée, elle se fit fort tranquillement l’artillerie, l’équipage et les vivres se sauvèrent.
En avant de Séez petit village au pied du Petit-Saint-Bernard, l’armée s’arrêta pour laisser prendre
le devant aux canons qui se traînaient avec peine sur la montagne. Au milieu de la nuit elle reçut
ordre d’allumer un grand

\marginpar{\footnotesize\raggedright 1793 et 94.} nombre de feux et de se mettre en marche pour
passer le Saint-Bernard. Je dis au soldat qui me servait d’amener mon cheval, il vint me dire qu’il
s’était détaché du piquet et que personne ne savait ce qu’il était devenu. En attendant toute la
troupe partit et je restai tout seul avec mon domestique sur cette plaine pour chercher mon cheval ;
la nuit était fort obscure, les feux que nous avions allumés étaient sur le point de s’éteindre
avant que je décou–vris la moindre trace de ma bête. Ce n’est qu’après bien des peines et des
inquiétudes, que tout à coup j’entendis arriver mon cheval au petit trot qui paraissait me chercher
à son tour ; je fus au comble de la joie, et je rejoignis le corps aussi vite que possible. À peine
arrivé à Saint-Germain village situé à la pente du mont, l’ennemi nous attaqua à l’aube du jour, on
se battit toute la journée sans se faire grand mal ; la nuit permit à notre armée de se retirer à
moitié montagne sans être inquiétée que par quelques grenades royales qui passaient sur nos têtes ;
on bivouaqua et le lendemain l’armée entra dans les retranchements du Petit-Saint-Bernard. L’ennemi
était resté en Savoie, notre régiment marcha à Morgex pour prendre ses quartiers d’hiver, nous
étions à la fin d’octobre. À peine les arrangements nécessaires furent-ils pris que nous reçûmes
l’ordre d’aller camper sur le Saint-Bernard, nous y restâmes jusqu’au milieu de novembre et de là
nous nous rendîmes à Aoste et quelques jours après à Ivrée. Chacun se trouvait heureux d’être rendu
à un état de tranquillité ; on oublia bientôt les désagrements de la campagne ; Un joli séjour, très
bonne société, des bals et le théâtre nous dédommageaient de tout ce que nous avions souffert. Ce
n’est qu’au commencement de 1794 que je fis la connaissance d’une jeune et jolie femme qui ne
contribua pas peu à me rendre le séjour d’Ivrée infi–niment agréable. Je devins éperduement
amoureux, et malgré un nombre de rivaux je résolus de toucher le coeur de ma belle. Ma passion
s’annonça sous de malheureux auspices ; Albert Zehender mon bon ami me demanda en duel à cause des
soins que je rendais à Mme Bonis, que lui aimait aussi. Je me servis à regret de mon épée, Zehender
fut grièvement blessé ; ce fâcheux accident se sut par la ville, et les autres rivaux s’éloignèrent
sans coup férir. Me voilà seul dans le coeur de ma charmante amie, après deux mois de sollicitudes
et de peines ; j’ai joui pendant mon règne de toutes les faveurs que l’amour connaît, excepté celle
qui seule rend heureux, elle devait être la

\marginpar{\footnotesize\raggedright 1794} récompense d’une fidélité reconnue. Les femmes italiennes
en général ne laissent pas longtemps soupirer les amants, et si j’avais été plus au fait alors des
us du pays, je n’aurais pas eu à ma plaindre des rigueurs de mon amie. L’hiver de 94 est celui qui
m’a offert le plus de jouissances et de satisfaction pendant le cours de mon service en Piémont ;
des plaisirs de tout genre absorbaient mon temps. À l’approche du printemps M. le Général de
Rochmondet me promut au grade de Sous-lieutenant des Chasseurs, je m’achetai un joli sabre qui
excita la jalousie de mon Lieutenant M. F. qui apparemment n’avait pas les moyens de s’en procurer
un ; il me dit un jour à la promenade, en présence d’une quantité d’officiers : Je crois que vous
portez là une arme dont vous auriez de la peine à vous servir ; je lui répondis qu’en tout cas mon
épée se trouvait à la maison ; il me comprit et chercha dès ce moment à me voir seul ; l’occasion
s’offrit d’abord après ; il m’aborda me disant : M. il y a longtemps que vos propos me déplaisent,
allez chercher votre épée et suivez moi ; je l’assurai qu’il venait de prévenir mes intentions. Nous
nous battîmes un bon moment, je reçus une bien légère blessure au bout du bras ; mon ennemi me dit
qu’il était content, et nous nous raccommodâmes tant bien que mal. Nous partîmes d’Ivrée au mois
d’avril pour la vallée d’Aoste, je quittai ma chère amie avec beaucoup de regrets et promis de lui
écrire ; nous campâmes pendant quelques jours à la villotte, de là on nous cantonna à La Salle ; le
24 fut la malheureuse affaire où nos troupes perdirent le Saint-Bernard par la négligence, et non la
trahison comme le prétendent des calomniateurs, du Capitaine Bégoz, qui fut fait prisonnier de
guerre à Montvalezan. L’armée piémontaise se retira à pierre taillée et sur ses hauteurs ; je ne me
rappelerai jamais sans pitié, de la nuit que nous passâmes sur la route des retranchements du Prince
Thomas, au village de Morgex ; l’ignorance de nos généraux et la découragement du soldat devaient
nous livrer à l’ennemi, heureusement que le sort en décida autrement ; les Français ne nous
poursuivirent point avec la rapidité que leurs succès auraient dû leur inspirer. Au moment de nos
désastres, je fus envoyé à Aoste pour prendre le prêt, en arrivant dans cette ville, je rencontre le
Duc de Montferrat fils du roi qui venait à l’armée, il me fit arrêter et me pria de lui faire le
détail

\marginpar{\footnotesize\raggedright 1794.} des revers que nous venions d’essuyer. Dans le courant
de mai et de juin nous reprîmes les positions que les Français nous avaient enlevées, ceux-ci se
tinrent tout l’été au Saint- Bernard. On se battit sérieusement le 12 juin et le 24 juillet sans
qu’aucun parti n’ait fait grand mal à l’autre ; presque tous les jours il y avait des affaires
d’avant-postes, qui furent rarement au désa–vantage des Piémontais. Je fis toute cette campagne
comme officier aux Grenadiers, notre camp était fort agréable ; le Duc de Montferrat qui commandait
l’armée, sous les auspices du Général Bachman, comblait de bien le soldat, outre sa ration ordinaire
il recevait une augmentation de pain de riz et de lard. Les officiers prirent les meilleurs
arrangements pour se rendre la vie du camp moins pénible ; à l’approche du froid, on bâtit une
grande baraque où une vaste cheminée offrait un chaud asile ; chacun se fit construire une petite
maisonnette qui le mettait à l’abri des injures de l’air. J’eus dans ce temps une dispute avec un
officier du régiment au sujet d’une fort petite dette de jeu ; après la lui avoir demandée plusieurs
fois en vain, je lui dis un jour en présence de beaucoup d’officiers, que je lui en faisais cadeau ;
ce propos lui déplut, et il m’en demanda satisfaction. Dufresne était grand, bel homme, il se
présenta sur le champ de bataille avec une suffisance et une persuasion de vaincre qui me
choquèrent, je fus assez heureux de le blesser au bras, après une convention réciproque de ne pas
nous tirer sur le corps. Il me dit avec étonnement que c’était la première fois qu’il avait été
touché avec une épée ; je lui tendis la main, et depuis ce temps j’ai très bien vécu avec lui. Le 24
octobre je résolus d’aller à Ivrée voir ma maîtresse (si on peut appeler ainsi une femme qu’on n’a
jamais touchée...), et de là aller passer quelques jours à Turin ; je manquai à mon insu aux formes
en demandant une permission, m’adressant au Duc de Montferrat, avant d’en parler au Colonel du
bataillon des Grenadiers, dont je faisais partie, mon ignorance ou plutôt mon étourderie me coûta
cher. Le marquis Carret Commandant des Grenadiers prévint le Duc que je m’étais adressé à Son
Altesse Royale avant d’avoir le consens–tement de mon chef, sur quoi le prince donna l’ordre de
m’arrêter à Turin, Ivrée ou au Fort de Bard et de me retenir aux arrêts.

\marginpar{\footnotesize\raggedright 1794 et 95.} En approchant de Bard je fis remarquer le Fort à
mon compagnon de voyage m’écriant : je ne voudrais pas y être en peinture, une 1/2 heure après j’y
étais embastillé. Ignorant à Turin le malheureux sort qui m’attendait, je me livrais à tous les
délices qu’offrait alors le séjour de cette charmante capitale ; je partis le 5 novembre avec
beaucoup de regret ; à Ivrée je vis encore ma maîtresse. En arrivant au Fort de Bard, le gouverneur
me fit arrêter, il m’assigna une petite chambre au haut du château, fort où j’ai passé cinq semaines
dans une paix profonde. Les reproches de ma conscience ne troublaient point mon repos, j’étais
parfaitement tranquille ; les officiers d’invalides qui composaient la garnison venaient
soigneusement me tenir compagnie, M. le Gouverneur me fit l’honneur de me voir à différentes
reprises et m’invita à dîner. Je regrettais cependant de passer en prison, un temps que j’avais voué
à un semestre, j’écrivis au Duc pour obtenir mon élar–gissement, et ne reçus point de réponse ; une
lettre au Général Bachman eut plus de succès, il s’intéressa en ma faveur, et l’ordre de ma
délivrance arriva enfin le 14 décembre. D’abord arrivé au bataillon de Grenadiers, je m’adressai à
mes supérieurs pour avoir un congé, il me fut accordé à condition que je serais de retour au corps
le 1er de février 95. Je partis d’Arvier avec l’ami Wyttenbach, amoureux alors de la jolie
Angélique, il me fit le plaisir de m’accompagner jusqu’à Aoste, et me prêta son cheval jusqu’à
Saint-Rémy ; j’avais vendu le mien pendant mon séjour à Turin. Le 29 décembre je me trou-vai à Berne
; mes parents me reçurent fort bien. La veille du Nouvel An il y eut un souper, dont les suites
m’ont été bien pénibles ; on cassa, comme à l’ordinaire bouteilles, assiettes, verres etc. et les
jeta par la fenêtre ; la sentinelle qui se trouvait près de l’abbaye des gentilshommes, fut
atteinte, elle appela la garde. En se retirant du souper, je voulus être le dernier qui sortit de la
maison afin d’éviter d’être arrêté ; d’autant plus que j’étais en uniforme. La garde arriva trop
tard pour se saisir des premiers, elle se plaça aux portes de l’abbaye pour prendre ceux qui
sortiraient les derniers ; au moment qu’elle m’apperçut elle fit un mouvement pour tomber sur moi ;
j’oubliais alors que j’étais dans un pays, où on ne rendait pas le respect qui en Piémont était dû à
l’uniforme ; je ne réfléchissais point qu’en Suisse je ne jouissais point de mes prérogatives
militaires ; seul occupé de l’affront que j’allais essuyer, je tire mon épée menaçant de percer le
premier qui m’approcherait ; les soldats n’en conçurent aucune inquiétude ils tombeèrent sur moi,
m’arrachèrent mon épée, la mirent en pièces et me

\marginpar{\footnotesize\raggedright 1795.} régalèrent de coups de crosses. Tout ce bruit attira mes
amis, ils vinrent à mon secours, et mirent la garde en pleine fuite, après en avoir rossé quelques
individus. Chacun se retira, moi seul je m’occupai des suites fâcheuses que cet événement pourrait
avoir ; je priai l’ami Sinner qui était avec moi d’aller voir le chef de poste de la garde qui avait
été insultée, nous y allâmes ensemble, et le conjurâmes de n’en faire aucun rapport le lendemain ;
mon nom, et un écu de 6 francs que je lui mis entre les mains le rendirent docile à mes
sollicita–tions, et j’allai me coucher fort tranquillement. L’année 1795 eut un com–mencement plus
orageux que je me l’imaginai d’abord ; toute la ville retentit du bruit qui avait eu lieu pendant la
nuit, et on rendit la chose beaucoup plus sérieuse qu’elle ne l’était en elle-même ; la cabale jugea
à propos d’y mêler son fiel venimeux ; il serait trop long de détailler ici tout ce qui fut mis en
jeu pour avoir prise sur moi ; on décida que l’affaire serait jugée en conseil de guerre ; M.
l’Avoyer Steiger, qui en était alors le président, me fit l’honneur de me dire, que je n’avais qu’à
m’amuser pendant mon semestre, me faisant entendre qu’il ferait en sorte que la chose ne serait
traitée qu’après mon départ. Je n’eus rien de plus empressé que de suivre un aussi généreux conseil,
le mois de janvier fut consacré à tous les plaisirs dont il était susceptible, et même à
quelques-uns qui auraient pu être dangereux. Je rejoignis le régiment au jour désigné, et malgré le
désagrement qui avait troublé mon bonheur, j’ai quitté la patrie avec les plus sincères regrets. En
arrivant à Aoste le 30 janvier, on était en Carnaval ; des fêtes de toute espèce, et surtout des
bals que les officiers payaient chèrement me distrayaient, mais ils ne pouvaient me faire oublier
l’impression qu’avait fait sur mon coeur la connaissance de Julie de Werdt de Toffen, déjà les
derniers jours de l’autre semestre, je la vis une fois, elle me plut du premier abord, et pendant
longtemps elle fut chère à mon coeur, comme on le verra plus bas. En général mon coeur a toujours eu
un penchant irrésistible pour le beau sexe ; il aurait envisagé les liens du mariage comme un
bonheur inappréciable, si les circonstances qui nous commandent ne lui avaiant refusé constamment de
s’unir à une épouse chérie ; encore aujourd’hui en écrivant ceci mes vues ne tendent qu’à ce but ;
mais toujours dans le principe qu’il est préférable d’être garçon heureux qu’époux malheureux ; et
qu’on ne doit se marier

\marginpar{\footnotesize\raggedright 1795.} si on n’est dans la conviction qu’on rendra son épouse,
ses enfants et soi-même heureux. Le mois de février se passa à Aoste ; je reçus beaucoup de
politesses dans la maison du préfet Foissa, je fréquentais Mme la Baronne du Noyer et sa charmante
fille, la marquise de la Vierne. Les capitaines et les lieutenants des Grenadiers se trouvaient
alors absents et comme le plus ancien Sous-lieutenant je commandais les trois compagnies. L’ami
Wyttenbach ne me quittait jamais dans ce temps-là, nous mangions ensemble, malgré que nous
demeurions fort loin l’un de l’autre, il venait chez moi ; je me rapelle que nous avons donné à
dîner un jour au vieux Ravet et à sa fille Angélique. Le bonheur dont nous avons réellement joui
pendant notre séjour à Aoste, fut interrompu par l’ordre du départ ; je reçus à la même époque mon
brevet de Lieutenant. Nous nous mîmes en marche pour Ivrée, où nous rencontrâmes le régiment Suisse
de Bachmann qui nous donna un excellent dîner ; notre destination était au Béatange, couvent aux
environs de Coni, le voyage se fit très agréablement ; nos Grenadiers ne restèrent que quelques
jours dans ce couvent, ils partirent pour la Crave ; je fus relevé par Gruber qui vint me remplacer
aux Grenadiers et je partis pour le corps (le 7 avril), qui était alors à Alexandrie. À peine arrivé
dans cette ville, j’appris que mon bataillon était parti, je le joignis à Gorzegno, nous y restâmes
quelques jours, et en repartî–mes pour nous rendre à Ceva le 1er de mai. Après un petit séjour dans
cette ville on nous fit camper sur les hauteurs qui l’entourent ; le 12 de ce mois je reçus six
doubles louis de mon père, qui me furent volés dans ma tente ; un funeste coup pour un officier
subalterne, surtout en temps de guerre. Dans le courant de l’été nous changeâmes fort souvent de
positions ; de Bagnasco on nous fit marcher au col de Pin, de là à la Sotte que nous quittâmes pour
aller occcuper la Spinarde qui fut emportée à l’ennemi sur la fin de juillet. Outre la perte de mon
argent j’eus encore un autre chagrin au camp de Ceva ; une grossièreté de la part de Wyttenbach me
força de me battre avec un de mes meilleurs amis, le hasard me favorisa de nouveau, je vainquis mon
adversaire ; je reçus aussi une égratignure au pouce droit, la poignée de mon épée, quoique en
argent fut percée — ce duel n’a pas eu d’autres suites, nous étions les meilleurs amis après.

\marginpar{\footnotesize\raggedright 1795.} Pendant cinq mois nous occupions la redoute espagnole ;
les deux bataillons du régiment réunirent dans ce camp tous les agréments, et toutes les commodités,
dont il peut être susceptible. Ma tente était surtout parfaitemant arrangée, les après-midi, le
Colonel Stettler et d’autres officiers y venaient faire leur partie aux tarot. Le 7 août un second
vol troubla de nouveau la tran–quillité et le bien-être dont je jouissais pendant mon séjour à la
redoute espagnole. On m’emporta mon portefeuille, qui contenait alors une somme de 15 louis en
billets royaux. Les premières années que j’ai passées au service, j’avais plus le goût de la dépense
que de l’économie ; en avançant en âge les représentations amicales de mes parents m’apprirent à
connaître l’utilité de cette dernière, et je commençais à avoir beaucoup d’ordre dans mes finances
lorsque tout le fruit m’en fut ravi, par les vols des mois de mai et d’août. Mes chers parents
reconnurent mes bonnes dispositions, ils furent assez généreux de rembourser en grande partie la
somme que je venais de perdre et m’encouragè́rent par leur confiance à éviter de plus en plus de
faire des dépenses inutiles. Au milieu de septembre je fis un trè́s joli voyage avec M. le Colonel
Tschiffeli, le jeune Jayet me prêta son cheval ; nous partîmes du camp pour nous rendre à
Bardinetto, village gênois, et quartier général de M. d’Argenteau, nous soupâmes avec ce Général, et
le lendemain nous prîmes la route de Loano. Je jouissais, pour la première fois, de voyager dans des
bois d’oliviers, mais je fus ravi du coup d’oeil de la mer ; ces rivages étaient alors couverts de
vaissaux anglais, et ses bords fourmillaient de troupes autrichiennes. nous côtoyâmes la mer pour
nous rendre à Finale ; une allée de citronniers, d’orangers, et de figuiers, produite sur un terrain
sec et aride, nous y conduisit ; ce qui me frappa le plus dans cette jolie ville de la Riviera de
Gênes, était un camp établi sur la place ; il y avait alors le régiment de Terzi ; les soldats se
plaignaient de la chaleur, je mis la tête dans une de leur tente, qui me parut en effet un four, et
je me réjouissais d’être campé sur la montagne, où nous avions un climat délicieux. Nous quittâmes
Finale pour nous rendre à Savone, un petit bâtiment nous y mena dans quelques heures. La journée fut
em–ployée pour voir cette ville qui a un joli port de mer, et un château assez

\marginpar{\footnotesize\raggedright 1795.} fort, aussi longtemps que l’ennemi ne parvient pas à
s’emparer des hauteurs qui les dominent du côté de terre. Les maisons de Savone sont bâties en
briques comme à peu près toutes celles de la Riviera, les rues sont pavées de même et fort étroites.
Le voyageur doit se plaindre de la malpropreté des auberges dans cette partie de l’Italie, si
quelque chose peut l’en dédommager c’est l’excellent poisson accommodé à l’huile vierge qu’on y
mange. Nos finances ne nous permirent pas de pousser jusqu’à Gênes, nous nous rembarquâmes et prîmes
la rade de Pado, le pavillon autrichien qui flottait sur ses murs se montrait de loin ; nous eûmes
beaucoup de peine pour débarquer ; le Général autrichien Lipthay avait donné les consignes les plus
sévères ; il nous parut cependant qu’elles n’auraient pas dû concerner des officiers du roi de
Sardaigne allié alors à l’Autriche. Ce Général nous reçut grossièrement pour un homme dont on attend
des procédés honnêtes, il nous donna cependant un passeport pour aller à bord de la flotte anglaise
qui croisait dans ces parages. Je montais l’Agammemnon de 64 canons, commandé par Nelson ; le
capitaine anglais n’était pas plus honnête que le Général autrichien, il m’adressa à peine la parole
et me donna un mousse pour me promener sur le vaisseau. Après avoir tout vu et tout admiré, je
rejoignis le Colonel Tschiffeli qui n’avait pas quitté la chaloupe ; nous arrivâmes heureusement à
Finale, où nous nous baig–nâmes, et le jour aprè́s nous fûmes de retour à notre camp, très contents
et satisfaits de notre course. La nuit du 4 au 5 octobre est une de celles dont le souvenir ne
s’effacera jamais. Tranquillement couché dans ma tente, je faisais des remarques sur la singularité
de ce que par le plus terrible orage une toile me mettait à l’abri du vent et de la pluie, lorsque
tout à coup ma bonne maison fut emportée avec duvet et couverture de mon lit, et qu’une pluie
affreuse m’inondait dans un instant ; je me sauvai en chemise dans la tente de mon voisin le
Capitaine de Watteville, qui me crie : n’entrez pas, sinon le vent m’emporte, je ne l’écoutai point
et lui dis que j’étais nu, il me laissa entrer, me couvrit d’une capote et à peine ma toilette
achevée que sa tente quoique une marquise, eut le sort de la mienne ; nous nous sauvâmes tous les
deux dans celle de Wyttenbach, qui tint quelques heures plus longtemps, elle succomba au milieu de
la nuit, et avec elle le reste du camp. Par bonheur que le lendemain le temps fut des plus beaux,
chacun se sécha et on affermit mieux ses tentes. Il est difficile de se faire une idée des orages
qui ont lieu sur les montagnes de l’Apennin, le vent avec un bruit de tonnerre déracine les plus
gros arbres, il enlève des mulets chargés et les porte à de longues distances de la route avec les
conducteurs. Au mont alegro, poste occupé par nos Chasseurs, les soldats se creusaient de larges
fossés et s’enterraient, ne pouvant se tenir

\marginpar{\footnotesize\raggedright 1795 et 96.} sur la surface de la terre ; les jours où il
faisait beau, le climat était d’une douceur extrême le ciel serein nous découvrait les montagnes de
l’île de Corse. À juger de la rapidité avec laquelle se suivent les récits de mes duels on pourrait
me croire querelleur, ce que je n’étais certainement pas ; il faut observer, que ce qui paraît
rapproché dans ce recueil l’est bien moins dans le fait, je ne parle que des événements les plus
intéressants de ma vie, plusieurs mois se sont souvent passés sans que rien de remar–quable me soit
arrivé. C’était à peu près à la fin d’octobre qu’un badinage mal interprété par François Z.
l’engagea à me demander satisfaction ; l’issue de la chose prouva que le tort était de son côté, il
reçut une forte blessure, et nous nous raccommodâmes fort bien. Au mois de novembre le séjour de la
montagne devint fort désagréable ; le 23 les Français attaquèrent sur tous les points, l’aile droite
et l’aile gauche, gardées par les Autrichiens se replièrent, les Piémontais qui défendirent le
centre conservèrent leur position. Par une bravade mal placée, le Général Colli, qui nous
commandait, ne voulait point céder le terrain quoique par l’insuccès des Autrichiens nous ne
pou–vions plus tenir avec avantage notre première ligne. Le 28 il se décida de commander le retraite
sur Ceva, l’ennemi qui a eu le temps de se renforcer nous inquiéta beau–coup et ce qui aurait pu
s’effectuer cinq jours plus tôt avec honneur et gloire, se fit après à la débandade, en désordre par
un temps affreux et en laissant à l’ennemi onze canons qu’on eut à peine le temps d’enclouer. Notre
régiment faillit geler dans cette fameuse expédition, les pieds de plusieurs soldats ont eu ce sort
pendant le temps que nous étions à la Sotte et à la Spinarde ; heureusement que nous reçûmes l’ordre
de marcher au Jouet où commandait le Prince de Carignan ; nous passâmes la nuit au bivouac et à
l’aube du jour nous prîmes le chemin de Ceva. Les Français ne nous poursuivirent point, ils
occupèrent les hauteurs que nous venions d’abandonner ; notre armée prit position sur les collines
aux environs de Ceva, les deux bataillons de Stettler campèrent à la Tête noire. Le Brigadier
Stettler reçut le commandement de l’aile gauche sous les ordres du Général Montafia, M. Stettler me
choisit pour son aide de camp, notre quartier était au Mondon. Ma situation était alors plus
agréable que celle de mes cama–rades, qui au mois de décembre se trouvaient encore sous la tente.
Sur la fin du mois, l’armée prit ses cantonnements nous fûmes destinés pour la Toricella petit
village à 2 petites lieues de Ceva ; le Brigadier partit pour la Suisse il me remit le détail et le
commandement de sa compagnie. Notre séjour à la Toricella fut très court nous y passâmes cependant
le Nouvel An 1796 où malgré le peu de ressour–ces qu’offrait ce pays nous nous amusions assez bien.
Le 2 janvier l’ordre d’entrer en quartier d’hiver nous surprit fort agréablement, nous quittâmes
avec plaisir ce village pour aller nous établir à Ceva, lieu de notre destination.

\marginpar{\footnotesize\raggedright 1796.} Ce n’est qu’après les fatigues et les désagréments de la
guerre qu’on est dans le cas d’apprécier le repos et la tranquillité. Le séjour de Ceva fort peu
attrayant par lui-même, nous parut des plus agréables ; le 3ème bataillon du régiment qui avait
passé tout l’été à Alexandrie, vint nous joindre et nous fûmes réunis pendant trois mois, hormis les
Grenadiers qui prirent leur quartier d’hiver à Caru. Je me suis beaucoup amusé à Ceva, j’avais mon
logement en commun avec le Major Jonquière, dont l’humeur joyeuse ne contribua pas peu à me rendre
mon temps agréable ; je fréquentais souvent la maison du Marquis de Ceva et comptais des douceurs à
la petite Fanchette qui était trè́s jolie mais fort sotte ; mon ami Wyttenbach en était éperduement
amoureux. Je voyais quelquefois Mlle Davic qui peut passer pour une des plus belles filles du
Piémont, Mme Bringin ne lui cédait rien en beauté. La garnison donnait assez fréquemment des bals ;
j’étais lié avec plusieurs officiers du régiment de Savoie qui était avec nous ; je voyais par
contre peu ceux de la reine. Le jeu a troublé en partie le bonheur de ce temps-là, j’ai perdu 300
livres tournois en billets et pour réparer ce déficit de ma bourse, je me me refusais beaucoup de
petites jouissances que je m’étais accordées jusqu’ici, et je parvins à me rendre cette perte moins
sensible. Janvier, février et mars s’écoulèrent dans ce mélange continuel de douceurs et de peines.
Le retour de la belle saison nous arracha au commencement d’avril du paisible séjour de Ceva.
L’arrivée du Général Buonaparte à l’armée d’Italie, dépourvue de tout et mourant de faim ; le refus
du roi de Sardaigne de faire la paix avec la France, faisait présumer que les hostilités ne
tarderaient par à commencer. Chacun se promettait les plus grands avantages de cette campagne, et ne
supposait pas de possibilité de pouvoir être vaincu, mais la suite a donné de cruelles preuves du
contraire. Buonaparte attaqua nos troupes avec un acharnement inconcevable ; les avant-postes
autrichiens l’avaient provoqué par des escarmouches continuelles, et son attaque tendait plus à
s’assurer sa retraite qu’à la conquête de toute l’Italie ; l’ordre d’embarquer la grosse artillerie
et de se retirer derrière le Var était déjà donné, l’armée manquant totalement de vivres. Du 10 au
20 avril on se battit journellement et les Français emportaient l’avantage ; M. Stettler reçut dans
ce temps le commandement d’une aile et me prit de nouveau pour son aide de camp. Le 21 avril il y
eut une attaque générale au Mondovi, l’armée du roi

\marginpar{\footnotesize\raggedright 1796.} fit des prodiges de valeur et fut néanmoins forcée à la
retraite ; notre Général Colli perdit la tête, et sans une charge de cavalerie qu’exécuta le
régiment des Dragons du roi, et où il mit celle de l’ennemi en fuite, l’armée piémontaise aurait
peut-être été toute prisonnière de guerre. Ce malheur ne tomba en partage qu’à six compagnies des
gardes, et deux bataillons de notre régiment, que Colli oublia dans la redoute du Mondovi d’où il
avait déjà fait amener les canons et les munitions de guerre. En vain nous nous attendions d’un
moment à l’autre à recevoir l’ordre de nous retirer avec le reste de l’armée qui s’en allait à
grands pas ; entièrement abandonnés à nous-mêmes, n’ayant pour toute défense que le peu de
cartouches qui restaient dans les gibernes du soldat, nous prévoyions le sort qui nous attendait. Le
Général Buonaparte fit sommer la ville de Mondovi de se rendre dans un quart d’heure, sinon la
garnison passerait au fil de l’épée ; le Chevalier Delléra qui commandait la place, le Brigadier
Stettler, et le Colonel du régiment des gardes qui avaient leurs régiments sous leurs ordres, furent
indignés des propositions de Buonaparte ; on renvoya le parlementaire avec une réponse noble et
ferme et des demandes honorables ; la lettre fut écrite par M. Stettler, elle lui fait honneur.
Pendant que nous étions ainsi en pourparlers, l’ennemi jeta plusieurs bombes dans la ville, et
quelques boulets dans la redoute que nous occupions, on étendit le drapeau blanc sur les remparts et
le feu ne cessa point. M. Stettler offensé de cette atteinte aux droits de la guerre, proposa au
gouverneur de la ville de refuser toute capitulation et de se faire jour à travers l’ennemi ; il
sortit du congrès sans attendre la réponse du Chevalier Delléra et vint à la redoute. Écumant de
rage il tint aux soldats un discours énergique et digne d’un Suisse, il les exhorta de se préparer à
repousser la force par la force, et à se rappeler la valeur de leurs ancêtres. Le régiment fut saisi
d’enthousiasme et chacun résolut de mourir pour le roi ; au moment d’un si beau zèle arrive un
officier des gardes dans notre redoute qui nous crie : d’ordre du roi, braves Suisses, posez les
armes. Il est difficile de se former une idée de notre consternation en recevant un pareil ordre ;
l’officier dit à M. Stettler que le gouverneur ne pouvait approuver l’avis qu’il lui avait donné de
refuser toute capitulation, que comme Suisse il aurait été du même sentiment

\marginpar{\footnotesize\raggedright 1796} mais le roi lui ayant confié une province, il se trouvait
obligé de ménager les sujets et de préserver la ville des malheurs inévitables qui l’attendraient si
elle fit une longue résistance, qui par le manque absolu de munitions de guerre était même
impossible. Il fallut obéir, et attendre la réponse du Général français sur la lettre du Brigadier
Stettler ; un autre parlementaire fut en-voyé à Buonaparte pour lui témoigner la surprise de l’armée
piémontaise de se voir canonné pendant qu’il s’agissait d’une capitulation ; un billet conçu dans
ces termes fut toute notre satisfaction : „J’avais donné ordres à„ „mon chef d’artillerie de dresser
une batterie de canons en face du Mondovi„ „pour la bombarder en cas qu’elle voulut faire
résistance, il exécuta mes ordres„ „avant que j’eus le temps de la faire prévenir que je venais de
proposer une „capitulation, je ferai cesser le feu sur-le-champ.“ Signé : Buonaparte„ Sa réponse à
nos propositions était telle qu’on pouvait l’attendre du vainqueur, pour le-quel rien n’est plus
doux que d’être généreux. La ville, les personnes, et propriétés respectées et la garnison
prisonnière avec les honneurs de la guerre. Buonaparte entra en triomphe en ville ; l’État général
piémontais alla à sa rencontre, en s’approchant de lui, il salua de son sabre, et dit au Chevalier
Delléra :„Je suis Bonaparte „c’est aujourd’hui la 1ère fois que j’ai eu l’honneur de me battre avec
les troupes „de votre maitre, leur bravoure augmente ma gloire.„ Le Général permit en même temps que
les officiers prisonniers conserveraient leurs épées. C’est à la demande du Général piémontais de
laisser l’épée aux officiers, que Bonaparte répondit : Je suis charmé de vous donner cette preuve de
la générosité française. Nous passâmes la nuit dans un tas de foin, je ne quittai pas M. le
Brigadier Stettler. Quel réveil que celui du 22 avril ! en sortant du foin on nous fit partir pour
Ceva ; à Lesegno, où Buonaparte avait établi son quartier général, notre chef voulut lui demander
l’agrément d’aller en Piémont sur parole, mais le Général ne le reçut pas, et nous fûmes obligés de
prendre le chemin de Nice. Comme je voyageais avec M. Stettler je fus bien traité partout, on eut
égard à son rang ; notre route allait par Bagnasco, Garessio, Ormeea, Piova Oneglia et San Remo où
nous nous em-barquâmes pour Nice. M. le Brigadier son neveu Théophile et moi étions les seuls de 23
qui ne souffrions pas de la mer, nos camarades avaient des malaises continuels. On nous laissa trois
jours en repos après notre arrivée à Nice ; nous eûmes le désagrément d’y voir célébrer la victoire
du Mondovi et les précédentes ; 22 drapeaux, tant autrichiens que piémontais et suisses, furent
étendus sur un balcon de la maison de ville, et on dansa la Carmagnole autour. J’eus le bonheur de
trouver de l’argent que

\marginpar{\footnotesize\raggedright 1796.} je pris chez M. Peillon sur la maison Porta à Lausanne,
j’en remis la moitié à M. le Brigadier. Manquant de tout, je m’achetai quelques habillements et me
pourvus du nécessaire pour continuer la route ; à Nice je vis la comédie, où on représentait des
pièces révolutionnaires ; ce qui me frappa le plus dans cette ville fut une bataille entre deux
capitaines français, ils se prirent par les cheveux, se roulèrent par terre saignèrent tous les deux
en plein café, sur le cours, en présence d’un grand nombre de spectateurs ; ma surprise diminua en
reconnaissant dans un des guerriers mon ancien perruquier d’Annecy. Nous partîmes de Nice le premier
de mai, pour nous rendre à Valence en Dauphiné ; j’ai souffert pendant le voyage d’un vent horrible
qui désole toutes les années cette partie de la France, c’est le mistral ; habit, frac, manteaux
étaient insuffisants pour nous en garantir. Notre route était par Cannes, Fréjus, Brignoles, Aix,
Orange, Avignon et Valence. ?À Avignon je fis la connaissance de Mme de Van Berchem née Polier,
Messieurs Jayet et de Brenles me procurèrent cette charmante occasion d’apprendre à connaître une
maison, où pendant trois jours je fus comblé de politesses. Mme de Van Berchem nous donna un joli
goûter à la Rosette, agréable campagne du gendre de Mme de Van Berchem qui dans les orages de la
révolution y fut assassiné par les jacobins. Un ami de Madame de Van Berchem certain riche polonais
qui vivait depuis longtemps à Avignon, nous donna un dîner splendide ; je fais mention de ce galant
homme autant par reconnaissance pour ses honnêtetés que pour me rappeler d’une singularité qui lui
était particulière, il ne pouvait achever une phrase sans y mêler 2 ou trois fois le mot « exésavier
», terme qui d’ailleurs n’a aucune signification quelconque. Je rends de même hommage à la beauté et
aux talents de Mlle Nanette de Van Berchem maintenant Mme Bazin, elle aida Mme sa mère à nous rendre
le court séjour auprès d’elle aussi agréable que possible. Nous quittâmes Avignon pleins des douces
sensations que nous y avions éprouvées, le voyage n’offrit plus rien qui méritât d’être remarqué.
Nous fûmes très bien reçus par les habitants de Valence ; M. Stettler fit écrire au Directoire de
Paris pour obtenir l’agrément de nous rendre en Suisse sur parole ; en attendant la réponse, nous
prenions les arrangements les plus convenables pour passer notre temps agréablement, une dame
s’offrit très gracieusement de soulager le sort des malheureux prisonniers, elle nous invita à venir
passer une partie de la journée chez elle. Ceux des officiers

\marginpar{\footnotesize\raggedright 1796.} qui préféraient la bonne société à la vie désagréable
des cafés, se rendirent avec empressement à une invitation aussi grâcieuse ; je fus de ce nombre.
Mme Beaude était une femme de 23 ans, jolie aux yeux de quelques-uns, très intéressante aux yeux de
tous. La haine pour la révolution, qui avait à peu près ruiné sa maison, lui inspirait beaucoup
d’intérêt pour ceux qui par leur état devaient être du même sentiment qu’elle. Des officiers suisses
bernois, au service du roi de Sardaigne, ne pouvaient pas être amis d’un système qui dès son origine
menaçait le trône de leur maître, et ne paraissait point propre à consolider la liberté helvétique
en général moins encore la république aristocratique de Berne en particulier. Douée d’une âme
excessivement sensible Mme Beaude fut touchée de notre situation ; éloignés de notre patrie,
inquiets sur le sort du roi, malheureux de nous trouver dans un pays, où le nouvel ordre de choses
avait fait couler tant de sang, nous aurions pu intéresser des coeurs moins généreux que le sien.
Elle avait deux soeurs dont une de 18 l’autre de 16 ans, et c’est avec ces trois aimables personnes
que j’ai passé la plus grande partie de mon séjour à Valence ; la cadette a su particulièrement me
fixer ; que de moments délicieux se sont écoulés dans cette société charmante, combien de fois mon
imagination s’est-elle occupée depuis de la félicité dont je jouissais alors ; ah ! si toutes les
prisons étaient de cette sorte mon plus ardent voeu serait de n’être jamais libre. Ce bonheur comme
tous les autres eut aussi sa fin, le Directoire français consentit à ce que les Suisses s’en
retournassent chez eux ; je quittai Valence avec les plus sincères regrets ; nous en partîmes à la
fin de mai. Arrivés à Vienne en Dauphiné, un petit garçon qui portait une quantité d’imprimés à la
main, qu’il criait pour les vendre à un sol la pièce, nous apprit la paix avec le roi de Sardaigne.
Elle est aussi honteuse pour celui qui a été forcé de la faire, que pour ceux qui ont eu la bassesse
de la dicter. Nous nous arrêtâmes quelques jours à Lyon, où le théâtre, et les décombres de la place
de Louis XV étaient les seuls objets qui attirèrent principalement mon attention. Je profitais de ce
temps pour écrire à Mme Beaude, je lui demandai des cheveux des trois soeurs, qui me furent refusés,
par un trop grand attachement aux préjugés de ce monde ; notre correspondance n’a jamais été
interrompue depuis, et je reçois toujours avec un vrai plaisir des nouvelles de cette intéressante
famille ; le père de Mme Beaude est à l’île de Malte, son mari est en Égypte, sa soeur Lucie vient
de se marier

\marginpar{\footnotesize\raggedright 1796.} à Bourg-en-Bresse, la cadette, Mlle Marthe vit à Valence
avec sa mère. Mme Beaude a deux charmants enfants, Jean-Jacques et Adèle, le garçon a actuellement 7
et la fille 5 ans. À Lyon nous prîmes une voiture entre le Brigadier, Stettler de Gotstadt, Herbort,
de Watteville, Dietzi et moi, pour nous conduire à Berne ; le voyage se fit très agréablement ; à
notre passage à Ferney nous allâmes voir le tombeau de Voltaire, et nous arrivâmes à Berne le 1er de
juin. L’époque que je passais en Suisse pendant les années 96 et 97 est la plus heureuse de ma vie ;
je restais à Berne le mois de juin, en juillet j’accompagnai ma tante Stettler de Zofingue et Mme
Dachselhofer Kirchberger au bain de Schinznach, de là je me rendis à Schaffisheim chez le cousin
Brutel ; le voisinage de Lenzbourg me procura le plaisir de voir Messieurs de Watteville, fils du
bailli et mes bons amis ; je quittai mon cousin pour aller demeurer au château de Lenzbourg ; il
arriva du monde, j’en sortis, et acceptai un logement chez M. Stettler secrétaire ballival ; je fis
connaissance avec Madame, une bien aimable femme, elle me mit dans sa confidence et m’avoua son
amour pour l’aîné des W. avec lequel elle se trouvait alors en froid, je me donnais toutes les
peines pour raccommoder les amants, j’y réussis ; cette bonne oeuvre me valut l’amitié de Mme
Stettler, qu’elle accordait facilement. Au mois d’août je partis de l’Argovie avec regret, les
lettres charmantes que je recevais de temps en temps de Mme Stettler me dédommagèrent un peu. Je
partageais ce mois entre Kirchberg et Schüpfen, au 1er endroit j’eus le bonheur de trouver Julie qui
était venue voir sa soeur, le peu de jours que je passais avec elle ralluma mes feux ; au
commencement de septembre je la revis à Toffen, mes yeux lui dirent déjà alors ce qu’éprouvait mon
coeur, j’allai voir mon cousin de Greyerz, à Kirchdorf, de là je partis pour Thoune où je
m’embarquai pour Interlaken. M. le Bailli de Bonstetten qui avait servi au même régiment que moi ne
reçut parfaitement, Mme son épouse s’empressa de même à me rendre mon séjour très agréable, elle a
beaucoup d’esprit et est d’un commerce charmant. Quinze jours se passèrent à Interlaken dans le
bonheur, je retournai à Berne prenant encore ma route par Kirchdorf et Toffen. Je reçus dans ce
temps un ordre de M. le Brigadier Stettler, qui me dit que les capitaines du régiment me chargeaient
de prolonger mon séjour en Suisse, pour faire des recrues ; tous les prisonniers de guerre furent
contraints de joindre leur corps, et le régiment fut formé de nouveau. Je prenais les arrangements
nécessaires pour rester chez moi ; je fréquentais le Leist de la Brunngass, dont je fus reçu membre
ainsi que de la petite société ; le premier me procura la

\marginpar{\footnotesize\raggedright 1796.} connaissance de plusieurs jeunes gens qui sont
maintenant mes meilleurs amis ; je fis de même alors celle de l’Abbé Chappû, émigré français d’une
vertu austère et d’un mérite rare. Mon père s’intéressait au sort de cet homme que la Révolution
avait rendu malheureux, pour moins blesser sa délicatesse, en lui accordant des secours pécuniaires,
je priai l’Abbé de me donner des leçons en italien, quoique je connaissais cette langue aussi bien
que lui. Par la suite je m’aperçus que je n’avais pas même récompensé les services qu’il me rendît ;
sa conversation était très instructive, il m’enseigna les principes de la poésie française pour
la-quelle je me sentais du goût, nous lisions les grands poètes de sa nation, et je parvins à
apprécier et à sentir leur mérite, en me perfectionnant en même temps dans la langue française. Cet
homme a gagné l’attachement et l’estime de tous ceux qui ont appris à le connaître, il a quitté
Berne d’après un ordre du Directoire français, qui enjoignait à tous les émigrés de sortir de la
Suisse ; j’ignore où ce respectable ami traîne des jours qui lui sont devenus amers, depuis qu’il
voit sa patrie malheureuse. En novembre je projetai de me fixer dans une coterie de femmes ; je
choisis sans balancer celle qui réunissait les amies de Julie, et la plus jeune alors. Des soirées,
des parties de danse se succédaient, je me vouais tout entier à l’objet de mon amour, avec une
retenue et une décence qui contraste avec la manière dont on courtise les femmes en Italie ; mes
parents me conseillaient de ne pas m’attacher sérieusement, vu que ma situation s’opposait à une
union avec Julie, mais j’écoutais peu ces froids raisonnements ; l’espérance, cet heureux soutien
des amants ne m’aban–donna jamais ; les obstacles qui s’offraient au premier abord à
l’accomplissement de mes voeux ne me paraissaient pas invincibles ; M. l’Abbé qui était dans ma
confidence m’encourageait encore ; Julie de son côté recevait mes hommages avec plaisir, elle
favorisait avec une délicatesse extrême un penchant qui était peut-être réciproque ; elle apprit
l’italien dont je lui avais fait le plus grand éloge. À la fin de décembre je partis pour Bâle, où
se faisait alors le siège de la tête du pont de Huningue par les Autrichiens ; je fis ce voyage avec
Stettler de Bipp l’aîné et en conduite de mon Sous-officier recruteur Vita, qui pendant le voyage
nous a beaucoup amusés.

\marginpar{\footnotesize\raggedright 1797.} L’année 1797. commença dans les réjouissances et les
dissipations ; je dansais à la Redoute, à la Cage et aux parties particulières ; Mme Stettler arriva
de Lenzbourg dans ce temps, je la fréquentais beaucoup et elle contribua à me rendre mon séjour des
plus agréables. Le peu d’accord qui avait régné entre mon frère et moi depuis notre première enfance
fut enfin sub–jugué par la raison, et sans être dans une intime liaison nous vivions en paix
ensemble. Les amis de ma tendre jeunesse avec lesquels je m’étais brouillé déjà en 89 venaient
d’être remplacés par d’autres aux-quels j’étais cher ; une conduite sans reproche m’assurait
l’attachement de mes parents, frè́res et soeurs, et ma réputation au tégiment était telle que je
n’en avais que des jouissances à attendre. Tant de satisfaction à laquelle se joignait encore
l’heureux amour pour Julie, ne pouvaient que me rendre l’homme le plus fortuné ; je sentis l’étendue
de mon bonheur, j’en ai souvent rendu grâces au grand dispensateur et je prévoyais qu’un tel
bien-être ne pouvait avoir une longue durée. Les mois de janvier, février et mars formèrent une
chaîne de félicités. Une grande partie de mon temps se vouait aux intérêts du régiment ; je passai
bien des heures chez ma soeur Dachselhofer et lui lisais quelque chose d’agréable, par intervalle
j’allais de temps en temps pour quelques jours à Schüpfen. En ville mes soirées se passaient au
Leist ou à la petite société, je jouais alors au billard et assez heureusement ; n’étais-je pas
entre hommes, je me trouvais en femmes près de ma chère Julie, je faisais la partie avec elle et
l’accompagnais à la maison. Combien de fois nos yeux se rencontraient-ils ! Combien nous nous sommes
dit par cet intéressant langage ! Époque fortunée avec quel délice retrace-je maintenant ton
bienfaisant souvenir ; qui aurait prévu alors, que tu serais le précurseur de tant de douleurs et de
peines ? À la fin de mars il y eut un bal et un souper splendide à l’hôtel de Musique, on s’y amusa
à merveille ; les hommes restèrent jusqu’au matin ; quand il fut jour ils résolurent d’aller
déjeuner chez la soi-disante Anneli ; je me mis de la partie ; en arrivant à la maison la plus
proche de celle où nous voulions nous arrêter, nous heurtâmes à la porte avec beaucoup de bruit, le
maître vint ouvrir et nous combla d’injures, nous le laissâmes crier et prîmes notre chemin, les
autres camarades au nombre de 12 à peu près, qui nous suivirent, arrivèrent immédiatement après
nous, ils

\marginpar{\footnotesize\raggedright 1797.} entendirent les injures de cet homme, et je ne sais ce
qu’il s’y passa, mais lui est allé chez le greffier me dénoncer que je l’avais battu jusqu’au sang ;
cette affaire fit beaucoup de bruit en ville ; une quantité des jeunes gens qui avaient été du
déjeuner comparurent devant le juge, la chose tira en longueur, je protestai vainement de mon
innocence, on me fit payer une forte amende. C’est aussi au mois de mars que je fis faire mon
portrait par Mme le Grand, je le destinai à Julie ; mais réflexion faite j’y voyais de l’absurdité,
et je l’ai donné à ma soeur Dachselhofer qui l’a reçu avec intérêt. Au commencement d’avril le
moment s’approchait où je devais renoncer aux innombrables jouissances que je goûtais à Berne ; mon
départ se fixa au 10 courant avec mon ami de Watteville de Lenzbourg, et M. Wetty de Muhlhausan,
notre nouveau ministre de camp. Dimanche le 9 au soir je vis encore l’aimable Julie, je passai la
soirée avec elle, chez ma belle-soeur, en la reconduisant à la maison, je lui dis que le regret de
la quitter augmentait la peine que j’éprouvais de partir de Berne, je voulus lui en dire davantage,
mais sa soeur Charlotte m’empêcha assez bêtement de m’expliquer. Je partis lundi le 10 avril avec un
pres–sentiment qui paraissait me dire que je ne reverrais plus ma chère patrie dans ce parfait
bien-être dont elle jouissait alors. Nous arrivâmes (sans avoir eu d’autres désagréments que ceux
que nous causaient les ridicules de M. Wettly) à Genève. Nos passeports n’étant pas signés par
l’ambasssadeur français en Suisse, on nous refusa le passage par la Savoie, et nous n’y parvînmes
qu’à force d’une infinité de ruses que je mis en jeu. Aux bains d’Aix on nous prit pour les fils du
comte d’Artois, avec leur instituteur, et il fallut prouver le contraire ; ce qui n’était pas bien
difficile. Le passage du Mont Cenis a été excessivement mauvais, M. Wettly qui en jugeait sur le
reste du Piémont se désolait d’avoir mis le pied dans cet abominable pays, et nous ennuyait beaucoup
par sa désagréable humeur ; une fois à Suze, il se trouva moins à plaindre et à Turin il commença à
croire qu’on pouvait vivre, en Piémont. Nous nous arrêtâmes un jour dans cette belle ville, on nous
apprit que le 1er bataillon était à Saluzzo et nous

\marginpar{\footnotesize\raggedright 1797.} partîmes pour cette ville. Notre voyage a été de dix
jours, la conversation roulait en général sur les objets intéressants que Watteville et moi avions
laissés en Suisse ; le ministre de son côté ne parlait que des onze ans qu’il avait passés à la Haye
et il comparait la propreté hollandaise avec le sale Piémont. Les camarades du régiment nous
reçurent fort bien, le séjour de Saluzzo était agréable surtout dans cette belle saison ; nous n’y
restâmes que huit jours, le bataillon reçut ordre de se rendre à Turin, où nous arrivâmes le 1er de
mai par une pluie des plus abondantes. Je me trouvais heureux d’être en garnison dans cette
charmante capitale, le 3ème bataillon vint nous joindre ; le mois de mai n’était pas agréable il n’a
pas discontinué de pleuvoir, le temps se passait à courir au théâtre, à rouler dans les cafés et à
monter -les gardes ; je fréquentais rarement les filles publiques, dont le nombre était énorme, et
que de jeunes garçons venaient offrir à tout moment aux militaires. Je fis connaissance avec la
comtesse Borgharet, la comtesse Saint-Alban, la marquise de Verdun, j’allais le plus souvent chez la
Saint-Alban, j’y trouvais toujours l’ambassadeur espagnol Don Uhloa. Dans les mois de juin, juillet,
août, et septembre mon genre de vie était à peu prés le même ; les troubles intérieurs du Piémont
donnaient les plus grandes inquiétudes sur le sort du roi ; le service était très fatigant les
officiers furent obligés de se caserner à un couvent près du quartier des soldats afin d’être
toujours à portée de leurs troupes. Des révoltes sérieuses s’éle–vèrent à Moncalieri, une république
se proclama à Asti ; l’armée vint à bout d’étouffer les unes et de détruire l’autre. L’ambassadeur
français Miot donnait toujours les plus belles assurances de la sincérité des sentiments de loyauté
de la grande nation ; sa femme d’une vanité antirépublicaine disputa le rang aux princesses et
oubliait qu’elle n’était que citoyenne ; je fus lui rendre visite à son arrivée à Turin, elle était
fort honnête et se plaisait dans le grand monde, aussi son époux fut-il rappelé par le Directoire,
qui lui reprochait de fréquenter l’aristocratie et la noblesse. Le train de vie des militaires qui
sont en garnison à Turin est très monotone ; le matin à 6 heures, en été, on se rend à l’exercice, à
8 on le quitte pour aller se préparer pour la parade, qui a lieu sur la grande place Saint-Charles ;
après on va au rapport chez le colonel et à midi tout le monde se rend à la cour. Les antichambres
se remplissent des premiers gentilshommes du royaume, des militaires supérieurs de tous grade, de
pages, et de valets de pied et au dernier appartement, des gardes du corps ; quand le roi et toute
la famille arrivent, on fait une profonde révérence, on les suit à la messe, où on attend

\marginpar{\footnotesize\raggedright 1797.} leur retour pour faire une seconde révérence après
laquelle un chacun se retire. À une heure après midi les officiers vont dîner, ils prolongent cette
occupation autant que possible, et la quittent pour aller prendre le café et la liqueur ; on se
réunit dans un café, où on reste jusqu’à l’heure de la promenade, quelquefois dans l’intervalle on
va faire des visites aux femmes de sa connaissance, après s’être promené en long et en large sur la
citadelle, le tambour appelle à la visite du soir, on se rend au quartier et de là au théâtre, qui
finit sur les 11 heures, que chacun est charmé de se reposer pour recommencer le même train le jour
après. Je voyais beaucoup à Turin la comtesse Langouse, une aimable Niçoise fort liée avec Mme
Barthés, une jolie française, épouse du Lieutenant-Colonel au régiment de Bachmann qui était alors
en garnison avec nous. Pendant toute le temps que j’ai servi, j’ai été une seule fois aux arrêts
d’ordre du régiment, M. le Major Bucher me les donna pour 24 heures, j’avais manqué d’une heure
l’exercice du matin, par la négligence de mon domestique Pirolet, âgé de 18 ans qui a été quitte
pour un bon soufflet que je lui appliquai en arrivant à la maison. Le 13 août j’eus une dispute
assez vive à table, avec mon ami de Watteville ; nous sortîmes de la ville pour aller nous battre,
souvent l’envie me prenait de lui faire des propositions pacifiques, et l’idée, qu’il m’en fera de
son côté, s’il en aurait le désir, me retint toujours, l’assaut commença, il fut blessé. Je lui dis
après, combien de fois j’avais été sur le point d’éviter de combatte avec un bon et ancien ami, il
m’assura qu’il avait eu le même sentiment et que la même idée l’avait retenu. De pareilles
dispositions nous firent de part et autre oublier le passé. Après avoir couru tous les traiteurs de
Turin nous résolûmes de faire ménage quelques amis ensemble, on me pria de me charger du détail. Un
bon cuisinier et des provisions de tout genre furent mes premières acquisitions ; le 1er dîner eut
lieu le 1er septembre pendant quelques jours la machine parut vouloir bien aller, peu à peu nous
nous aperçûmes de l’infidélité du cusinier et de la mésintelligence entre les domestiques, et le
tout fut rompu à la fin du mois. Quelques jours auparavant des propos malhonêtes du Lieutenant Ernst
me forcèrent de lui en demander raison, nous nous battîmes, je lui portai, après un long combat, un
coup sur le plat du haut de la main qui le désarma entièrement, il ne put plus porter l’épée et ne
perdit cependant point de sang ; la chose en resta là. Le 27 septembre je reçus l’ordre de partir
pour le détachement de Chivasso, où je ma trouvais commandé par le Lieutenant Barthés, officier de
Bachmann, un jeune homme de mon âge, et mon aîné de service.

\marginpar{\footnotesize\raggedright 1797.} Pendant mon séjour à Turin j’ai toujours reçu
régulièrement des nouvelles de mes chers parents, leurs lettres me parlaient souvent de Julie et me
représentaient les difficultés qui s’opposaient à une union ; mon ami de Muralt, que j’avais prié de
me rappeler fréquemment au souvenir de ma chère amie, m’écrivait que j’étais aimé, et que Julie
serait à jamais fidèle. Les jolies lettres de Mme Stettler Fischer et celles de Mme Beaude que je
recevais dans ce temps me firent aussi le plus grand plaisir. Les mois d’octobre et novembre se
passèrent à Chivasso, où je voyais ces Dlles Capra. J’étais logé chez Mme Canis vieille et laide
femme, qui avait des attentions pour moi. Je mangeais à l’Agneau chez un bien brave hôte qui me
traitait fort bien ; Barthés vint aussi y manger. Après quelque séjour dans cette fort petite ville,
je fis la connaissance de Mme Lançon, j’oubliais pour un moment ma Julie et devins amoureux de cette
belle femme, qui est la première dont j’ai été trompé ; je croyais voir en elle un exemple de vertu,
des sentiments élevés et un caractère charmant, et j’ai trouvé à la suite qu’elle était un
assemblage du contraire. Après lui avoir fait une cour très assidue, elle se renia un jour que je me
fis annoncer et se mit à la fenêtre, pour me voir partir, je pris la peine de lui écrire une lettre
fort touchante et de me plaindre de ses procédés, elle me répondit et nous entrâmes dans une étrange
correspondance, que je finis en rompant avec elle ; je regrette de n’avoir conservé ses lettres elle
écrivait supérieurement bien le français. À son honneur, Barthès et moi avions donné un bal le 21
octobre jour de la foire, avant lequel elle nous servit un très joli souper. À la fin de novembre
les deux bataillons qui étaient en gar–nison à Turin, reçurent ordre de partir pour Alexandrie ; ils
s’embarquè–rent sur le Pô, où une barque fit naufrage, il ne périt qu’un homme, mais une quantité
d’effets de compagnie furent à jamais perdus. Je montai avec mon détachement sur une de ces barques,
à son passage à Chivasso, qui est situé au voisinage de ce fleuve ; nous arrivâmes fort heureusement
à Casal, où le marquis de la Tore invita le corps d’officiers d’aller passer la soirée chez lui, je
m’y rendis avec plusieurs de mes camarades, et nous vîmes une des plus brillantes assemblées du
Piémont. Le 3ème de décembre nous étions à Alexandrie, où il fut décidé que le 1er bataillon se
rendrait à Voghera, il y arriva le 6. Je m’arrangeai d’abord avec Stettler pour nous loger ensemble,
nous trouvâmes un plain-pied à fort bon prix

\marginpar{\footnotesize\raggedright 1798.} et les premiers mois de notre séjour dans cette ville,
se passèrent très agréablement. Au Nouvel An 1798 les officiers du bataillon se réunirent tous chez
M. le Colonel Zehender nous y eûmes un excellent dîner ; la fête fut un peu troublée par le mauvais
effet que le vin produisit chez Stettler et Zehender surtout ce dernier était dans un état
dangereux. Pendant le carnaval il y avait des bals masqués où les jeunes officiers se trouvaient
tous ; j’ai eu un soir des désagréments avec des employés français qui étaient masqués en
perruquiers, et nous poudraient les habits, l’affaire aurait pu devenir sérieuse, si le Gouverneur
ne s’était donné la peine de tranquilliser les parties. L’hiver était superbe, tous les jours nous
allions nous promener en corps, Stettler et moi sommes allés quelquefois à Casteggio où le cousin de
Gotstadt était en détachement, il y donna un bal qui fut très joli. C’est dans ce temps que nous
recevions de Suisse les nouvelles les plus alarmantes, tous les courriers la tranquillité dont nous
jouissions à Voghera était interrompue par des relations de nos camarades qui se trouvaient en
Suisse et qui voyaient de près l’orage qui grondait sur notre chère patrie. Le machiavélisme des
Français, leurs manigances voilées des démonstrations les plus amicales, sous l’aspect le plus
flatteur, séduisirent une partie des esprits dans le corps même du gouvernement de Berne ; le pays
de Vaud, par le secours d’un traître (La Harpe) qui était accueilli et soutenu par le Directoire de
France, fit sa révo–lution ; le Général Weiss que le gouvernement envoya dans cette contrée pour
maintenir l’ordre, ne servit qu’à produire le contraire. Les autres cantons de l’Helvétie loin de
manifester un intérêt commun dans une époque aussi généralement dangereuse la saisirent pour
satisfaire les sentiments de jalousie que la puissance et le bonheur de la république de Berne leur
avait depuis longtemps suggérés ; ils ne voulurent point consentir que leurs contingents marchassent
au pays de Vaud, ils l’envisagèrent comme ne faisant aucunement partie du territoire helvétique, et
l’abandonnèrent ainsi à son propre sort. Le patriotisme le plus pur animait les officiers du
bataillon de Voghera ; heureux de verser leur sang pour leur patrie, ils résolurent d’adresser un
mémoire à Leurs Excellences pour les prier de nous permettre à voler à son secours ; je fus chargé
de la rédaction de ce témoignage solennel de notre entier dévouement, et je le transmets ici, tel
qu’il a été expédié.

\marginpar{\footnotesize\raggedright 1798.} Titulierten Grössten den Wagen als die Sache zu führen,
finden den getreuesten Untertanen des Regiment Stettler, ein scharfe Ausdruck ihren hohen Obrigkeit
die die aufrichtigsten Empfindungen der treuen Soldaten und den unüberwindlichsten Gegenstand wohl
an den Tag zu legen. Die äusseren und inneren Gefahren, die allen Nachrichten zufolge der Freiheit
und Sicherheit des teuren Vaterlandes drohen, erwecken in uns allen den feurigsten Wunsch, zur
Verteidigung desselben zu eilen. Wir dürfen umso mehr auf die Erfüllung desselben hoffen, da nach
dem Vertrag, den Euer Titulierten mit dem Könige von Sardinien haben, Hochdieselben das Regiment auf
Euer Titulierten Frage sogleich verabfolgen lassen muss. Wenn es also der Wille des Allmächtigen
ist, uns Krieg auf der Schweiz Ortsteil sein soll, so sind wir einer für alle willig und bereit, den
schönen Tod fürs Vaterland, uns trösten, gnädige Obrigkeit, zu sterben, und haben erwägen dürfen,
Euer Titulierten hier unsere Empfindungen und Gesinnungen aufs segentlichste anzugeloben. Wir enden
mit inbrünstigsten Bitten zu Gott, Er möge die weise und gütige Regierung Euer Titulierten ferner
mit seinen Tagen krönen, und deren getreuen Untertanen bis in die härtesten Zeiten die Ruhe, das
Glück und die Wohlfahrt gewähren, die sie unter Euer Titulierten liebreichen Zepter seit schon so
langen Jahren reichlich geniessen. Pénétrés des sentiments que nous confessons dans ce peu de
lignes, mes camarades me priè́rent de partir pour Alexandrie où se trouvaient les autres bataillons
du régiment ; de faire signer par eux la lettre que nous voulions adresser au souverain, et de les
aiguillonner au même zèle si cela était nécessaire. Arrivé à Alexandrie plusieurs de mes amis,
particulièrement le ministre de Camp et F. Zehender applaudissaient à ce que je leur proposai ; mais
je trouvai parmi la plupart des officiers moins d’empressement que je me flattais ; Groos et
Bourgeois quoique du Pays de Vaud, se distinguaient néanmoins pour faire cause commune avec les amis
de Voghera. Après bien des peines et mille désagréments qui seraient trop longs à détailler je
parvins cependant à faire signer le mémoire, M. Collet seul prit sa démission lorsqu’on lui présenta
la lettre ; elle partit à l’adresse de M. le Brigadier Stettler à la fin de février, et n’arriva à
Berne que le 17 mars ; époque où notre pauvre ville était déjà au pouvoir de l’ennemi. La chaleur
que j’ai mise à opérer la réussite de mon projet, m’a tiré plusieurs camarades à dos, les uns
étaient assez mauvais sujets pour détester leur légitime gouvernement

\marginpar{\footnotesize\raggedright 1798.} d’autres préféraient leurs intérêts pécuniaires aux
sentiments d’honneur, ainsi leur perte ne pouvait m’être sensible. Après la lettre partie je
retournai à Voghera ; pendant mon séjour à Alexandrie, le cousin Bucher m’avait comblé d’amitiés,
j’étais logé et nourri chez lui, sa jolie Catherine faisait le ménage, elle prenait aussi une
sensible part aux malheurs qui menaçaient la patrie. Incertains sur notre sort nous étions à Voghera
dans les plus cruelles inquiétudes ; je conso–lai mes amis avec l’idée que les Français à force
d’argent renonceraient à faire la guerre à la Suisse ; Watteville ne voulait point que la nation
s’avilisse à acheter la paix, il préférait la mort à ce qu’il appelait déshonneur ; je croyais par
contre que dans un moment où ce noble sentiment n’animait pas la multitude, il serait plus sage de
saisir tous les moyens possibles pour éviter une rupture avec la France qui alors n’avait point
d’ennemi déclaré et dont les forces étaient d’une supériorité contre laquelle il était extravagant
de vouloir lutter, dans la situation désastreuse où se trouvait la Suisse à cette époque. Tout le
monde connaît l’issue de la catastrophe, nous l’apprîmes le 18 mars par une lettre de M. le
Brigadier Stettler au Colonel Zehender qui nous fit le détail de la mort de son frère M. le Bailli
de Bipp. On me chargea d’annoncer cette funeste nouvelle à mon ami, son fils ; la perte du meilleur
des pères, d’une maniè́re aussi cruelle, car il fut assassiné par ses gens, l’affecta vivement, tous
les amis se réunirent auprès de lui pour pleurer avec lui les malheurs de la chérissime patrie. Je
reçus consécu–tivement des lettres qui contenaient la suite des événements qui se sont succédés
avant et au moment de l’entrée des troupes françaises dans la ville de Berne. Mes parents
m’exhortaient à la résignation et à une confiance aveugle dans les décrets de la Providence ; ma
mère me dit dans une de ses lettres : „On n’est jamais entièrement malheureux par la chute“ ”méritée
de quelques vains honneurs et une diminution de ses biens, quand on „n’est dévoré ni par l’ambition
ni par la soif des richesses, et qu’on apprend „ à se contenter de peu. Je me faisais dès lors mon
propre système ; je me rangeai sous la loi de la nécessité, tout changement quelconque dans l’état
me paraissait impossible, et je croyais qu’il était du devoir de tout honnête homme de faire ses
efforts pour diminuer les malheurs qu’une révolution funeste traîne toujours à sa suite. Je me
flattais qu’une fois que les rênes du nouveau gouvernement se trouveraient entre les mains d’hommes
sages, nos

\marginpar{\footnotesize\raggedright 1798.} maux iraient à leur fin. Ces sentiments étaient ceux de
la majorité des officiers raisonnables, les nouvelles de la Suisse nous intéressaient toujours
vivement, nous recevions de temps en temps les Annales de Haller, qui représentaient le nouvel état
des choses sous son vrai point de vue blâmant d’un côté ses grands défauts et applaudissant de
l’autre à ses avantages. Éloignés du foyer de tous ces changements nous étions moins faits à
apprécier leur heureuse ou funeste influence que ceux qui se trouvaient sur les lieux mêmes ; on
nous assurait que tout gouvernement quelconque avec des hommes sages, vertueux et désintéressés
était bon, et nous passions notre temps dans la douce espérance d’un avenir plus heureux. Les
Français non contents du désespoir où ils avaient plongé les Suisses, voulaient encore révolutionner
les heureux habitants du Piémont ; il est hors de mon plan d’étaler ici les infamies que les
gouvernants atroces de cette grande nation mirent en feu pour parve–nir à leur but ; je ne touche
qu’à celles qui me regardaient personellement. Au commencement d’avril les agents français se
servirent de tout ce qu’il y avait de gens abjects en Piémont pour former dans le sein de l’État une
armée révo–lutionnaire ; je reçus ordre de partir avec la compagnie des Grenadiers que je commandais
en absence du capitaine, pour marcher contre ce que nous appelions les brigands ; je quittai Voghera
le 18 avril, le même jour j’arrivai harassé de fatigues au Bosco, d’où je partis le lendemain pour
le Cairo, j’y restais jusqu’à la fin du mois et en repartis pour me rendre au Castellazo par Spigno
et Acqui. Le 1er de mai j’arrivai au Castellazo, et j’y joignis une grande partie du régiment. Le 6
de ce mois une femme me vola mon portefeuille avec deux cent cinquante livres quelque temps aprè́s
on la découvrit et 100 livres tournois furent sauvés. Dans ce temps je me trouvais très malheureux ;
la situation du Roi de Sardaigne devenait de jour épuisaient le trésor qui était déjà fort mal garni
; plusieurs officiers mais surtout Camarés le cadet toujours l’ennemi des Bernois me tint des propos
très désobligeants au sujet du mémoire que j’avais fait partir, pour à ce qu’il disait, casser le
coup aux officiers du pays de Vaud. Les circonstances furent cause que nous ne nous battîmes pas,
plu-sieurs amis me prièrent d’oublier ces injures et de conserver mes jours pour la défense de mon
souverain ; mécontent, découragé, je pris la plume pour demander à mes parents le consentement de
prendre ma démission. Je leur remarquai que le roi de Sardaigne serait détrôné avant trois mois et
qu’il serait disgracieux d’être estropié par ces brigands, sans n’avoir rien à attendre de la
reconnaissance royale.

\marginpar{\footnotesize\raggedright 1798.} Livré aux plus sinistres réflexions je reçus l’ordre de
partir pour Serravalle avec mes grenadiers. Le nouveau et joli pays, l’éloignement de ceux qui
avaient causé mes chagrins, et l’arrivée de mon ami Stettler qui vint prendre le comman–dement de ma
compagnie servit à me distraire ; j’oubliai le passé et je résolus de partager le sort du régiment
dans ces circonstances pénibles, une lettre de M. le Brigadier et la réponse de mes parents
servirent encore à me décider davantage. Au commencement de juin Stettler et moi fûmes détachés sur
les montagnes de Gênes, nour passâmes six jours au bivouac par un temps affreux, le particulier de
Serravalle chez lequel j’avais été logé me combla de bontés en m’envoyant des vivres de tout genre à
mon poste. La guerre avec les brigands se faisait tous les jours de la manière la plus étrange,
protégés clandestinement par les Génois et les Français, il n’y a pas d’outrage que le roi n’ait
essuyé de la mauvaise foi des derniers ; l’argent sut encore retarder un moment la chute du trône,
le Général Brune alors à l’armée d’Italie voulut bien em–ployer ses bons offices pour terminer ce
qu’il appelait les différends entre la république de Gênes et le roi de Sardaigne ; sous ce prétexte
il mit des troupes françaises au Fort de Serravalle, qui avait été assiégé par les Génois, les
brigands et les Français et que le gouverneur piémontais rendit par capitu– lation aux premiers,
après trois jours de siège. Les négociations de Brune nous procurèrent l’avantage de quitter le
bivouac le 7 de juin, nous finîmes ce mois dans les bons villages de Montaldo, Carpenetto et
Frugarolo ; dans ce dernier j’étais logé avec Armand mon Sous- lieutenant dont j’ai fait un ami
depuis ; nos lits dans ce cantonnement étaient si remplis de puces et de poux que nous quittâmes nos
chambres pous aller coucher sur la paille. C’était au commencement de juillet que les brigands se
réunirent de nouveau pour nous inquiéter. Le 4 du mois le gouverneur d’Alexandrie fit prévenir
l’armée, qu’ils tenteraient un coup sur cette ville ; nous marchâmes à la petite aube du jour pour
Spinetta où, d’abord arrivés on nous dit qu’il n’y avait pas un ennemi ; on se moqua du vieux
gouverneur et s’en retourna à Frugarolo. Au 5 de bon matin la même alerte vint nous réveiller, on se
mit sur-le-champ en marche, cette fois nous fûmes plus heureux ; les bri–gands étaient en effet
seulement à une lieue d’Alexandrie où la garnison les attendait pour faire la révolution ; ce projet
fut déjoué, nous tombâmes

\marginpar{\footnotesize\raggedright 1798.} sur l’ennemi la baïonnette au bout du fusil, il était au
nombre de 1200 dont 300 mordirent la poussière et 600 furent faits prisonniers ; on eut lieu à sa
repentir de la générosité à laisser la vie à des traîtres ; les Français les prirent sous leur haute
protection, les firent conduire à Milan, par la 69ème demi-brigade dont les chefs permirent, que ces
scélérats sur leur route dans les États du roi, chantassent les chansons les plus révolutionnaires,
criassent des sottises pour mettre l’épouvante dans les coeurs des honnêtes piémontais. Le 20
juillet je rejoignis le bataillon de Voghera, content d’avoir heureusement fini cette infâme
campagne ; je passai l’été fort agréablement, je fis la connaissance du Baron de la Chambre
capitaine aux Dragons de Piémont, qui était en détachement avec nous, il n’a pas peu contribué à me
rendre le séjour de Voghera fort agréable. Je m’occupais dans ce temps à traduire les lettres de la
Marquise de Lambry que la Chambre m’avait prêtées, je lisais aussi les campagnes du Maréchal de
Saxe, et je jouais beaucoup au trictrac ou aux poches ; ainsi finit le mois d’août. En septembre la
chasse aux oiseaux, et les promenades à cheval dans les vignes, avec la Chambre, variaient mes
plaisirs. Le 1er d’octobre nous partî–mes pour Alexandrie pour nous réunir aux deux autres
bataillons ; la situ–ation de nos affaires politique continuait à être très malheureuse ; on se
flattait toujours que les Français n’auraient par assez peu de bonne fois, pour révolutionner le
Piémont, après les promesses solennelles que ses ambassadeurs étalaient à la cour, avec une loyauté
apparente ; ils s’emparèrent dans ce temps de la citadelle de Turin. Je fis d’agréables
connaissances à Alexandrie, je fréquentais souvent la maison Ghilini, où la marquise me témoigna des
politesses, je passais mes soirées chez Donna Laura ou au café avec mes camarades ; j’ai aussi fait
quelques courses à Valenza, pour voir mes connaissances du régiment de Zimmerman, qui étaient venues
me faire visite à Alexandrie. On s’occupait alors d’une grande réforme au régiment et dans
l’organisation des Suisses en général ; tous les colonels de cette nation furent appelés à Turin ;
la révolution de notre patrie rompait les capitulations que le roi avait contractées avec les
cantons séparément, et on voulait singer dans les services étrangers l’unité qui faisait une des
premières bases de la nouvelle constitution helvétique. L’intérêt

\marginpar{\footnotesize\raggedright 1798.} que le ministère piémontais paraissait mettre dans cette
nouvelle formation nous inspirait une espèce de confiance et nous faisait croire à la conservation
du roi. Beaucoup de vieux officiers devaient obtenir leur retraite, et faire place aux jeunes, par
cet arrangement je parvenais au grade de Capitaine-commandant avec 100 louis d’appointements ; les
premiers jours de décembre nous ap–prîmes que le Major Ernst avait été nommé chef du régiment, et
que Messieurs Stettler, Tschiffeli, Bucher, Zehender et Jonquière l’aîné étaient accommodés, cette
nouvelle nous frappa d’étonnement ; chacun combina d’avance quelle serait sa situation future et
attendait avec la plus vive impatience les nominations aux nouvelles places, lorsque le 6 du mois,
la plus infâme révolution vint trancher le fil aux plus belles espérances. La veille les Français
firent courir le bruit qu’ils partiraient le lendemain pour marcher du côté de Rome. Les officiers
de l’armée piémontaise, qui étaient en garnison à Alexandrie, prirent même congé des officiers
français de leur connaissance et ne se doutaient aucunement des projets du jour suivant. À l’aube du
jour toute la troupe se rendit sur la place, d’où on détacha une partie pour s’emparer du gouverneur
et des différentes gardes piémontaises qui étaient aux portes de la ville. Je fus surpris à la Porte
de Gênes par une soixantaine de carabiniers de la 16ème légère qui avaient à leur tête le Capitaine
Couturier ; le caporal de ma garde vint me dire qu’un détachement français demandait l’ouverture des
barrières, je m’étonnais beaucoup de ce qu’on voulait passer par ma porte pour aller à Rome, et je
fis dire au capitaine, que j’en ferais prévenir le gouvernement ; j’ordonnais en attendant d’aller
prendre les clefs à la garde de la place, et j’ouvris les barrières. Le capitaine français demanda à
me voir ; en entrant au corps de garde je lui dis que j’étais fâché de le faire attendre et que
j’avais déjà envoyé prendre les clefs de la porte ; il me répondit avec une espèce d’émotion : mon
cher ami je ne veux point sortir de la ville, il s’agit de tout autre chose ; le gouvernement
français a cru une mesure nécessaire de s’emparer du Piémont, et dans ce moment votre roi sera déjà
arrêté, le gouverneur d’ici a déja été conduit en citadelle ce matin de très bonne heure ; en
conséquence j’ai ordre de me saisir de toute votre garde. Voici mes instructions, avec une
proclamation que vous lirez à la tête de votre troupe, je désire de tout mon coeur que vous ne vous
opposiez point à tout ce que je suis obligé d’exiger de vous, pensez que la citadelle est à nous,
que les portes de la ville sont fermées, et que dix mille Français l’entourent

\marginpar{\footnotesize\raggedright 1798.} Je vous promets qu’il ne vous arrivera aucun mal, vous
entrerez au service de la République avec les mêmes avantages que vous aviez dans celui du roi. Je
sens, ajouta-t-il avec bonté, tout le pénible de votre situation, et j’exécute à regret un ordre de
cette espèce. Je sortis de mon corps de garde dans une confusion extrême, je lisais à ma garde
l’infâme calomnie que le gouvernement français crut inventer pour couvrir d’un voile sa sourde
trahison ; le soldat frémit d’horreur et posa les armes sur mes ordres ; il y avait deux caporaux et
douze hommes avec un tambour. L’officier français fit tout au monde pour adoucir l’amertume de mon
sort ; et je n’ai qu’à me louer de ses procédés ; il rejoignit sa troupe, pour aller surprendre les
soldats qui se trouvaient au quartier. Sur les 9 heures du matin je reçus la permission de quitter
le poste, on me rendit mon épée, et me donna l’ordre de me rendre à une église où se devaient réunir
les officiers piémontais ; j’y trouvai en effet la plupart de mes camarades autour d’un flacon de
liqueur, qui devait les guérir de leur étourdissement ; chacun racontait la manière dont il avait
été surpris, et tous assuraient qu’ils ne s’y attendaient pas ; je pouvais en dire autant pour mon
compte, car le soir avant cette fameuse catastrophe j’écrivis une lettre à mon frère dans laquelle
je lui fais un tableau très avantageux de mes espé–rances pour l’avenir, au service du roi. À midi
on permit aux officiers d’aller dîner et on leur donna les arrêts chez la Marquise Pra qui avait une
vaste maison, où elle nous assigna deux chambres ; sur le soir tout le monde s’y trouve, les
officiers de Stettler et de la reine étaient pêle-mêle couchés sur des canapés, des chaises, ou par
terre ; la marquise nous fit servir de la polenta, du riz, quelques chapons et du vin, qui nous
égaya un peu ; je fis ce soir une partie au piquet avec le Capitaine Jonquière, ceux qui ne
pouvaient dormir jouèrent au pharaon ; ainsi se passèrent trois jours au bout desquels plusieurs
officiers, moi du nombre, firent porter leur lits dans la maison Pra ; on s’adressa au Général
Montrichard pour obtenir l’agrément de se loger où on jugerait à propos, il le permit quelques jours
après. Je priai Stettler de Bipp de partager son logement avec moi, pour plus d’économie, il y
consentit et nous demeurâmes de nouveau ensemble. Nous apprîmes de toutes parts que la révolution
s’était opérée de la même manière à Turin et dans les autres villes du Piémont ; le 13 décembre nous
vîmes le roi et toute sa famille à son passage à Alexandrie, Sa Majesté était escortée par 180
chevaux dont 90 Français, et 90 chevaux légers du roi avec 800 hommes

\marginpar{\footnotesize\raggedright 1798.} d’infanterie républicaine. On traita la cour avec égard,
partout sur son passage les drapeaux français saluaient le carrosse où se trouvaient le roi et la
reine. À Alexandrie les militaires se rangèrent sur la place, et firent une profonde révérence à
tout ce qui appartenait à la famille royale. D’abord après la révolution du Piémont, le gouvernement
helvétique envoya deux commissaires, pris dans le Corps législatif, pour procurer aux régiments
suisses un établissement favorable, au service de la République française ; on nous déclara corps
helvétiques, et nous arborâmes notre cocarde nationale avec beaucoup de plaisir, pour quitter la
française, que nous avions prise avec tant de répugnance ; mais nos régiments étaient presque
dissous, les Français engageaient impu–nément nos gens, qui par manque de pain étaient forcés de
prendre parti ; un ordre intima aux corps républicains de nous rendre nos gens, mais il fut très mal
observé. Je crois que c’est au milieu de décembre à peu près, que nous partîmes d’Alexandrie passant
par Voghera et Pavie, pour nous rendre à Milan et de là à Codogno, je fis ce voyage en voiture avec
plusieurs camarades ; il ne s’y passa rien d’intéressant ; nous étions fort peu disposés aux
remarques agréables et utiles ; je me reproche encore à ce jour, de ne pas avoir été à l’académie de
Pavie. À Milan le temps se passa à courir d’un général à l’autre, chez Suchet chef de l’état-major
on nous fit signer un acte par lequel on s’engageait à servir la République. Nous fîmes visite à
Albert Haller qui en l’absence de son frère était ministre helvétique près la Cisalpine, il invita
trois de nos Messieurs à dîner, disant aux autres que son ménage de garçon mettait obstacle à
traiter un plus grand nombre, ce qui nous donna une petite idée de notre représentation nationale
près les puissances étrangères, d’autant plus que notre compatriote Haller était prévenu de notre
arrivée. Nous partîmes de Milan le lendemain, nous couchâmes à Lodi, on nous fit voir le pont où
Buonaparte s’était couvert de gloire, le jour après nous arrivâmes à Codogno, où je fus logé chez un
chanoine qui me combla de bontés ; si j’ai oublié son nom, je conserve

\marginpar{\footnotesize\raggedright 1798 et 99.} néanmoins la plus sincère reconnaissance pour ses
bienfaits. Notre séjour dans ce village ne fut pas long, nous le quittâmes pour nous rendre à
Treviglio, petite ville cisalpine, on me logea dans une belle et grande maison, avec mon ami
Stettler, le maître était un chanoine, Suisse italien de naissance, qui nous témoigna beaucoup
d’amitié et nous fit des politesses. C’est là que nous célébrâmes le 1er jour de l’an 1799. La fête
n’était pas brillante, l’incertitude sur notre sort, le souvenir des mal–heurs passés et présents et
le manque de moyens nous empêchèrent de nous livrer à des amusements frivoles ; quelques flacons de
vin qui se vidèrent entre quelques amis étaient tout ce qui pouvait nous faire remarquer le Nouvel
An. Je crois que c’est au 6 de janvier que tous les Suisses furent appelés à se réunir à Cassano, où
les commissaires français et helvétiques voulaient procéder à la formation de deux légions suisses
au service de la République française, à l’armée d’Italie. Les régiments Streng, Ernst, Bachmann
Zimmerman et Peyer Im Hof, arrivèrent au jour indiqué sur la place de Cassano, et par un froid très
vif ; toute la journée les commissaires respectifs s’occupèrent de l’amalgame de ces différents
corps, les soldats qui n’étaient ni nourris ni payés, souhaitaient les deux Républiques à tous les
diables, se voyant arrachés à leurs compagnies et leurs régiments, pour être incorporés avec des
gens qu’ils n’avaient vus des leur vie ; l’officier par les mêmes raisons éprouvait le même
mécontentement ; j’ai été un de ceux qui a le plus souffert de cette nouvelle formation ; dans mon
régiment j’aurais été nommé capitaine, si le roi avait été détrôné huit jours plus tard, et dans ces
deux légions on me plaça le 36ème Lieutenant, malgré les repré– sentations que j’adressai aux
commissaires ; je les priai de m’accorder un congé pour aller en Suisse en dédommagement du
préjudice que je venais d’essuyer, ils me le refusèrent. Après que cet amalgame, vicieux sous tant
de rapports, fut achevé, les commissaires repartirent et destinèrent la 1ère légion à Treviglio et
la seconde à Bergame ; je me mis en marche pour cette dernière ville ; en chemin je fis connaissance
dans un bouchon avec un Capitaine de Streng, Wipf de Schaffouse ; arrivés à Bergame il me proposa de
l’accompagner dans une auberge où il me ferait faire connaissance avec plusieurs bons enfants de la
nouvelle

\marginpar{\footnotesize\raggedright 1799. Janvier.} légion, je le suivis et il tint parole, je pris
d’abord en affection l’aîné des Bucher Capitaine dans Peyer. Peu à peu on apprit à se connaître, on
se lia avec ses nouveaux frères d’armes, je regrettais cependant beaucoup de mes anciens camarades,
que le sort avait placés dans l’autre légion, surtout Stettler de Bipp. Le hasard me destina un fort
bon logement à Bergame chez Don Salvagni, les gens de la maison étaient des plus honnêtes et
s’efforçaient à m’obliger dans toutes les occasions. Le 21 janvier le Général Moreau qui était dans
ce temps Inspecteur de l’armée d’Italie, ordonna à la garnison de Bergame de faire la parade ce
jour, pour célébrer l’horrible assassinat de Louis XVI et prêter le serment de haine à la royauté !?
Je ne veux peindre les sen–timents qu’éprouvèrent les Suisses qui furent appelés à cette infâme
fête, je sais qu’après un discours analogue à la circonstance, le Général Moreau cria haine à la
royauté, vive la République, aucun de nous n’ouvrit la bouche*. Quelques jours après je fis
connaissance avec *à l’exception d’un cri de petite voix moqueuse à dents serrées : vive la
République la belle-soeur de mon maître de maison ; cette femme eut la bonté de prendre de
l’inclination pour moi, elle me vit pour la première fois ce grand jour de parade et pria son mari
de m’introduire chez elle. Un soir ce galant homme me dit au café qu’il était le frè́re de mon hôte
et qu’il serait charmé de me voir chez lui ; je le remerciai de sa politesse et lui promis d’en
profiter ; le jour suivant il me vit à la même heure, il avait amené sa femme au café, suivant
l’usage du pays, et me présenta, me proposant de la conduire chez elle, pour y passer la soirée.
Nous sortîmes tous les trois, et nous rendîmes à la maison de Don Salvagni ; on se plaça autour de
la cheminée, deux autres messieurs de Bergame y vinrent quelques moments après et firent avec nous
la belle conversation. La jeune dame témoigna beaucoup de plaisir à faire la connaissance d’un
Suisse, elle me demandait si j’étais bon républicain, et lorsque je l’assurai du contraire, elle me
dit qu’elle était charmée de voir que je pensais comme tout ce qui appartenait à leur maison ; non
seulement n’aimons-nous pas la république, mais nous désirons les

\marginpar{\footnotesize\raggedright 1799. Février.} Autrichiens avec la plus grande impatience. Sur
les 9 heures je me retirai ; en sortant de la porte, le domestique de la maison m’arrêta et me dit
que sa maîtresse lui avait ordonné de me prier de venir chez elle le lendemain matin à 11 heures,
qu’elle serait seule ; je lui répondis que oui. Arrivé chez moi je réfléchissais à ce qui venait de
se passer, fermement résolu de ne pas manquer un rendez-vous dont je me promettais mille charmes ;
l’objet qui devait me les procurer ne m’avait cependant point inspiré de sentiment tendre. À l’heure
indiquée je me rendis à mon poste, Madame me reçut avec transport, elle me supplia de garder le
secret sur ce qui arrivait ; vous avez, me dit-elle, une si grande ressemblance avec un noble
vénitien que j’ai beaucoup connu, qu’au moment où je vous vis, Vous m’inspirâtes de l’amour, excusez
ma faiblesse cher Suisse et aimez moi comme je vous aime. Ce disours prononcé en bon italien, avec
une vivacité extrême, m’était fort étrange et nouveau, mais il ne me déconcerta pas ; le caractère
des Italiennes m’était connu et je sus d’abord quelle conduite tenir ; Mme Angélique trompa mon
attente, elle ne voulut point m’accorder ses faveurs, et me disait qu’elles ne seraient que la
récompense d’une constance longtemps éprouvée ; il me fallut attendre encore trois jours, après
lesquels je fus très heureux, et ce bonheur continua jusqu’à l’époque de mon départ. Jusqu’ici
j’avais eu fort peu de commerce avec les femmes, à Turin j’allais très rarement aux filles, à
Chevasso je vis une seule fois une femme du régiment de Christ, à Voghera, en campagne et à
Alexandrie j’étais toujours d’une sobriété extrême. Mais à Bergame je m’en dédommageais un peu je
voyais la belle Angélique, deux fois par jour ; âgée de 18 ans, supérieu–rement bien faite, car elle
mettait souvent des habits d’hommes, de grands beaux yeux noirs, et une figure très intéressante,
était faite pour réveiller les sens les plus endormis ; M. Paolo son mari ne se doutait point de
notre intimité, il me prit en grande affection, et m’invita souvent à dîner et à souper. Je passais
à côté de cela mon temps très agréablement à Bergame, je mangeais avec Jonquière Wipf, Knecht, et
Bucher, (c’est-à-dire de Peyer) le soir nous faisions souvent des petits soupers avec le maître de
maison de Wipf, un joli homme et

\marginpar{\footnotesize\raggedright 1799. Février.} qui politiquement pensait comme nous ; combien
d’heureuses soirées se sont passées à l’auberge du Chat, que de jolis soupers, tant entre hommes,
qu’avec des dames de notre connaissance, nous y avons fait. L’unique désagrément que j’ai éprouvé
alors provenait de mon insensibilité pour Angélique, il m’était impossible de l’aimer et il fallait
néanmoins le lui faire croire. au mois de février j’allais souvent danser au théâtre des
marionnettes, j’y fis la connaissance d’une petite Suissesse, femme d’un culottier et libertine à un
point qu’il suffisait de lui dire de sortir du bal pour se faire suivre ; en général j’ai été
scandalisé de l’indécence de ces parties de danse, je n’avais rien vu encore qui approche ce genre ;
Zehender le Jäger et Zingg me rappelleront toujours les souvenirs les plus bizarres. Le 11 du
courant à 4 heures du soir au café, je regardais jouer au 21 ; une remarque que je fis au banquier
irrita un chef de bataillon de la 31ème légère, un certain Favre, joueur ; il me dit une
impertinence, je convins de mon tort d’avoir parlé au jeu, et lui jetai un regard qui lui disait que
j’étais peu disposé à souffrir une insulte, il me comprit et nous sortîmes. Je courus à la maison
chercher mon épée, en arrivant je n’y pus entrer, le domestique avait emporté la clef, je demandai
un passe-partout au maître qui n’en eut point ; ce retard me mit en fièvre ; les regrets de faire
attendre un officier républicain qui pourrait me taxer de poltronerie, la nuit encore qui approchait
à grands pas augmentait mes inquiétudes ; à la fin le domestique arriva et je pris mon épée. Je
joignis en courant mon ennemi, qui sortait de la ville avec son second et le mien, que j’avais
envoyé en avant, pour prévenir mon adversaire de l’accident qui m’était survenu. À peu de distance
de la ville, nous nous plaçâmes sur un champ dont je me rappellerai la mollesse, il pleuvait et
neigeait par un gros vent et on se voyait à peine ; pendant un bon quart d’heure nous nous chassions
sur ce champ tendre sans nous faire de mal ; je me battais généreusement, car toutes les fois que
j’étais près de lui, il se disait désarmé ou voulait se reposer, à quoi je consentis toujours ; je
reçus une légère blessure à la cuisse dont il n’en sut rien, je lui fis de mon côté plusieurs
égratignures aux

\marginpar{\footnotesize\raggedright 1799. Février.} mains, mais il ne s’en plaignit point. Le
mauvais temps et la nuit enga–gèrent les secondants à nous faire des propositions de paix, je dis à
mon ami Schumacher, que je ne céderai que sur l’invitation de l’adver–saire même ; après un nombre
de difficultés réciproques, le chef de bataillon me dit : je vois que vous êtes brave comme un
Suisse doit l’être, et si vous êtes content je le suis aussi ; nous nous embrassâmes et reprîmes le
chemin de la ville. À mon retour chez moi, j’y trouvai plusieurs de mes amis qui trè́s impatients
sur l’issue du combat étaient venus m’attendre, mon maître de maison, sa femme, les enfants avaient
une joie extrême de me voir arriver sain et sauf. Quelques moments aprés, je reparus au café où tout
le monde me prouvait par un intérêt marquant, que le public avait fait des voeux en ma faveur ;
Favre vint aussi, il joua toute la soirée sans tirer les gants, ce qui donnait à croire qu’il avait
été légèrement blessé. Cette affaire, quoique très insignifiante me donna du relief aux yeux des
Bergamasques, qui n’auraient jamais cru qu’un officier qui ne fut Français oserait se battre avec un
militaire de cette nation, qu’on craignait autant qu’on la détestait. Mon Angélique surtout en
conçut une satisfaction réelle, elle me dit le lendemain, mon cher tu es courageux, s’il était
possible je t’en aimerais davantage. Je voyais encore à Bergame un capitaine français aussi de la
31ème légère, j’allais souvent chez sa femme qui était fort laide, avec beaucoup d’esprit ; elle
avait une petite fille de 9 ans qui était charmante ; le nom de cette famille m’est échap–pé. Petit.
Le séjour de Bergame a été pour moi aussi agréable que possible, je m’y suis mis au nouvel uniforme
helvétique, que nos commissaires nous avaient choisi, et qui était bleu avec les revers jaunes,
paramans de même avec une battelette verte, collet et doublure rouges boutons en or avec le n°2. Le
12 de février en revenant de la parade je me rendis au café où je vis une affiche par laquelle le
Corps législatif de la République helvétique demandait un interprète italien, auquel on donnerait
150 louis d’appointements annuels

\marginpar{\footnotesize\raggedright 1799. Février.} Quelques jours auparavant, la faveur du chef de
légion et le sort m’avaient destinés à rester au bataillon de réserve, qui, à ce qu’on prétendait,
passerait le reste de l’hiver à Milan ; je conçus dès lors l’idée de m’en aller en Suisse, et de
quitter un service si contraire à mes principes, et où je me trouvais fort arriéré pour l’avancement
mon projet n’était cependant point d’aller battre le pavé chez moi, je me proposais plutôt de voir
par mes yeux le nouvel ordre de choses et de tâcher à me placer de manière à pouvoir faire le bien,
et de tirer un revenu qui dispenserait mes parents à me secourir à l’avenir ; ma connaissance des
langues allemande et française et les principes que j’avais de l’italienne me firent voir une
possibilité d’obtenir cette place d’interprète que je ne connaissais d’ailleurs ni de près ni de
loin, mais les 150 louis qu’on y attachait m’en donnaient une bonne opinion et je résolus de faire
tous mes efforts pour pouvoir partir de Bergame ; mon ami Barthés officier du même corps me
sollicita encore, ses affaires de famille l’obligèrent de passer en Suisse et il désirait beaucoup
d’avoir un compagnon de voyage. Notre chef Ernst se prêta très gracieu--sement à consentir à notre
départ et nous remit un billet pour le Général Montrichard qui devait nous donner une permission
jusqu’à Milan, que nous obtînmes sans difficulté ; nous partîmes le 15 février après dîner, je me
trouvais l’homme le plus heureux du monde, j’emportai tous mes équipages, dans la ferme espérance
que le général et chef de l’armée d’Italie m’accorderait un congé de quelques mois, après les-quels
je ne comptais plus revenir ; j’avais assez d’argent, le quartier- maître m’avait payé jusqu’aux
derniers jours et peu de temps avant j’avais reçu quelques louis de la maison. Je pris congé de mes
amis avec un sentiment pénible, je craignais ne les plus revoir ; mon Angélique était excessivement
affligée et pour la dernière fois il fallut feindre de l’être, je l’assurai que je reviendrais au
bout d’un mois, elle me fit cadeau d’une charmante corbeille en noyaux de melon, et d’un joli
cure-dent en argent et eut la générosité de m’offrir sa bourse ! je me séparai d’elle pénétré de
reconnaissance. Un moment bien agréable pour moi était celui du départ, une quantité de mes amis se
rangèrent autour de ma voiture et m’accompagnè́rent de leurs amicaux voeux ; j’ai été

\marginpar{\footnotesize\raggedright 1799. Février.} surtout sensible aux larmes que j’ai vu
répandre à Wipf. Mon brave domestique Ruegger ne me vit point partir avec plaisir ; en quittant mon
service je l’ai récompensé de manière qu’il ne se plaindra pas de moi ; je n’ai pas manqué non plus
de manifester à la maison Salvagni que mon coeur n’oubliera jamais les bontés qu’il a reçus chez
eux. Le voyage se fit heureusement jusqu’à un village où nous quittâmes la voiture, pour nous
embarquer sur un canal qui conduit tout près de Milan. Quelle mauvaise nuit ! sur ce bateau avec une
quantité de gens de toute espèce ; nous arrivâmes enfin à Milan, un de nos compagnons de voyage nous
donna un asile, et nous mit dans une chambre à trois lits, dont il occupait un avec sa femme, nous
nous couchâmes jusqu’au jour, que nous fîmes toilette pour aller chez le général en chef, un fiacre
nous y conduisit, mais il nous renvoya jusqu’à 3 heures après midi ; après bien des peines nous
eûmes la permission d’un mois. J’étais le plus content du monde, le terme m’importait fort peu, une
fois au-delà des monts. Le lendemain fut encore employé à courir par la ville, nous allâmes chez le
Général Bachmann, qui avait alors ad interim le commandement des Suisses, il nous reçut très bien,
Barthés sortait du régiment de son nom. Je rendis visite à M. Haller dont j’ai parlé plus haut,
encore cette fois il ne me reçut point comme entre compatriotes on doit se recevoir quand on se
retrouve dans un pays étranger ; je pris aussi congé de Knecht, qui se trouvait alors à Milan pour
les intérêts du corps, il me témoigna des regrets de me voir quitter la légion. Le 17 au matin nous
nous mîmes en voiture avec un Vénitien capitaine au service de la Cisalpine et un bon vieux qui s’en
retournait chez lui à un village au bord du lac Majeur, nous dînâmes à Batassano et arrivâmes sur le
soir à Côme où nos compagnons de voyage nous quittèrent : pour tuer le temps nous entrâmes dans un
joli petit théâtre où il y avait un concert après la fin du-quel les jeunes gens de la ville
s’exercèrent en rhétorique. Sur les 9 heures nous reprîmes le chemin de l’auberge, on nous y traita
fort bien, elle est située au bout du lac sur una jolie place*. Le 18 nous nous fîmes conduire à
Capolago. Que de réflexions s’élevèrent dans mon coeur en entrant sur le territoire helvétique et en
y voyant ces arbres qui ne portent que des fruits de misère et de désespoir ; je m’apercevais avec
plaisir qu’au lieu du *je m’y suis retrouvé le 7 octobre 1854.

\marginpar{\footnotesize\raggedright 1799. Février.} bonnet des jacobins, du chapeau de
Guillaume-Tell, tant il est vrai que ces objets peuvent frapper l’imagination ; j’ai eu occasion de
m’apercevoir depuis que le tout est prestige. Arrivés à Capolago nous voulûmes sur-le-champ partir
pour Lugano en bateau ; quoique nous avions à traiter avec nos compatriotes, ils nous écorchèrent
d’une manière abominable, je pensais que la révolution n’avait pas amené dans cette contrée la
loyauté des moeurs ; au moment où nous nous voyions forcés d’être à la merci de ces Juifs, arrive un
bateau de Lugano, qui comptait repartir sur-le-champ et nous fîmes par ce moyen le voyage à très bon
compte ; à une heure après midi nous débarquâmes à Lugano ; l’aubergiste qui demeure sur la place de
cette jolie petite ville nous donna un cheval pour nous en aller ce jour encore jusqu’à Bellinzone,
il se chargea en outre de mes coffres, et promit de les faire passer à Lucerne. Le trajet est
infiniment désagréable et très long, nous montâmes tour à tour notre rosse, Barthés me fit une
confi–dence lamentable, sa maîtresse de Bergame dont il m’avait parlé comme d’une chaste Susanne
l’avait régalé peu amicalement, il commençait à s’en ressentir. Ce fâcheux accident mit mon
compagnon de voyage d’assez mauvaise humeur ; à Catenazzo nous fîmes un halte et donnâmes de
l’avoine à notre bête, il était déjà huit heures du soir ; sur les dix nous arrivâmes enfin à
Bellinzone, après quelques difficultés on nous ouvrit les portes ; un excellent souper qu’on nous
fit payer au poids de l’or, nous remit un peu de la fameuse course de l’après-midi. Le 19 nous
continuâmes notre route de la même manière, montant tour à tour le cheval que nous avions pris au
dernier gîte ; à Oleggio on s’arrêta pour dîner, une rencontre d’un chasseur à cheval du 9ème
facilita notre voyage, il prit toujours en groupe celui qui devait marcher à pied, c’est ainsi que
j’arrivai à l’auberge du Péage à Airolo pour y coucher. Nous y trouvâmes deux jeunes Luganais qui
allaient à Lucerne pour entrer comme officiers dans les troupes auxiliaires, que l’Helvétie devait
fournir à la France, ils continuèrent leur voyage avec nous. Le 20 de fort bonne heure on partit
pour Airolo, où après le déjeuner on se mit en route pour passer le Saint-Gothard, la montagne fut
assez bonne, on se restaura à l’hospice et avant la nuit nous arrivâmes à Urseren, nous prîmes notre
logis chez le président du district, qui nous donna un fort bon soupér, pour notre argent s’entend ;
le 21 je passai par le trou d’Urseren et sur le pont du diable, chaque pas dans ce berceau de la
vraie liberté helvétique faisait naître dans mon coeur des sensations douloureuses, que les troupes
françaises qui fourmillaient partout n’étouffaient

\marginpar{\footnotesize\raggedright 1799. Février.} point, comme il est facile à concevoir. Je pris
ce jour là les devants, Barthés avait plus besoin du cheval que moi, à Steg je dînais avec
l’officier français qui y était en détachement, et sur les 4 heures j’arrivai heureusement à
Altdorf, où je trouvai mes compagnons que je croyais fort loin de moi, à table. Ce bourg qui dans la
même année fut la proie des flammes était dans ce temps encore dans toute sa beauté ; sur le soir
nous poussâmes jusqu’a Flüelen. Le 22 nous partîmes sur le bateau messager pour aller à Lucerne il
faisait froid et nous étions charmés de nous mettre dans une chaise à porteur qui avait par hasard
été chargée sur le bateau pour la faire raccommoder en ville. Barthés me quitta à cette auberge à la
gauche du lac en partant de Flüelen, où les bateliers se restaurent toujours ; il se fit conduire à
Brunnen pour aller de là à Zoug où je trouvais son père. Le vent était si contraire que nous
restâmes plus de 12 heures sur l’eau ; arrivés à Lucerne sur les dix heures du soir nous ne
trouvâmes plus de place à l’Aigle, il nous fallut aller au Cheval blanc, où nous fûmes assez mal ;
je trouvai surtout un contraste entre l’honnêteté dont on était reçu autrefois dans toutes les
auberges de la Suisse et la grossièreté actuelle. Le 23 je fis arranger mon uniforme avant de partir
afin de paraître honnêtement dans la capitale de l’Helvétie ; je courus d’abord au Grand Conseil, où
me conduisit mon cousin Stettler de Zofingue, il me mit en même temps au fait des principaux points
de la nouvelle Constitution et m’encouragea à me placer, il parla en ma faveur au ministre des
finances Finsler, je fis mes preuves, mais on me refusa, alléguant qu’il y avait déjà tant de
Bernois que le ministre ne pourrait en prendre un nouveau du même nom d’un de ses premiers
secrétaires. J’eus le plaisir de trouver à Lucerne, mon ami Pfyffer capitaine dans Zimmerman qui a
eu le bonheur d’être en semestre pendant la révolution du Piémont ; maintenant il faisait les
fonctions de quartier-maître ; il m’offrit d’abord un billet de logement, en me disant, je sais ce
qu’il te faut, je te mettrai bien, tu seras à 7 minutes de la ville, mais d’autres choses te
dédommageront de ce désagrément. Je m’acheminai tranquillement vers le Guggi, campagne de M.
Fleckenstein M. me reçut froidement, mais Mme n’en était que plus honnête et m’invita à souper, je
m’en retournai en ville pour remercier Pfyffer de m’avoir logé chez une très jolie femme d’une
vingtaine d’années et je me trouvai heureux de commencer mon nouveau séjour sous de si bons
auspices. Le soir je voyais au café une quantité de Bernois plus ou moins de ma connaissance, des
jeunes gens fils de nos anciens magistrats, employés dans tous les bureaux des différentes autorités
et qui me parlaient du nouvel ordre des choses comme très supportable et bon même sous différents
rapports. L’interprète français du Grand Conseil, M. Sprungli, avant la révolution officier en
Hollande me disait de ne pas perdre de temps pour me présenter pour la place d’in–terprète italien,
il m’assurait qu’au commencement il avait été aussi timide que moi pour faire ses preuves, que
l’habitude seule d’ailleurs facilitait cet emploi, et m’offrant tout ce qu’il ferait dans son
pouvoir pour me faire à ma nouvelle vocation ; cette place ajouta-t-il, n’est pas à bien loin aussi
pénible que la mienne ; les représentants italiens n’exigent que rarement la traduction des
discussions, afin de ne pas perdre trop de temps, tandis qu’il n’y en a pas une que je ne sois
obligé de traduire. Tous ces motifs, les conseils de mes compatriotes qui se trouvaient alors à
Lucerne et l’appât des 150 louis me décidèrent à me mettre sur les rangs ; j’allai avec Sprungli
chez le Président du Grand Conseil, Schlumpf, pour demander l’accès, qui me fut honnêtement accordé.
Le soir auparavant je composais un petit discours dans les trois langues pour le tenir à l’ouverture
de la séance, je m’y rendis le 28 à 9 heures du matin, je savais très bien par coeur ce que je
voulais leur dire, mais à ma honte je perdis la tramontane et fus obligé de tirer mon papier de la
poche pour le lire ; ce pas d’école me décontenança pour le reste de la matinée, plusieurs
représentants m’encouragèrent néanmoins et je repris un peu mon assiette ; on me fit traduire
quelque chose et je m’ acquittai médiocrement. Plusieurs jours de suite je vins passer la matinée en
Grand Conseil, il y avait encore avec moi un M. am Rhyn ancien bailli faisait aussi ses épreuves
pour la place d’interprète, il ne s’y prit pas mieux que moi, mais il avait pour lui son âge, son
ancien rang et une nombreuse famille sans fortune. À vue du pays il me paraissait et Sprungli
m’assurait même que la majorité du Corps législatif était en

\marginpar{\footnotesize\raggedright 1799. Mars} ma faveur. D’un autre côté mon ami Clavel qui était
alors à Lucerne me proposa d’entrer au bureau de la guerre ; Barthés qui était venu dans cette ville
et qui voyait plusieurs représentants me dit que quelques-uns d’eux ne m’aideraient pas pour la
place que je sollicitais et que surtout ceux qui avaient été commissaires à l’organisation de nos
légions ne me voulaient pas du bien. Ces obstacles et les difficultés attachéees à la charge
d’interprète italien me décidèrent d’y renoncer ; Clavel m’invita à dîner, il faisait ordinaire avec
Carrard, je fis la connaissance du ministre de la guerre pendant le repas et le lendemain on me
reçut dans son bureau comme secrétaire rédacteur. Je ne disais à personne que j’étais placé audit
ministère ; les représentants qui me favorisaient me demandaient pourquoi je ne venais plus faire
mes preuves ; Secrétan surtout se loua infiniment des procédés de mon père qu’il avait reçus sous
l’ancien régime et il me dit qu’il ne négligerait rien pour m’en témoigner sa reconnaissance, il
m’assura même que la plupart de ses collègues étaient disposés à me donner leur voix. Ces
démonstrations amicales m’embarrassèrent dans un moment où j’avais déjà renoncé à ma prétention, je
disais à M. am Rhyn qu’il devait par une lettre me prier de ne pas me mettre sur les rangs avec lui,
il le fit et je montrai ce billet à tous ceux qui voulaient savoir pourquoi j’abandonnais le projet
de devenir interprète. Me voilà donc à la 1ère division du bureau de la guerre où je me trouvais bec
à bec avec Jolimai ci-devant officier dans notre corps, et qui par une suite d’inconduite avait été
obligé de prendre sa démission ; un jour que je dînais chez le chanoine ...... à Lucerne, avec Ernst
et Sprungli, le Capitaine de Bon, et les deux frères Snell me demandèrent des renseignements sur la
conduite de Jolimai, je leur fit sentir que je préférerais ne pas entrer en matière, ils me
comprirent, et n’eurent rien de plus empressé que de faire les rapporteurs ; Jolimai m’envoya un
cartel, je me rendis au rendez-vous avec Barthés mon secondant et un jeune secrétaire fut celui de
Jolimai ; j’avoue qu’à mon regret on vint à bout de nous pacifier, nous

\marginpar{\footnotesize\raggedright 1799. Mars.} nous en retournâmes en ville, on but une bouteille
au couvent des Urselines, et depuis j’ai vécu amicalement avec Jolimai. Je me trouvais fort bien au
bureau de la guerre, Lanther qui était alors le chef de la 1ère division paraissait me vouloir du
bien, et le ministre Répond me témoignait de l’amitié ; on me reprocha souvent manque de pru–dence
dans mes propos, qui prouvaient à l’évidence que je n’étais pas chaud partisan du nouveau régime.
Tous les soirs les Bernois se réunissaient au café, me trouvant alors avec mes compatriotes je
croyais pouvoir me livrer entièrement, mais il s’est trouvé qu’eux-mêmes désapprouvaient ma
franchise ; je leur faisais naturellement ma confession de foi, je témoignais mes regrets sur la
perte de notre ancien bien-être, je détaillais les vices de la nouvelle constitution qui par le
système représentatif nous mettaient des idiots à la tête des premières autorités, je blâmais les
extravagants qui frappaient de mépris tous ceux qui ne courbaient pas leurs têtes devant la déesse
de la liberté, que les Français paraissaient avoir chassée de chez eux, pour porter le désordre
parmi nous ; d’un autre côté je ne désirais point les troupes autrichiennes et avec eux la
contre-révolution ; je jugeais d’après Lucerne, le reste de la Suisse ; dans cette ville on
paraissait en général vouloir se faire au nouvel ordre de choses, et on redoutait l’arrivée des
coalisés, qui porteraient le fléau de la guerre sur le territoire helvétique. Après avoir fait les
détails de ma vie publique, je viens à ceux de ma vie privée ; j’ai déjà dit plus haut que mon ami
Pfyffer m’avait procuré un très agréable logement hors de ville ; un des premiers soirs que je m’en
retournais chez moi, je me préci--pitai de mon long dans le canal qui se trouve d’abord hors des
portes, j’arrivai à la maison dans un état pitoyable et mes hôtes se donnèrent toutes les peines
pour me sécher. Rien n’était plus bizarre que le ménage de M. Fleckenstein, ci-devant sénateur de la
république de Lucerne âgé de 28 à 30 ans, et marié avec Catherine Gilli, fille d’une petite
bourgeoise de la ville qui tenait un café. Mme Fleckenstein avait à peu près 22 ans

\marginpar{\footnotesize\raggedright 1799. Mars.} fort jolie et d’une coquetterie extrême, qui
poussait son mari à une jalousie peu commune elle avait autant d’amitié pour les Français que son
époux de haine et d’abhorration, en général le caractère de cette femme était l’opposé de celui de
son mari. Il serait im-possible de faire le tableau des scènes bizarres dont j’ai été témoin dans ce
singulier ménage ; Mme conservait toujours un attachement sincère à un capitaine français, qui avait
logé six mois chez eux et au quel Mme a peut-être sauvé la vie en prodiguant tous ses soins à une
dangereuse blessure que cet officier avait reçue à la guerre des petits cantons. Je voulais effacer
de son coeur un souvenir qui s’opposait à mes voeux de lui inspirer de l’amour ; je lui vouais le
peu d’heures que me laissaient mes nombreuses occcupations, j’employais tous les moyens en mon
pouvoir, et quand je concevais des espérances, il arrivait une lettre du Capitaine supérieur Bigarré
qui en un moment les détruisait toutes ; le mari d’ailleurs paraissait me craindre, il n’allait
jamais en ville sans obliger sa femme de l’accompagner et à la maison il ne la laissait jamais seule
dans la chambre ; une manière aussi cruelle de traiter une femme irritait Mme Fleckenstein au point
que souvent elle en faisait des reproches à son mari, qui alors se fâchait, Mme commençait à
pleurer, et il m’est arrivé quelquefois de voir verser des pleurs à tous les deux à la fois, pendant
que nous étions à table. Je passais dans cette maison près de deux mois, M. Fleckenstein était aussi
désagréable que sa femme était jolie, j’aurais certainement conçu de l’amour pour elle, si ses
manières et ses discours avaient été plus modestes et plus décents. Pendant mon séjour à Lucerne
j’écrivis régulièrement à mes parents et toutes mes lettres respiraient le désir de les voir ; j’eus
occasion de rendre quelques petits services à Benford qui par ses lettres m’assurait de sa
reconnaissance ; je reçus aussi une lettre de Mme de Toffen, qui me priait de procurer par le canal
du ministre français un passeport à une émigrée, je me suis donné toutes les peines à ce sujet sans
avoir entièrement réussi. Ainsi s’écoulèrent mes jours jusqu’au 20 avril, que je sollicitai une
permission du ministre de la guerre pour aller passer quelques jours à Berne ; je ne l’aurais point
obtenue si dans ce temps le ministre Répond n’avait pris sa démission, je profitais de cette
circonstance pour avoir un congé que Répond m’accorda le dernier jour de ses

\marginpar{\footnotesize\raggedright 1799. Avril.} fonctions. Je quittai Lucerne sans regret, malgré
les concerts et les bals qui y étaient assez fréquents ; mon chef de bureau Lanther ne voyait pas
mon départ avec plaisir ; je pris des arrangements pour partir avec M. Répond ; avant de me séparer
de Mme Fleckenstein elle me remit une épingle (à mon grand chagrin elle s’en alla en poudre pendant
le voyage)que je devais faire arranger ici à Berne, et me la recommanda comme l’objet le plus cher
qu’elle possédait, je lui fit cadeau de cette petite corbeille qu’Angélique m’avait donnée à
Bergame. La veille de mon départ j’eus un moment d’oubli le soir en m’en retournant chez moi, ce
qu’il ne m’était pas arrivé depuis Bergame ; le lendemain je me mis en route avec Répond qui
s’applaudissait de ne plus être ministre de la guerre, nous nous arrêtâmes à Sursee et Zofingue et
pour le coucher au Morgenthal ; le lendemain nous arrivâmes pour le dîner à Kirchberg. Répond
continua sa route jusqu’à Berne, et moi je restais chez mon frère. Pendant le voyage j’ai eu les
plus grands agréments de la société de mon compagnon, j’ai appris à connaître un homme, qui quoique
insuffisant pour être à la tête d’un ministère, était incapable de faire le mal et dont les opinions
politiques ne se trouvaient pas aussi éloignées des miennes que la charge qu’il avait revêtue aurait
dû le faire croire. Après deux ans d’absence j’éprouvais un plaisir sincère à revoir mon frère et sa
femme ; à peine les premiers épanchements de l’amitié satisfaits, on vint sur le chapitre de la
politique et dans l’innocence de mon coeur je déve–loppais ma façon de penser, et ma manière de voir
les choses. Les cheveux que j’avais fait couper, la chaleur que je mettais à la défense de plusieurs
membres du Corps législatif, l’espérance que je manifestais de voir le retour du bonheur et de la
tranquillité, la déclaration que je ne désirais pas les Autrichiens, mirent ma belle-soeur dans le
plus grand étonnement et dès lors je fus proclamé sinon jacobin, au moins archidémocrate. Mon frère
par contre ne paraissait pas bien éloigné de mon avis et prenait mon parti contre sa femme qui était
irritée ; en allant me coucher je m’abandonnais à un tas de réflexions sur ce qui venait de se
passer. Comment ? me dis-je, à Lucerne l’aristocratie que tu professes t’a fait des

\marginpar{\footnotesize\raggedright 1799. Avril.} ennemis et dans l’opinion de tes plus proches, on
te reproche les sentiments contraires. J’arrivai le jour suivant à Berne, mes parents me reçurent
amicalement je me trouvais heureux d’être rendu à ma ville natale ; on m’en avait fait un si
déplorable tableau, que je la trouvai mieux que je ne l’espérais. Je revis toutes mes connaissances
; la plupart blâmait mes cheveux coupés et de ce que j’avais eu l’idée d’étre interprète au Grand
Conseil, on ne me pardonnait pas même d’avoir travaillé au bureau de la guerre. Malgré tant de
reproches on me reçut comme membre honoraire au cercle du Faucon, j’y allais fort souvent, et
m’apercevais bien que plusieurs jeunes gens se méfiaient de mes principes, souvent on se plaisait à
déchirer des personnes honnêtes, dont le probité des intentions m’était connue, je prenais alors
leur parti et me faisais ainsi un grand nombre d’ennemis. Les premiers jours de mon arrivée à Berne
je fis visite à Mme de Toffen, j’y trouvai Julie qui ralluma sur-le-champ des feux mal éteints ; je
fréquentais de nouveau la coterie que je voyais pendant mon dernier semestre, elle s’était accrue
des deux demoiselles Ernst de Romainmôtier ; la ville retentissait des éloges de la beauté de la
cadette je la vis pour la première fois sans en être frappé, une voix rauque et un maintien qui ne
me parut pas assez modeste pour une fille de 15 ans produisirent cette indifférence ; la soeur aînée
me plut mieux, la douceur et la modestie étaient préférables au surplus d’éclat qu’avait la cadette.
À peu près dans la même époque je fis aussi la connaissance de Charlotte Bürkli, jeune jolie petite
et riche zurichoise. Je passais une grande partie de mon temps en société de ces demoiselles, la
politique était aussi le sujet non interrompu de nos entretiens ; je m’apercevais facilement que ma
confession de foi que je croyais cependant être celle de la raison déplaisait fort à Julie ; de
nouvelles protestations de mes parents que nous ne nous conveni–vons pas et surtout un surplus
d’embonpoint qui la rendait toute ronde me refroidirent peu à peu ; la dernière preuve que je lui
donnais de mon attachement fut à l’occasion des instances qu’elle me fit de lui procurer Métastase,
je priai M. Stapfer Gross de me le prêter, je le lui envoyai avec une petite lettre délicate, et
depuis l’indifférence est venue à la place d’une inclination qui pendant fort longtemps a été pour
moi une source de Mai. Dans cet intervalle j’écrivis une lettre à Lanther, qui en vacance du
ministre de la guerre, avait été chargé du portefeuille ; je lui demandai s’il voulait que je revins
à Lucerne ; il me répondit vaguement que j’étais le maître de faire ce que je jugeais à propos ; je
sentis fort bien que c’était me dire honnêtement qu’il n’avait plus besoin de mes services ;
j’attribuai cette espèce de disgrâce à l’envahissement de la plus grande partie de l’Helvétie par
les troupes autrichiennes, qui diminuaient les travaux de ce ministère, à mes principes que Lanther
trouvait trop aristocratiques et au grief de ce que j’étais parti de Lucerne contre son gré. Je me
trouvais ainsi sans emploi, exposé à être mis en réquisition militaire et pour l’éviter je changeai
la permission du général français qui était d’un mois, et mis six à la place d’un, de manière que
personne ne pouvait s’en aper–cevoir ; cette ruse fut néanmoins inutile. À la fin de mai, le
gouvernement helvétique se sauva à Berne, nous logeâmes pendant deux jours le représentant Mohr du
Siebenthal et après M. Wetstein de Bâle, chef de la 1ère division du bureau de la guerre en
remplacement de M. Lanther, qui fut promu ministère ; j’appris de M. Wetstein que des Directeurs
même avaient trouvé dangereux que je continuasse mes fonctions de secrétaire et qu’eux étaient cause
de la manière d’agir de Lanther vis-à-vis de moi, une visite que je fis à ce dernier me prouva qu’il
avait encore de l’amitié pour moi. Le changement total des circonstances politiques produisit en
partie celle de mon opinion ; dans le temps que les succès des armées françaises paraissaient
indubitables et que même dans le lointain il n’y avait aucune lueur d’espérance de voir changer
l’état des choses à l’avantage de ceux qui avaient tout perdu par la révolution, je me soumettais
aux ordres du destin mais je ne désirais pas, comme je l’ai dit plus haut, que les troupes
impériales s’approchassent trop de nos frontières, afin que la Suisse ne devienne le théâtre de la
guerre ; maintenant que le sort des armes en avait autrement décidé, et qu’une moitié de l’Helvétie
venait d’être rendue à son ancien régime, je conçus de même l’espérance qu’un pareil bonheur
pourrait nous tomber en partage et je fis des voeux pour la prompte arrivée des Autrichiens ; non
que je me blousais au point, comme tant d’autres de mes compatriotes le firent, de croire
qu’aussitôt à l’entrée de l’archiduc Charles à Berne, l’ancien gouvernement et ses prérogatives
seraient rétablies, mais je me flattais à voir renaître dans

\marginpar{\footnotesize\raggedright 1799. Juin.} - ma patrie un ordre de choses préférable à celui
que les baïonnettes françaises nous avaient forcé d’accepter. Mais l’illusion de bonheur qu’on se
faisait à Berne de voir les coalisés a été cruellement trompée jusqu’ici, et ceux qui ne s’y fiaient
pas aveuglément sont moins malheureux que l’est le grand nombre qui ne prévoyait pas même la
possibilité d’une non- réussite de ses voeux. Après la prise de Zurich, qui eut lieu le 7 juin, on
était pendant six semaines dans la persuasion de voir un jour ou l’autre le retour de M. l’Avoyer
Steiger, qui à la révolution de la Suisse avait échappé aux plus grands périls, et a pu
miraculeusement sauver une vie qu’il voulait sacrifier au rétablissement de la République de Berne,
dont il était le chef. Le gouvernement helvétique était un sujet de mépris pour la plupart des
Bernois, qui s’imaginaient de reprendre les rênes et d’user de représailles envers les représentants
du peuple ; eux qui avaient été ménagés, à l’exception de quelques otages, se proposaient de
commencer leur règne par des fusillades et des dépor–tations, et quand on reprochait à ses jeunes
compatriotes l’immoralité et la terrible inconséquence d’une telle conduite ; on vous accusait de
modération suite d’un caractère faible et sans système, et qui ne peut être adoptée pour base quand
il s’agit de punir des démocrates. Que de disputes n’ai-je eu dans ce temps ? que d’ennemis ne me
suis-fait ? et que de méfiance on a mis dans mes principes, parce que je soutenais le droit des
gens, la morale et la vertu ? un sentiment intérieur applaudit néanmoins à la justesse de ces
raisonnements et le suffrage des gens qui pensent bien ne peut être mis en paragon avec le blâme des
sots. Le 12 juillet je reçus une invitation du commissaire du canton de Berne M. Bondeli qui me
priait de me rendre chez lui, j’y fus le 13. Il me fit la proposition de partir pour Aarau où je
devais surveiller un versement de 10 mille quintaux de foin que le canton avait à fournir aux
Français. jusqu’ici, ajouta-t-il, nos paysans furent toujours dupés par les commissaires français,
chargés à recevoir ces fourrages et d’en donner les reçus ; la Chambre adminstrative voudrait avoir
un homme de confiance pour la remise de ces 10 mille quintaux entre les mains des magasiniers
français. J’étais un moment indécis, si je voulais m’employer à cette commission désagréable

\marginpar{\footnotesize\raggedright 1799. Juillet.} - l’état de désoeuvrement où je me trouvais,
avec les encouragements de mon père fixèrent ma détermination. Je me rendis le lendemain chez M.
Bondeli, commissaire du canton qui me donna les instructions nécessaires ; le 14 je me mis en voyage
aux frais de la république, sur un petit char de réquisition ; j’avais reçu l’ordre de passer chez
le cantonnier Wyss de Brittnau, qui devait me seconder, je vis dans ce village le ministre M.
Desgouttes et son épouse, qui me témoignèrent beaucoup d’amitié, M. prêta même sa chaise pour me
conduire jusqu’à Aarbourg. Je continuai ma course avec Wyss jusqu’à Lenzbourg, où se trouvait alors
le quartier général français de là je poussai jusqu’à Brugg, pour y prendre des arrangements
concernant ma mission, après quoi je m’en retournai à Aarau. Cette capitale du nouveau canton
d’Argovie passait pour très jacobine et je ne me promettais aucun agrément de séjourner avec des
gens dont les principes étaient si opposés aux miens ; on me logea chez un M. Wydler, surnommé le
Marseillais et allié Streng à juger de ses propos il ne me paraissait nullement partisan de la
révolution mais j’ai appris depuis et encore pendant mon séjour auprès de lui, qu’il était hypocrite
; sa femme se cachait moins et sans être à la hauteur des prin–cipes du jour, elle préférait le
nouveau à l’ancien ordre de choses. Malgré la différence de nos opinions je me trouvais passablement
chez M. Wydler, la seule chose que me gênait, par un manque d’habitude, était de manger avec les
servantes ; en me mettant en pension chez mon hôte, je ne voulais pas m’y refuser pour la rareté du
fait et comme ennemi des préjugés. Rien de plus désagréable que mes occupations dans cette petite
ville, où se réunissaient un tas de patriotes que je ne voulais fréquenter dans mes heures de loisir
; après que les livraisons de foins étaient faites, j’allais me promener tout seul dans les superbes
campagnes des environs, et comme j’avais un équipage à ma disposition, qui m’était fourni par un
vendeur de fourrage, l’aubergiste du Sauvage, je faisais souvent des courses chez mon employé Wyss à
Lenzbourg où je trouvais mon cousin Stettler ci-devant secrétaire baillival, avec lequel je me
mettais à politiquer à tort et à travers ; j’ai passé chez lui des heures bien agréables, il fut
attaqué pendant mon séjour au voisinage, de terribles rhumatismes qui le mirent au tombeau au regret
général Août. À Aarau les versements de foin malgré mes sollicitations fréquentes se faisaient
lentement, et pour ne pas me laisser dans l’inaction, la Chambre administrative de Berne me créa son
chargé d’affaires, pour négocier près des autorités françaises l’entrée des sommes énormes qui nous
étaient dues pour des fournitures faites à l’arrivée. J’aurais de la peine à retracer ici les peines
inconcevables que je me suis données pour extorquer un argent légitimement dû, tout fut sans succès,
et ces commissaires croyaient accorder des grâces, quand ils mettaient leur signature sur un acte
qu’ils ne pouvaient nier de devoir. Rien de plus insupportable que les employés français, leur
hauteur, leur morgue, leur insolence ne peuvent qu’indisposer tous ceux qui ont le malheur d’avoir à
démêler avec eux. Ma désagréable mission m’a mis dans le cas de voir souvent M. Robert, qui était
alors comissaire du gouvernement et ordonnateur helvétique ; cet homme très républicain m’a toujours
témoigné de l’amitié malgré qu’il connaissait fort bien ma façon de penser ; souvent nous avons
parlé politique, chacun avec le même enthousiasme pour sa cause, si la mienne par rapport aux
circonstances exigeait plus de retenue, je ne lui laissais rien perdre de sa force ; j’ai reconnu en
Robert un homme qui voulait le bien, et nous ne différions que dans les moyens d’y parvenir ;
Zimmerli qui était son 1er secrétaire m’a déplu, déjà son extérieur est insolent et le connaît-on de
plus près, on le trouve de même ; l’autre jeune homme qui travaillait aussi chez Robert, est enivré
des nouveaux principes à un point qui le rend insociable pour les gens qui ne pensent pas comme lui
; j’ai eu dans ce temps de fréquentes entrevues avec l’ordonnateur supérieur Daru ; il est fâcheux
que cet homme infiniment instruit et d’une société agréable soit commissaire, il est à craindre que
ce métier l’abrutisse. Mais une connaissance bien différente que je fis à Aarau fut celle de M.
Hurner ci-devant capitaine d’artillerie au service de la république de Berne ; de tout ce que je vis
à Aarau cet homme seul était vraiment aristocrate ; quoique âgé de plus de 50 ans je l’ai vu verser
des pleurs sur la chute de son légitime souverain ; c’est auprès de cet ami que j’ai passé les plus
agréables heures de mon séjour à Aarau j’ai souvent dîné et goûté chez lui ; il logeait et
nourrissait 2 magasiniers français qui étaient d’assez bons enfants, Gabriel et La Motte

\marginpar{\footnotesize\raggedright 1799. Septembre.} Le 31 août ma grand-mère Herbort est morte et
a été enterrée le 4 septembre à l’âge de 90 ans. Au commencement de septembre mes affaires
m’appelaient à Bâle, l’ordonnateur français Bonnemain me combla de politesses, me donna à dîner,
mais il ne me satisfit point au sujet que j’avais à traiter avec lui et je revins à Aarau sans avoir
pu réussir. Un personnage qui m’était bien à charge et que je voyais fréquemment est ce capitaine
Lafléchéra de la 5ème demi-brigade d’infanterie cet homme un peu fol avait logé dans notre maison à
Berne, et dès qu’il apprit que j’étais à Aarau, où il restait pour se faire guérir de la gale, il
m’accabla de visites. J’aurais déjà dû dire plus haut que dans mes courses à Lenzbourg où j’avais à
faire avec Mathieu Favier ordonnateur français en chef, je voyais souvent Hünerwadel dans la maison
duquel l’ordonnateur logeait, et où je reçus des honnêtetés ; je fus même un jour rendre visite à
Hünerwadel à Ammerswil chez le ministre Wild, où il était vicaire, je fis cette partie avec mon
cousin Gabriel Stettler, qui était alors à Schaffisheim. À peu près à la même époque nous nous
proposâmes d’aller voir les avant-postes autrichiens, au confluent de la Limmat et de l’Aar ; je
vins un beau matin à Schaffisheim avec mon équipage, les cousins Stettler et Brutel s’y mirent, en
chemin nous trouvâmes Hünerwadel, qui nous joignit, et ainsi nous cheminâmes jusqu’à Königsfelden,
où nous visitâmes le tombeau de l’archiduc Albert, qui avait été tué à Wildenstein et l’hôpital
militaire français nous nous rendîmes sur la hauteur de Windisch, pour voir les positions
autrichieñes, Mme la ministre Ernst et sa fille, qui est très jolie, vinrent nous donner tous les
renseignements, finalement le ministre parut aussi, il me fit une description lamentable des
nombreuses réquisitions auxquelles il se trouvait exposé. Notre carrosse nous conduisit à Brugg d’où
après un dîner frugal, nous nous rendîmes à Rhein, et c’est de là que nous avons pu prendre une idée
de la situation des deux armées. Satisfaits d’avoir été aussi près des Autrichiens, nous nous en
retournâmes à Brugg, où on nous prépara un bon goûter à l’auberge de l’É́toile ; sur le soir
Hünerwadel resta à Lenzbourg Stettler et Brutel à Schaffisheim et moi je m’en retournai à Aarau fort
content de ma course. Dans l’intervalle de mon séjour dans cette ville, j’ai fait une promenade à
Berne, l’ennui m’avait tourmenté au point que je n’y tenais plus. Un jour à midi Wita surnommmé le
Judenbub m’apporta des papiers de Berne que concernaient ma mission, je lui dis, quand partez-vous ?
il me répondit : sur les 2 heures je résolus de faire voyage avec lui, nous partîmes en effet à
l’heure indiquée

\marginpar{\footnotesize\raggedright 1799. Septembre.} sur un mauvais char de réquisition, à 8
heures du soir nous arrivâmes à Herzo–genbuchsee, où Wita avait un cheval de relais ; comme j’avais
grande envie d’aller souper chez mon frère à Kirchberg, je laissai Wita sur le char, et montai son
cheval, sans cependant lui dire mes projets ; la bête alla si bien qu’avant 9 heures je fus à
Kirchberg où mon frère fut surpris de me voir, après que je lui dis qu’à 2 heures et demi j’étais
encore à Aarau. Le matin Wita monta son cheval, et je me mis à la diligence pour être de bonne heure
à Berne. Ce qui m’avait particulièrement engagé de quitter Aarau était une lettre de mon père, qui
m’annonçait que M. le Baron de Chambrier cherchait un jeune homme de famille qui sut les deux
langues et qui aurait du goût pour la diplomatie, et que M. le Principal Wegner s’était adressé à
mon père, pour savoir si un de ses fils avait du goût pour cette carrière. Mon père envisagea
d’abord la chose comme très avantageuse et me conseilla d’écrire à M. de Chambrier, qui dans sa
réponse me témoigna le désir de me voir, et c’est dans ce but que je vins à Berne, d’où je partis le
12 septembre pour me rendre à Neuchâtel ; j’y arrivai à 7 heures et demi du soir ; en sortant de la
diligence un laquais demanda si je n’étais point M. Stettler ? qu’il avait ordre de me conduire dans
la maison de M. le Baron de Chambrier je le suivis ; M. le Baron me reçut avec les civilités d’un
homme de Cour, il me dit que nous partirions dans un moment pour nous rendre à Cormondrèche sa
campagne, qui était à une petite lieue de la ville ; peu après le carrosse vint nous prendre, il
s’arrêta devant une autre maison, et deux messieurs montèrent, je leur offris ma place, M. le Baron
me dit, Monsieur restez je vous prie, c’est mon fils adoptif et son instituteur. Pendant le voyage
la conversation était très suivie, M. de Chambrier ministre du roi de Prusse près la Cour de Turin,
depuis 21 ans, me parla beaucoup du service de Piémont, et des agréments de ce pays-là ; il me parut
que dans un instant nous étions arrivés à notre destination ; on m’introduisit dans une belle salle,
et je pus enfin considérer et voir les personnes avec lesquelles j’avais voyagé. On ne pouvait
qu’être d’abord prévenu en faveur de M. de Chambrier cet homme respectable est de l’âge de 46 ans,
il est grand, mince et méfait, mais des yeux noirs d’une vivacité extrême et un râtelier éblouissant
qui involontairement se voit

\marginpar{\footnotesize\raggedright 1799. Septembre.} donnent à sa figure une expression qui
inspire un intérêt qui s’augmente encore par la douceur et l’honnêteté qui l’accompagnent ; les
manières de M. le Baron sont tant soit peu affectées, mais ce défaut, si c’en est un, se perd dans
le nombre d’inno centes et rares qualités qui l’accompagnent, et dont je parlerai plus bas. On
servit un bon souper, la conversation devint générale, M. Poinsieux l’insti–tuteur du jeune de
Chambrier parlait assez bien, mais comme il étudiait d’imiter la suaveur et le mielleux que M. le
Baron met avec un peu d’affectation dans sa manière de parler, tout ce qu’il disait devenait
désagréable par le ton singulier qu’il se forçait d’adopter ; le seul jugement que je pouvais porter
du jeune homme de 14 ans, fut un bon appetit. D’abord après les 10 heures chacun se retira dans son
appartement. Livré à mes réflexions je re–passai tout ce que j’avaiss vu ; le train de maison de M.
de Chambrier, où six valets étaient aux ordres du maître faisait un disparate avec le ménage plus
uni de M. Wydler, où je mangeais avec les servantes ; une observation que je n’aurais pas même voulu
faire à M. de Chambrier lui qui n’était jamais sorti de l’ordre de l’ancien régime aurait pu trouver
la chose bizarre. Je reposais dans un lit tendre, des sens qui avaient à se préparer à de nouvelles
impressions agréables. Le matin après les 8 heures on vint m’appeler pour le déjeuner, sur les 9 le
Baron me proposa de passer au salon où il saisit l’occasion de parler du sujet qui m’avait amené
auprès de lui. „Depuis longtemps, me dit-il, M. le Marquis de Gallo, ministre de Naples à la Cour de
Vienne, désirait un secrétaire privé qui sut les deux langues, comme suisse il se confia à moi pour
lui chercher un de mes compatriotes ; la malheureuse révolution qui met une infinité de jeunes
patriciens bernois dans une situation à ne plus pouvoir employer leurs talents au bien de leur
patrie, par l’aversion naturelle que doit leur inspirer le nouvel ordre de choses m’a donné l’idée
de prier M. Wagner de s’informer, si la place en question pouvait peut-être convenir à un d’eux, et
voilà, M. ce qui me procure l’honneur de votre connaissance. Après cela M. le Baron me fit un long
détail des avantages et désavantages de cette charge ; il mettait au nombre des derniers l’étiquette
des cours catholiques, la difficulté pour un protestant de voir le grand monde, et le rôle
subalterne d’un secrétaire privé

\marginpar{\footnotesize\raggedright 1799. Septembre.} sans doute ajouta-t-il, aurez-vous de bons
appointements, et vous serez attaché à un homme d’un grand mérite, je connais le marquis de Gallo
particulièrement et je puis vous répondre que vous auriez lieu d’être content de lui, d’ailleurs
vous avez tout le temps de réflechir là-dessus je ne lui écrirai que lorsque vous serez décidé ;
dites moi maintenant M. ce qu’il vous en semble ? Cette ouverture m’avait mis tout à coup au fait
dont il s’agissait, et je fus trompé dans mes espérances ; je m’ima–ginais que M. de Chambrier
cherchait un secrétaire de légation pour sa mission à Turin et je ne doutais point, qu’aussitôt
arrivé à Cormandrèche, je serais installé à cette place ; je me réjouissais d’avoir un caractère
public et de signifier quelque chose en politique ; je dis à M. de Chambrier que sous bien des
rapports il me serait agréable d’être attaché à la personne de M. le Marquis de Gallo, mais qu’avant
qu’il m’eut fait l’honneur de me parler, j’avais conçu d’autres espérances je me flattais,
ajoute-je, que vous Monsieur, auriez peut-être besoin d’un secrétaire de législation, et il m’aurait
été doux de vous accompagner à Turin. Après un moment de réflexion, il me répondit que si en effet
cela me faisait plaisir il serait possible que nous nous arrangeassions ensemble, vu que depuis fort
longtemps il n’avait reçu de nouvelles de son secrétaire de légation qui avait suivi le roi, en
Sardaigne, qu’il craignait même qu’il était mort. Comme j’aurais voulu finir l’affaire et m’engager
sur-le-champ, je lui dis que je serais charmé d’avoir quelque chose par écrit pour pouvoir dire à
Berne que j’étais attaché ailleurs, en cas qu’on voulût me prendre pour la réquisi–tion militaire,
ce que je faisais semblant de redouter beaucoup ; pour me tranquilliser il me remit un papier en
forme de lettre, où il me dit, que vu le manque de nouvelles de son secrétaire, et mes dispositions
d’entrer dans la carrière diplomatique nécessitaient de fréquentes visites à Neuchâtel etc. par ce
moyen il me donnait une autorisation d’aller me réfugier à Cormondrèche toute fois que je le
jugerais à propos ; M. le Baron ajoutait encore qu’en cas que M. Interluttner son secrétaire vint à
se retrouver, il tâcherait à l’engager de chercher une place plus avantageuse ; que lui-même
s’intéresserait à la Cour pour la lui obtenir, ce qui ne serait pas difficile, puisqu’il avait en
bien des occasions d’obliger le roi, pendant son séjour à Cagliari.

\marginpar{\footnotesize\raggedright 1799. Septembre.} Les affaires politiques de ce temps
permettaient d’espérer que sous peu le roi de Sardaigne rentrerait dans ses États, et que M. le
Baron ne tarderait pas à partir pour Turin. L’après dîner de ce jour où nous eûmes cette
expelication M. de Chambrier me conduisit en carrosse chez M. le Gouverneur de Béville, auquel il
dit qu’il projetait à m’attacher à la légation de Prusse, Son Excellence me reçut fort bien et me
disait à cette occasion des choses obligeantes. À notre retour à Cormondrèche, je convins avec M. le
Baron, que dans un mois je retournerais auprès de lui et qu’alors je partagerais mon temps entre ce
séjour et celui de Berne, jusqu’à l’époque où je m’attacherais plus particulièrement à la légation
prussienne. Je me séparai de M. le Baron le 3ème jour, il me donna son carrosse jusqu’en ville et
m’accompagna à la moitié du chemin, en me quittant il me dit : M. je ne prends point congé de vous ;
au plaisir de vous revoir.... Mon imagination était frappée des procédés charmants de cet homme
respectable, depuis plusieurs mois je n’avais vécu qu’avec des républicains, et cette manière
obligeante et délicate me rappelait si fort l’ancien régime, que je me croyais transporté au temps
heureux qui précéda la révolution. J’arrivai à Neuchâtel, je donnais un écu de six francs au cocher
et je sortis du carrosse chez Mme Bourgeois Meuron, ma soeur Dachselhofer avait été 12 ans
auparavant en pension chez cette dame et depuis elle conservait toujours avec elle et sa fille Mlle
Rosette une liaison intime, j’ai eu occasion le voir cette dernière plusieurs fois à Berne chez ma
soeur. Ces dames me reçurent le plus amicalement du monde, j’y passai la soirée avec une soeur de
Mme Bourgeois avec Mme et M. Sandoz le Secrétaire d’État, et une personne dont j’ignore le nom, mais
qui a un talent rare pour le dessin ; je fus même invité à souper, au milieu du repas je vis arriver
le jeune M. du Roveray de Genève grand ami de ces dames ; rien n’est plus naturel, on ne pourrait
trouver à l’âge de 18 ans avec un extérieur très agréable, autant de solidité (Où était-elle ?) de
caractère avec beaucoup de modestie et d’usage du monde. Sur les onze heures je quittai ces dames
pour aller me mettre dans la diligence, ce jeune homme eut la complaisance de m’accompagner et ne me
quitta qu’au moment

\marginpar{\footnotesize\raggedright 1799. Septembre.} du départ, contre minuit. Content et
satisfait de tout ce qui m’était survenu dans ce pays, encore l’asile fortune de cette abondance de
bonheur qui inondait toute la Suisse avant que les prôneurs de la liberté vinssent souiller son
territoire. Le plaisir de faire la narration d’un espace de trois jours qui s’est écoulé aussi
agréablement, m’a fait omettre deux autres circonstances dont je ne voudrais pas perdre le souvenir.
Déjà le 3 de juillet la nature paraissait demander un nouveau soulagement que la raison lui avait
refusé depuis Lucerne ; mon ami S. me proposa d’aller à cette maison de campagne située au Rhein,
nous nous fîmes servir un goûter auquel deux filles assistèrent, nous passâmes la soirée ensemble on
s’égara de part et d’autre et la chose n’eut plus aucune suite fâcheuse ; pendant quelque temps j’en
avais des inquiétudes, elles servirent à me retenir dans les bornes de la sagesse dont pendant le
reste de l’année je ne me suis plus écarté que quatre fois. Tout le mois de juillet et d’août ma
grand- mére maternelle dépérissait considérablement ; parvenue à l’âge de 90 ans, avec les
faiblesses qui y sont attachées, on ne pouvait que faire des voeux pour sa prompte délivrance, et le
samedi à 3 heures du soir au 31 août il plut au Seigneur de la retirer après plusieurs violentes
convulsions. Mon coeur dépose ici les sentiments de reconnaissance pour les bienfaits que cette
chère grand-maman n’a discontinué à nous rendre pendant tout le temps qu’elle nous a connus. De
retour à Berne je n’eus rien de plus empressé que de dire à mes parents l’heureuse issue de ma
visite et je repartis soudain pour aller finir mes affaires à Aarau ; je fis ce voyage avec un
employé français fort honnête et point républicain ; à Kirchberg mon frère me donna sa chaise et
m’accompagna avec sa femme et son petit Fritz jusqu’à Herzogenbuchsee où je fis servir du café pour
l’aimable compagnie ; le Français et moi continuâmes notre route et à 11 heures du soir nous fûmes à
Aarau ; ne me proposant point d’y rester longtemps je ne voulus pas retourner chez mon hôte et
préférai loger à l’auberge du Sauvage ; mon ami Hurner m’offrit sa table le plus obligeamment du
monde et j’en ai profité quelquefois dans ce peu de jours ; la mort lui enleva dans ce temps sa
soeur à laquelle je fermis les yeux, tourmentée

\marginpar{\footnotesize\raggedright 1799. Septembre.} depuis longtemps d’une maladie bien pénible,
on dut remercier la providence d’avoir raccourci ses maux, mais la tendresse fraternelle pleura
néanmoins sa mort. Je quittai Aarau au 22 septembre, j’espère que ce sera pour toujours, déjà
pendant mon premier séjour, Mme Wydler était allée faire une cure au bain de Lostdorf ; je lui
rendis une visite avec son mari et j’y payai le goûter, je pouvais d’autant mieux le faire que pour
le temps que j’ai séjourné chez M. Wydler, on n’a exigé de moi qu’une pension excessivement modique.
Immédiatement après mon arrivée à Berne, on apprit la malheureuse reprise de Zurich par les troupes
françaises, et la mort du valeureux Général Hotz. L’effet que d’aussi funestes nouvelles
produisirent sur l’esprit des Bernois est difficile à dépeindre ; autant qu’on se flattait
auparavant des succès de l’armée coalisée, autant redoutait-on maintenant les progrès des
républicains. Ses soi-disant patriotes se croyaient au plus haut degré de prospérité et la
destruction des trônes ne leur paraissait qu’une fort légère tâche. Je voyais dans cet événement le
sort de la guerre, et dans la défaite totale de l’armée russe le peu d’accord qui régnait parmi les
coalisés et les grands retards qu’éprouverait encore notre délivrance du joug français et c’est là
où je laissai la politique. Les derniers jours de septembre et la moitié du mois d’octobre se
passaient néanmoins dans les plaisirs et les réjouissances. L’agréable saison nous engageait à de
fréquentes promenades avec nos demoiselles. Le goût de la nouveauté attachait alors mes sens à la
beauté de Mlle Emilie Ernst, la jeunesse et l’esprit de cette belle rose naissante paraissaient
mériter mes soins, je les lui vouais pendant quelques semaines et lui adressai même des vers sur un
sujet indifférent, ils furent très bien accueillis. Après Emilie je distinguais le plus sa soeur
Mélanie et souvent la modestie et la douceur me parurent préférables à la beauté qui s’en trouvait
dépourvue. Cette indécision prouve assez que l’amour n’avait aucune part dans cet espèce de
penchant, car il ne

\marginpar{\footnotesize\raggedright 1799. Septembre.} m’aurait pas laissé dans cette incertitude
sur le choix que je devais faire. En parlant des personnes du beau sexe qui faisaient plus ou moins
d’impression sur mon coeur ; je ne dois pas oublier les amis qui dans ce temps contribuaient à ma
satisfaction. Je voyais le plus souvent Muralt, Ernst Zeerleder et Crousaz ; ce dernier est membre
du canton Léman au tribunal suprême. Le temps s’approchait où je devais retourner à Neuchâtel comme
je l’avais convenu avec M. le Baron ; je partis de Berne le 18 octobre, Mlle Mélanie Ernst eut la
complaisance de porter elle-même à la maison, une lettre de recommandation pour Mme de Pierre
Bosset, et l’ami Ernst m’en remit plusieurs pour ses amis. J’arrivai à Cormondrèche sans que rien
d’intéressant me survint ; j’étais alors fort inquiété de maux de dents ; M. le Baron me reçut très
bien ; pendant la première semaine de mon séjour auprès de lui, mes douleurs me rendirent peu
susceptible aux jouissances que je me promettais ; il me conduisit peu de jours après mon arrivée
chez Mme Chambrier sa tante, où on m’accueillit au mieux, il m’accompagna de même chez M. le
Procureur général de Pierre et de là chez Mme de Pierre Bosset la bonne amie de Mlle Ernst. Je me
rendis seul chez Mme Bourgeois et chez M. le Chancelier Boive, un ancien ami de mon père, qui me
témoigna beaucoup d’amitié. J’écrivis dans ce temps deux lettres à mes parents, où je leur détaillai
les moindres circonstances de ce qui s’était passé depuis mon départ de Berne ; M. le Baron me fit
faire la connaissance de la famille Chambrier d’Auvernier, où je rencontrai Mlle Wagner soeur du
Principal. J’ai reçu beaucoup de politesses dans cette maison, ainsi que dans celle de M. le Colonel
Bedeau à Cormondrèche où j’allais passer plusieurs soirées–. M. de Chambrier toujours occupé à me
procurer du plaisir me proposa d’aller voir la fête des Armo–rins à Neuchâtel. Cette fête se célèbre
toutes les années, 24 jeunes gens de famille habillés en cuirassiers armés de hallebardes, vont au
château complimenter le Gouverneur et boire à la santé du roi. Les uns prétendent que cette
cérémonie doit faire valoir le droit des Neuchâtelois d’entrer armés au château de leur ville,
d’autres croient que c’est un hommage qu’on rend au souverain ; il me paraît qu’elle peut être l’un
et l’autre ; quarante jeunes garçons de 4 à 10 ans habillés en blanc et garnis de rubans, portent
des flambeaux pour éclairer les cuirassiers qui s’acheminent à nuit close embellissent cette fête.
Je jouissai peu de ce joli coup d’oeil à cause de mes maux de dents, j’en avais d’autant plus de
regrets que j’étais entouré de la plus aimable société qui fêtait réunie chez M. de Pierre, tout le
monde était gai, joyeux et content, je souffrais tout seul. J’allai passer

\marginpar{\footnotesize\raggedright 1799. Octobre.} la nuit dans la maison Chambrier où on
s’empressa à calmer ma douleur par les plus grâcieux soins, le tabac, de l’opium et un bain de pieds
me procurèrent une nuit tranquille, mais avec le jour mes tourments revinrent ; j’eus recours à un
arracheur de dents qui m’opéra assez bien, sans m’ôter la douleur qui ne me quitta qu’après des
remèdes stomatiques. Un autre événement vint troubler la joie à la-quelle je m’abandonnais depuis ma
guérison ; un jour à la promenade M. le Baron me dit qu’il avait reçu des nouveller de Florence et
que M. Interlattner son secrétaire de légation persistait à conserver une place qu’il desservait
depuis 15 ans au contentement du ministre, mais qu’il désirerait que celui-ci s’employât en sa
faveur pour obtenir de Berlin une augmentation d’appointements. Je fus stupéfait au premier abord
d’entendre un arrêt qui tranchait le fil à toutes mes espérances ; je voyais en même temps qu’il en
coûtait à M. le Baron qui m’avait déjà mis au fait du chiffre de cabinet, et auquel cette réponse du
secrétaire n’était pas indifférente après les relations entamées avec moi. Je lui dis que dans ce
cas ma présence à Cormondrèche devenait inutile ; après une conversation très longue où M. de
Chambrier démontrait toujours à mon égard une bienveillance et une générosité qui inquiétèrent ma
délicatesse il finit à peu près ainsi. Je serais fâché Monsieur, d’interrompre les liaisons d’amitié
ue j’ai contracté avec vous et quoique les circonstances viennent de changer, vos services me seront
toujours nécessaires ; j’ai besoin d’une personne de la plus grande confiance pour soigner mes
intérêts politiques en Suisse, m’arrive très souvent d’avoir des objets d’importance à faire passer
de Turin à Berlin, à Paris etc. et il m’importe d’avoir quelqu’un à Berne pour ce service-là, il
m’est d’ailleurs très agréable de recevoir de temps en temps des nouvelles politiques de la Suisse ;
vous serez à portée M. de me satisfaire à cet égard, le gouvernement se trouvant à Berne, il vous
sera facile à vous mettre au fait de tout ce qui s’y passe d’intéressant ; ma maison continuera à
vous être toujours ouverte et plus vous viendrez me voir, plus vous aurez de droit à ma
reconnaissance ; il ne s’agit plus M. que de savoir quelle indemnité vous désirez pour les peines
que je vous donne ? Je répondis à M. le Baron que comme j’étais déjà son débiteur pour les bienfaits
dont il m’avait comblés depuis l’instant que j’avais eu l’honneur de faire sa connaissance, je me
faisais un plaisir et un devoir de faire en ami et sans indemnité quelconque, tout

\marginpar{\footnotesize\raggedright 1799. Octobre.} ce qui pourrait lui être agréable et que
pendant tout le temps que je serais à Berne il n’avait qu’à disposer de moi. M. me répliqua-t-il, on
ne travaille pas pour rien, quand il s’agit de servir une Cour ; on ne vous en aurait aucune
obligation, ainsi ne vous faites aucun scrupule de fixer quel dédommagement vous demandez pour la
complaisance qu’on exige de vous ; après mille difficultés que je fis et qui n’eurent aucun succès,
je priai M. le Baron de me dispenser à prescrire la somme qui devait être la récompense d’un aussi
petit service, l’assurant que je me contenterais de tout ce qu’il jugerait à propos ; il me demanda
si trente louis me paraissaient suffisants, et malgré mes protestations que les deux tiers
suffiraient il ne voulut point se désister du nombre de trente et me les promit pour tout le temps
qu’il serait ministre ; quand je serai à Turin ajouta-t-il vous viendrez m’y trouver, vous logerez
chez moi à l’hôtel de Prusse,vVous reverrez toutes vos connaissances et si l’envie vous prend de
servir, je me donnerai les peines nécessaires pour vous procurer une compagnie au service du roi de
Sardaigne. Ces généreuses offres, présentées de la manière la plus obligeante, en abaissant toujours
le prix que je devais en faire, pro-duisirent sur mon coeur des sensations jusqu’ici inconnues ;
sans les manifester avec transport, je témoignais à M. le Baron ma plus vive reconnaissance. Déjà
avant tout ceci je lui avais fait sentir que ses procédés généreux coûtaient à ma délicatesse, sur
quoi, il me dit, M. je ne suis point généreux, et soyez convaincu que si je voulais l’être, ce ne
serait jamais vis-à-vis de vous. Que pouvais-je répliquer ; je consentis à tout ce qu’il voulut et
je me félicitais d’un destin qui me mettait en relations avec un homme aussi rare et aussi
respectable que M. de Chambrier. Après cette longue explication il ne fut jamais plus question
d’intérêts matériels les arran–gements étaient pris, et dès ce temps je me trouvais attaché à la
mission de Prusse ; je copiais moi-même le chiffre dont je devais me servir pour la correspondance
avec M. le Baron depuis Berne, pour les objets qui exigeraient le secret. Le genre de vie de
Cormondrèche, quoique bien différent de celui de Berne, ne me convenait pas moins, le temps
s’écoulait avec rapidité. Le matin à huit heures je me levais pour aller déjeuner.

\marginpar{\footnotesize\raggedright 1799. Novembre.} Après neuf heures chacun se retirait dans sa
chambre et s’occupait jusqu’à l’heure du dîner d’agréables instructions et lectures, M. le Baron me
donna à lire l’ouvrage de l’Abbé Barruel, sur l’origine et les causes de la révolution. Cet auteur
remonte au temps des premiers philosophes de ce siècle et attribue surtout à Voltaire et à Knigge
les progrès des principes révolutionnaires. Ce livre qui est assez bien écrit, dévoile une infinité
de sources secrètes où on ne serait pas allé chercher le germe de ce système dévastateur ; je lus en
même tepms les Mémoires de Fréderic second publiés après sa mort ; M. Bedeau me prêta la
Bibliothéque britannique où je lisais des traductions de Mme de Radclif ; je n’eûs pas le temps
d’achever les lettres sur la Sardaigne par l’Abbé D. Rien ne frappa mon imagination comme le récit
de Ramel au sujet des déportés français à la Guyane. On ne quittait son appartement que pour aller
faire un excellent repas, où chacun animait la conversation par le récit de ce qu’il lui avait paru
le plus intéressant de sa lecture ; après midi nous restions réunis dans la chambre si le temps ne
permettait pas d’aller faire un tour de promenade ; les jours de courrier, M. le Baron recevait
régulièrement des lettres et les papiers français, qui d’ordinaire nous donnaient matière à la
conversation ; souvent les soirées se passaient chez Bedeau, où j’étais très honnêtement reçu ; Mlle
Rosette Bourgeois se trouvait alors en visite chez Mme Bedeau et aug–mentait par sa société
charmante les agréments de celle de la maison Bedeau. M. le Baron sortait rarement, et il m’a même
paru que souvent il m’accompagnait par complaisance. Les moments qui n’étaient point employés à la
lecture de livres utiles servaient à ma correspondance ; j’écrivis à Mme Beaude pour lui donner des
nouvelles du changement de ma carrière, je reçus sa réponse ; j’adressai deux lettres à mon ami R.
Stettler l’une à Venise l’autre à Turin, sans apprendre autre chose que son séjour dans cette
dernière ville, d’où il doit être parti depuis ; j’étais en commerce de lettres suivi avec Ernst qui
m’écrivait les plus amicales et joviales épîtres, j’écrivis le 1er à Zeerleder qui m’en avait prié ;
je fus très content de sa réponse, elle fut suivie d’une seconde lettre de ma part

\marginpar{\footnotesize\raggedright 1799. Novembre.} peu de jours avant de quitter Cormondrèche ;
je lui disais ainsi qu’à Ernst plusieurs agréables choses pour le compte de ses soeurs. Un soir, que
par manière de conversation je disais à M. de Chambrier que feue ma grand-mère avait laissé un
journal de sa vie sur des feuilles volantes, qui nous mettait au fait de ce qui s’était passé il y a
70 ans ; il me dit avec beaucoup d’intérêt : comment M. Mme votre grand-mère a fait son journal ? Et
ne suivez-vous point un aussi bon exemple ? non seulement, répliquai-je, ma grand-mére mais ma mère
et ma soeur aînée ont entrepris cet ouvrage. M. le Baron parut charmé de ce que je lui disais, et me
pria de monter dans sa chambre où il ouvrit un buffet qui contenait 13 volumes in quarto parmi
d’autres livres de toute espèce, il prit un des premiers et me dit : Voilà mon journal ! je ne puis
assez vous dire ajouta-t-il combien de satisfaction il me procure, et je veux absolument vous
engager à faire le vôtre, suivez le bon exemple de Mme votre mère. J’avais beau lui remarquer que ma
vie était trop obscure pour mettre par écrit les événements qui y survenaient, que d’ailleurs je me
trouvais à un âge trop avancé pour commencer un ouvrage aussi étendu, qui ne pourrait qu’être
insipide ; il m’assura qu’il avait commencé le sien à peu près à mon âge, et sur les instan–ces d’un
digne et respectable ami le Proferseur Formey qui pendant son séjour à l’académie de Berlin avait
pris soin de son éducation, que pour composer le premier volume il s’était rappelé à sa mémoire les
époques majeures de sa vie, et que rien ne me serait plus facile que d’en faire de même ; dans ma
maison chacun fait le sien, l’instituteur, mon fils adoptif ; allons M. Stettler achetez des livres
et suivez mon conseil, vous verrez que vous ne vous en repentirez pas ; je résolus d’obéir et le
lendemain de cette conversation je me rendis en ville pour empletter le livre où j’écris et le même
jour j’en remplissai plusieurs pages. La première fois que j’écrivis à mes parents, je leur fis part
de ma résolution, en demandant leurs sentiments là-dessûs, ma mère me répondit que certainement un
journal pouvait être très utile, qu’elle l’envisageait comme un compte moral qu’on se rendait de ses
actions, et que dès qu’on l’entreprenait dans le but de devenir meilleur en se dévoilant ses propres
défauts avec franchise, cet ouvrage ne pourrait

\marginpar{\footnotesize\raggedright 1799. Novembre.} qu’être salutaire ; que le sien avait et était
encore pour elle une source de satisfactions. Tant de motifs encouragèrent mon coeur à un travail
que je voyais entouré de mille difficultés ; je me proposai de commencer mon vrai journal au
commencement de l’an 1800 mais je n’ai pu finir l’histoire de ma vie jusqu’à ce jour-là. Les jours
de novembre se passaient agréablement au milieu des vendanges ; j’assistai à une fête champêtre un
dimanche des vignerons à Auvernier, où il eut un bal dans une grange. J’éprouvai des sensations
pénibles en mêlant à ces innocents plaisirs l’idée des malheurs qui affligeaient ma patrie ; cette
même rélexion trouble souvent les récréations les plus permises que je me procure à Berne, mais la
légèreté reprend bien vite tous ses droits. Une lettre que je reçus de mon ami Ernst m’apprit
l’arrivée du jeune de Muralt, officier au service de Prusse, à Neuchâtel ; je quittai Cormondrèche
pour aller à la rencontre de cet ami, et j’allai loger chez ces dames Bourgeois qui me reçurent
aussi amicalement que possible ; j’attendis Muralt pendant deux jours, Mlle Bourgeois m’introduisit
dans sa société et le Baron Steinwehr officier recruteur me présenta à celle de la belle Mme de
Pierre Bosset ; le 3ème jour Muralt arriva enfin ; j’eus un grand plaisir de revoir ce jeune homme
en élégant uniforme, qui trois ans auparavant était jaloux quand il voyait le mien. Je courus toute
la ville avec mon ami, M. Steinwehr nous donna à dîner aux Balances ; je l’invitai en échange avec
Muralt à venir avec moi à Cormondrèche ce que M. le Baron permettait et désirait ; Mesdames
Bourgeois nous donnèrent une jolie soirée, où il y avait les deux officiers prussiens, du Roveray et
de Brandis, qui était venu avec Muralt, et moi, on fit la belle conversation avec ces deux dames
tout le soir. Le lendemain le petit char de Muralt nous conduisit à Cormondrèche Steinwehr Muralt et
moi ; M. le Baron nous donna un excellent dîner, on parla beaucoup de Berlin, après nous fîmes une
visite à la maison Bedeau et sur le tard mes amis repartirent. Le jour après je fus en ville avec M.
le Baron pour aller dîner chez M. le Gouverneur M. de Béville ; Steinwehr qui en majordome en fit
les honneurs me plaça à côté de Mme De Luze Messeract, femme du monde et belle encore quoiqu’à un
âge où ces frivoles avantages commencent à disparaître ; nous étions 24 personnes à table et le
repas fut gai. Muralt projetait de repartir pour Berne et m’engagea à faire le voyage avec lui.
Comme je m’apercevais

\marginpar{\footnotesize\raggedright 1799. Novembre.} fort bien que ma présence à Cormondrèche était
absolument inutile et qu’en me trouvant à Berne je serai par contre dans le cas de donner à M. le
Baron les nouvelles de ce pays là, je résolus de partir avec Muralt, et pris toutes mes mesures pour
quitter ce séjour. Je veux avant de m’en éloigner parler encore de M. Poinsieux l’instituteur et du
jeune Chambrier ; le premier est à l’âge de 30 ans, fils du ministre de Clarans ; élevé à Lausanne,
parmi la classe des ecclésiastiques je lui suppose beaucoup de bonnes connaissances ; l’avantage
d’être placé dans une bonne maison, et le séjour qu’il a fait à Turin flattent sa vanité, qui avec
l’air doucereux dont j’ai déjà parlé diminuent les agréments que sans cela sa societé pourrait
procurer ; j’ai conçu néanmoins beaucoup d’estime pour cet homme, ses sen–timens politiques
s’approchent infiniment des miens. M. de Chambrier fils adoptif de M. le Baron est à l’âge le plus
pénible de la vie à 14 ans. Les plaisirs enfantins dont on voudrait jouir encore sont toujours
entravés par les reproches que fait la raison naissante avec laquelle ces enfan–tillages se trouvent
en contradiction. Pour diminuer ce malaise autant que possible il est bon de tenir les jeunes gens
de cet âge assidu aux occupations utiles ; aussi M. le Baron est-il convaincu de la vérité de cette
maxime, et son fils passe la plus grande partie du jour à l’étude, et le soir pour sa récréation il
travaille à son journal qu’il a déjà commencé à l’âge de 12 ans. Frédéric est un beau jeune homme
rempli d’espérance, il est froid de son naturel et paraît n’avoir de goût décidé pour rien, son père
s’aperçoit avec plaisir de cette indifférence, il espère pouvoir à son temps le fixer à loisir ; je
crois qu’il le destine à entrer dans la carrière diplomatique et je suis persuadé que pour celle-ci
beaucoup de sang froid et une grande réserve sont des qualités essentielles ; il faut à la Cour des
gens qui soient difficiles à définir, ne dussent-ils souvent n’avoir rien à cacher il est bon encore
de faire paraître le contraire. Je témoignais toujours beaucoup d’amitié au jeune homme, il ne parut
cependant point rechercher ma société, j’ignore si c’est par timidité par modestie ou par ce que je
ne lui plaisais pas ; quelquefois le soir nous faisions une partie au trictrac à laquelle il
paraissait prendre plaisir, mais il lui en coûtait toujours de me proposer la partie le premier

\marginpar{\footnotesize\raggedright 1799. Novembre.} Je quittai Cormondrèche sans regret ; la
conviction que je n’y étais d’aucune utilité la persuasion que Berne m’offrait des agréments qu’un
séjour à la campagne été et hiver ne pouvait me procurer, et le plaisir que j’avais de revoir mes
parents, amis et connaissances après une absence de six semaines, me consolait de la perte que je
faisais en quittant M. de Chambrier dont la société était à la fois agréable et instructive ; en me
séparant de lui il fut aussi honnête à exiger de moi une promesse positive que je ne tarderais pas
trop longtemps à retourner chez lui. J’allai faire ma visite de congé à la maison Bedeau, et dans
l’après-midi du 22 novembre je partis pour Neuchâtel dans la voiture de M. le Baron je passai la
soirée chez Mlle Bourgeois où je fis la partie au whist avec Mme Sandoz Borel, la maîtresse de la
maison et le jeune Du Roveray. À 9 heures je me retirai avec Muralt à l’auberge du Faucon ; nous
restâmes tard près du feu et allâmes nous coucher ; que de rires et de farces de cette nuit je me le
rappellerai toujours, surtout de l’horrible servante. Le matin après un mauvais déjeuner nous
partîmes en char par Berne où nous arrivâmes sur les 6 heures du soir, sans que rien d’in–téressant
nous soit survenu ; Muralt m’entretenait beaucoup des agréments du service de Prusse et en même
temps des regrets qu’il avait à retourner à Berlin tandis qu’il se trouvait si bien dans sa chère
patrie. La fin de novembre et le commencement de décembre se passa fort agréablement, des soirées,
des parties de danse des parties de whist et le Falkenlaist prenaient mon temps, je travaillais de
même tous les jours à mon journal, je lisais des extraits que mon frère avait fait sur l’histoire
d’Angleterre. Dans nos sociétés je donnais encore la préférence à ces Dlles Ernst, je m’occupais à
faire des vers que je n’ai cepandant jamais voulu produire ; je voyais beaucoup le frère de ces
Dlles son humeur gaie, son esprit, et d’autres bonnes qualités lui ont gagné toute mon amitié ; je
lui reproche pourtant un caractère trop léger, quoique je le croie susceptible de sentiments
solides, il me paraît néanmoins que son coeur tient plus au frivole et à l’inconsidéré

\marginpar{\footnotesize\raggedright 1799. Décembre.} Le 16 décembre on arrangea un souper pour le
jeune de Muralt, qui comptait rejoindre son régiment ; nous étions 18 à table, au période des
santés, le jeune Tscharner du Lohn porta celle de M. l’Avoyer Steiger, dont le même soir on avait
appris la mort ; je trouvai cette santé malplacée par la raison que je viens d’alléguer et en second
lieu il y avait à ce souper un nombre de jeunes employés au gouvernement qui s’exposaient plus ou
moins de la boire, ceux qui la portèrent sautèrent sur la table et firent un vacarme terrible ; je
leur dis qu’il me paraissait qu’on aurait dû se dispenser de boire à la santé d’un homme que nous
devions plutôt pleurer puisqu’il venait de mourir, et je touchai en silence mon verre avec mes
voisins qui étaient tous secrétaire dans les bureaux des différentes autorités, et celui qui devait
se trouver le plus embarrassé était Crousaz le Grand juge. La modération que la réflexion me fit
mettre à un acte auquel mon coeur aurait applaudi dans toute autre circonstance, excita le blâme de
tous ceux qui voulaient mettre le raisonnement de côté, et me fit taxer l’ennemi de Son Excellence
Steiger ; cependant personne ne dit rien, et un moment après le même Tscharner s’écria en me fixant,
Citoyens et Messieurs je vous invite à boire à la santé de l’archiduc Charles. Cette dénomination de
citoyen qui s’adressait particuli–èrement à moi par le plus jeune convive que je connaissais fort
peu me choqua beaucoup, je me fâchai et lui dis, il n’y a point de distinction pareille à faire
entre nous, et ce ne serait jamais à vous qu’elle appartiendrait etc. On s’échauffa de part et
d’autre R. Fischer s’en mêla aussi et je me proposai déjà d’en tirer satisfaction le matin.
Tschiffeli de Roche, qui connaissait mon humeur et combien peu j’étais accoutumé à me laisser
marcher sur le pied, s’aperçut sur-le-champ de l’impression que tout ceci faisait sur moi ; il me
prit à part et me parla raison ; après bien des efforts et une généreuse condescendance de R.
Fischer, à laquelle je répondis, on vint à bout de me réconcilier avec lui et Tscharner ; nous bûmes
ensemble Décembre. et on parut oublier le passé pour la fin de la fête. Je proposai à R. Fischer de
nous tutoyer à l’avenir, pour lui témoigner que je ne portais nullement rancune ; je fis la même
offre à Zeerleder et Watteville de Malessert me l’offrit de son côté. Après avoir bien tablé nous
nous rendîmes au n°3 où avec Ernst j’eus un moment d’oubli et sur le matin je m’en allai coucher. Je
ne m’imaginais point que l’affaire que j’avais eue au souper aurait des suites, et ce ne fut pas
sans étonnement que j’appris que dans toute le ville le bruit courait qu’avec plusieurs autres je
m’étais refusé à boire la santé de Son Execellence Steiger, je contai pour toute justification la
chose telle qu’elle s’était passée et recueillis des uns le blâme, des autres l’approbation. Peu de
jours après je partis pour Schüpfen, où je passais des jours tranquilles et heureux ; le lende–main
de mon arrivée, je me rendis au soir avec mon beau-frère à Meikirch chez M. le ministre Leu, nous y
fîmes une partie de whist, y couchant et le jour suivant nous reprîmes le chemin de la maison. Ma
soeur s’occupait avec moi à faire des vers pour les demoiselles de ma société, que je me proposais
leur donner au Nouvel An, nous faisions aussi beaucoup de charades, je m’amusais de temps en temps
avec Henriette ma nièce et son gentil frère Fritz, tous les deux avaient un grand plaisir de m’avoir
auprès d’eux, ainsi que la petite Marianne Langhans nièce de mon beau-frère. Le jour de Noël je fus
au sermon et à la Communion, et le lendemain je m’en retournai à Berne très content de mon bien
agréable séjour. C’était un jeudi, le mercredi qui le précédait était le jour fixé pour la réception
des membres du Falkenleist, il y avait sur les rangs : Ernst, Muralt, Haller, Fischer, Zehender, mon
frère et moi ; je l’avais fréquenté depuis longtemps déjà je connaissais l’esprit qui y régnait et
j’ai pensé une ou deux fois à Schüpfen, qu’il serait possible que je fusse flambé ; alors une voix
intérieure m’assurait que non ; ne me connaissant aucun tort vis-à-vis d’aucun des membres, mais je
savais que le bruit qu’on avait débité après ce souper me mettrait en très mauvaise odeur près les
membres

\marginpar{\footnotesize\raggedright 1799. Décembre} de cette société. Arrivant à la maison je
trouvai mon ami Ernst chez mon frère, ils vinrent tous deux dans ma chambre et me dirent que sur 33
voix j’avais eu 16 négatives et mon frère 12 ainsi que nous étions tous les deux flambés ; Ernst par
contre avait été reçu avec 11 contre. Au premier moment cette nouvelle me surprit plus qu’elle ne me
peina, je sus d’abord que le souper en question en avait été la cause ; je ne comprenais par contre
point les raisons qui pouvaient avoir engagé les dignes membres à flamber mon frère qui ne s’était
point comme moi rendu coupable à leurs yeux, cependant il paraissait, ainsi que moi, prendre
gaiement son parti. Mon ami Ernst a surtout dans cette occasion manifesté ses sentiments à mon égard
; il résolut de renoncer à ce Leist malgré qu’il avait été reçu, mais on ne voulut point accepter
cette renonciation, et néanmoins il n’y a jamais mis le pied. Le même soir que j’appris cet accident
je fus à la société où je craignis que n’ayant entendu que par mes ennemis les raisons qui ont
occasionné cette mésaventure, ces Dlles me jugeraient mal et me verraient de mauvais oeil ; je me
suis bien trompé ; elles me reçurent fort honnêtement et me dirent qu’elles espéraient que je ne
m’affecterais pas de ce qu’il m’était arrivé ; je leur répondis que ne me connaissant aucun tort,
j’avais été bien vite consolé ; mon père dit à mon frère et moi que pour lui il était fort content
que nous n’étions pas de ce Leist, où, par ce qu’il venait d’y arriver, il paraissait régner un
mauvais ton. Je conviens cependant qu’il est toujours désagréable d’être exclu par ses plus
anciennes connaissances mêmes, d’une société qu’on a pendant longtemps fréquenté dans sa ville
natale et s’il n’y existait une conviction entière que ce n’est point par sa façon de penser et sa
conduite qu’on s’est attiré ce désagrément ; que la conscience nous dit : tu penses comme un honnête
homme tu agis de même, et c’est à tort que les hommes te jugent différemment ; je conviens dis-je,
que j’aurais eu le plus sensible chagrin de ce qui s’était passé ; en échange, ces différentes
réflexions m’ont absolument fait oublier un tort que mon coeur ne croit point

\marginpar{\footnotesize\raggedright 1799. Décembre.} avoir mérité. Le dernier jour de l’année, qui
d’ordinaire est un jour de fête, se passa chez Mme Freudenreich de Thorberg, où on se rendit à 6
heures du soir on y fit un très bon goûter, après lequel les jeunes se mirent à courir par les
arcades pour voir les boutiques qui étaient illuminées, on s’arrêta surtout chez le confisseur Bay
où nous achetâmes des biscomes avec de beaux ours pour régaler nos Dlles. Sur le tard nous nous en
retournâmes chez Mme Freudenreich où je fis la partie avec M.E.G.J. et Muralt ; on finit la soirée
en faisant le jeu du juge, je jouais ce rôle au contentement de la société, à ce qu’on m’a dit.
Avant de terminer cette année je veux remarquer encore que le désagrément d’avoir été refusé au
Leist a eu pour moi la suite que même mon ancien ami L. de Muralt s’est mis en froid avec moi. Il me
paraissait singulier qu’au seul moment où il aurait eu occasion de me donner des preuves d’amitié,
me témoignant ses regrets de ce qui venait de m’arriver, il s’est détaché de moi comme si je l’avais
sérieusement offensé. Trop fier de manifester la moindre sensibilité de la perte d’un ami dont j’ai
cru devoir faire cas pendant trois ans, j’ai gémi en secret sur le peu de fond qu’on pouvait mettre
aujourd’hui dans de pareilles liaisons et je me suis consolé avec l’idée que puisque cette petite
disgrâce qui m’était survenue avait pu ané–antir les sentiments d’amitié qui paraissaient nous unir,
ils n’auraient jamais été bien sincères de son côté, ainsi la perte ne doit et ne peut en être
sensible. Je me trouvais donc à la fin de cette année dans une situation assez désagréable pour le
moment ; privé d’une société d’amis où je me proposai passer une grande partie de l’hiver, désuni
avec un de ceux que je préférais, décrié par la ville pour un homme à sen–timents suspects, par mes
plus anciennes connaissances, puisque j’avais bu à regret à la santé de Son Excellence Steiger ;
incertain par-dessus cela de l’espèce d’établissement auprès M. de Chambrier, avec lequel je
corres–pondais néanmoins toujours, je n’avais que peu de jouissances à attendre de l’avenir ;
j’étais à côté de cela dépourvu d’argent, à l’époque où je devais payer tous mes comptes ; mon père
se trouvant peu disposé à me

\marginpar{\footnotesize\raggedright 1800. Janvier.} pourvoir de tout, dans un âge où il devait
s’attendre que je serais dans le cas de gagner mon pain moi-même.

\end{document}
